% Utilisation des structures de contrôle — Informatique 3ème — Cours signé OnBuch+
% Programme officiel MINESEC. Compiler avec : tectonic N14-utilisation-structures-controle.tex
\newcommand{\DOCMATIERE}{Informatique}
\newcommand{\DOCNIVEAU}{3ème}
\newcommand{\DOCMODULE}{Informatique – classe de 3ème}
\newcommand{\DOCLECON}{N14}
\newcommand{\DOCTITRE}{Utilisation des structures de contrôle}
% OnBuch+ — préambule « cours signés » ENRICHI (compilé avec tectonic / XeLaTeX)
% Système visuel « Pop » d'OnBuch+ : fond crème (jamais blanc), bordures encre
% épaisses, ombres « dures » décalées, titres Archivo Black en capitales.
% Version MATHÉMATIQUES : graphes (pgfplots/tikz), boîtes pédagogiques
% (définition, propriété, méthode, attention, le savais-tu, exercice résolu,
% exercice, TP, bilan), page de garde et signature OnBuch+ ancrée en bas.
%
% Macros attendues dans main.tex AVANT \input{preamble} :
%   \DOCMATIERE  \DOCNIVEAU  \DOCMODULE  \DOCLECON (numéro)  \DOCTITRE
\documentclass[11pt]{article}

\usepackage{fontspec}
\usepackage[french]{babel}
\usepackage[a4paper,margin=2cm,top=3.4cm,bottom=2.8cm,headheight=1.4cm,footskip=1.3cm]{geometry}
\usepackage{amsmath,amssymb,amsfonts}
\usepackage{mathtools} % \xmapsto, \coloneqq... souvent utilisés par les modèles
\usepackage{cancel} % simplifications barrées (bilans de forces)
% \abs{z} (module, valeur absolue) et \norm{u} (norme), utilisés par les modèles
\providecommand{\abs}{}\let\abs\relax\DeclarePairedDelimiter{\abs}{\lvert}{\rvert}
\providecommand{\norm}{}\let\norm\relax\DeclarePairedDelimiter{\norm}{\lVert}{\rVert}
\usepackage[table]{xcolor} % [table] : fournit aussi \rowcolors (alternance de lignes)
\usepackage{pagecolor}
\usepackage{tikz}
\usetikzlibrary{arrows.meta,positioning,calc,shapes.geometric,shapes.misc,shapes.arrows,shapes.symbols,decorations.pathmorphing,decorations.pathreplacing,decorations.markings,patterns,shadows,mindmap,trees,fit,backgrounds,matrix,intersections,angles,quotes}
\usepackage{pgfplots}
\usepackage[version=4]{mhchem} % équations nucléaires \ce{^{235}_{92}U}
\usepackage[european,siunitx]{circuitikz} % schémas électriques
\pgfplotsset{compat=1.18}
\usepgfplotslibrary{fillbetween,groupplots} % aires entre courbes (calcul intégral)
\usepackage{enumitem}
\usepackage{fancyhdr}
% [most] charge toutes les bibliothèques tcolorbox (listings, minted, raster,
% documentation...) et gonfle inutilement la mémoire TeX sur un document long
% (~300 boîtes) : on ne charge que ce qui est réellement utilisé.
\usepackage[skins,breakable]{tcolorbox}
\usepackage{graphicx}
\usepackage{colortbl}
\usepackage{booktabs}
\usepackage{tabularx}
\usepackage{diagbox} % case d'en-tête diagonale des tableaux à double entrée
% Colonnes extensibles centrée / gauche / droite (C, L, R), souvent utilisées par les modèles
\newcolumntype{C}{>{\centering\arraybackslash}X}
\newcolumntype{L}{>{\raggedright\arraybackslash}X}
\newcolumntype{R}{>{\raggedleft\arraybackslash}X}
\usepackage{array}
\usepackage{multirow}
\usepackage{multicol}
\usepackage{siunitx}
% parse-numbers=false : accepte aussi \SI{2,5 \times 10^{-3}}{...}
\sisetup{locale=FR,output-decimal-marker={,},group-minimum-digits=4,parse-numbers=false}
\DeclareSIUnit\ohm{\ensuremath{\symup{\Omega}}} % U+2126 absent de la police
\DeclareSIUnit\division{div} % réglages d'oscilloscope (ms/div, V/div)
\DeclareSIUnit\tr{tr} % tours (vitesse de rotation, ex. \SI{1500}{\tr\per\minute})
\DeclareSIUnit\bit{bit}
\DeclareSIUnit\byte{o} % octet
\DeclareSIUnit\inch{po} % pouce
\DeclareSIUnit\ppp{ppp} % points par pouce (résolution d'image)
\usepackage{qrcode}
% Blocs de code source (PHP, C, HTML, JavaScript, SQL) — spécifique à la
% matière Informatique : listings + habillage aux couleurs OnBuch+ (voir le
% style \lstset ci-dessous, à côté des autres boîtes pédagogiques).
\usepackage{listings}
% \degree : commande fréquemment utilisée par les modèles pour un angle ou une
% température en dehors de \SI{}{} (non fournie par les paquets chargés).
\providecommand{\degree}{^{\circ}}
% \qed (fin de démonstration), utilisé par les modèles sans amsthm
\providecommand{\qed}{\hfill$\square$}
% \barre : double barre des valeurs interdites dans un tableau de variations (tkz-tab)
\providecommand{\barre}{\ensuremath{\|}}
% environnement proof (amsthm non chargé) utilisé par les modèles
\ifdefined\proof\else\newenvironment{proof}[1][Démonstration]{\par\noindent\textit{#1.}\ }{\qed\par}\fi
\usepackage{lastpage}
\usepackage{hyperref}
\hypersetup{hidelinks}

% ============================================================
% Palette « Pop » OnBuch+ — mêmes hex que la classe P (plus_ui.dart)
% ============================================================
\definecolor{poporange}{HTML}{F59321}
\definecolor{popdark}{HTML}{E07A0C}
\definecolor{popcream}{HTML}{FFF6E8}
\definecolor{popink}{HTML}{1C1714}
\definecolor{poppurple}{HTML}{7A5AE0}
\definecolor{popgreen}{HTML}{2E9E6A}
\definecolor{poppink}{HTML}{E0457B}
\definecolor{popblue}{HTML}{2F7DE1}
\definecolor{popgold}{HTML}{C98A12}
\definecolor{popmuted}{HTML}{8A7F76}
\definecolor{popline}{HTML}{EADFD2}
\definecolor{popgray}{HTML}{8A7F76}
\colorlet{popgrey}{popgray}
% Alias tolérants (le modèle nomme parfois les couleurs différemment)
\colorlet{popwhite}{white}
\colorlet{popblack}{popink}
\colorlet{popred}{poppink}
\colorlet{popyellow}{popgold}
% Teintes claires pour fonds de tableaux / zones
\colorlet{popblueL}{popblue!10!white}
\colorlet{poporangeL}{poporange!14!white}
\colorlet{popgreenL}{popgreen!12!white}
\colorlet{poppurpleL}{poppurple!12!white}
\colorlet{poppinkL}{poppink!12!white}
\colorlet{popgoldL}{popgold!14!white}

\pagecolor{popcream}

% ============================================================
% Typographie
% ============================================================
\defaultfontfeatures{Ligatures=TeX}
\setmainfont{PlusJakartaSans-Medium.ttf}[Path=../../../fonts/,BoldFont=PlusJakartaSans-Bold.ttf]
% Maths en sans-serif (Fira Math) avec chiffres et lettres droites de Plus
% Jakarta, pour que les formules se fondent dans le texte.
\usepackage[math-style=ISO,bold-style=ISO]{unicode-math}
\setmathfont{FiraMath-Regular.otf}[Path=../../../fonts/]
\setmathfont{PlusJakartaSans-Medium.ttf}[Path=../../../fonts/,range={up/num,up/{Latin,latin},\mathup}]
\usepackage{icomma} % virgule décimale sans espace en mode math
% Caractères mathématiques Unicode absents des polices de texte (Plus Jakarta,
% Archivo Black) : rendus par leur équivalent mathématique, en texte comme en
% formule — sinon ils s'affichent en carré vide (ex. « Arithmétique dans ℤ »).
\usepackage{newunicodechar}
\newunicodechar{ℝ}{\ensuremath{\mathbb{R}}}
\newunicodechar{ℤ}{\ensuremath{\mathbb{Z}}}
\newunicodechar{ℂ}{\ensuremath{\mathbb{C}}}
\newunicodechar{ℕ}{\ensuremath{\mathbb{N}}}
\newunicodechar{ℚ}{\ensuremath{\mathbb{Q}}}
\newunicodechar{𝔻}{\ensuremath{\mathbb{D}}}
\newunicodechar{α}{\ensuremath{\alpha}}
\newunicodechar{β}{\ensuremath{\beta}}
\newunicodechar{γ}{\ensuremath{\gamma}}
\newunicodechar{δ}{\ensuremath{\delta}}
\newunicodechar{ε}{\ensuremath{\varepsilon}}
\newunicodechar{θ}{\ensuremath{\theta}}
\newunicodechar{λ}{\ensuremath{\lambda}}
\newunicodechar{μ}{\ensuremath{\mu}}
\newunicodechar{σ}{\ensuremath{\sigma}}
\newunicodechar{τ}{\ensuremath{\tau}}
\newunicodechar{φ}{\ensuremath{\varphi}}
\newunicodechar{ψ}{\ensuremath{\psi}}
\newunicodechar{ω}{\ensuremath{\omega}}
\newunicodechar{Σ}{\ensuremath{\Sigma}}
% majuscules produites par \MakeUppercase dans les titres (α-aminés -> Α)
\newunicodechar{Α}{\ensuremath{\alpha}}
\newunicodechar{Β}{\ensuremath{\beta}}
\newunicodechar{Γ}{\ensuremath{\Gamma}}
\newunicodechar{Δ}{\ensuremath{\Delta}}
\newunicodechar{Ε}{\ensuremath{\varepsilon}}
\newunicodechar{Θ}{\ensuremath{\Theta}}
\newunicodechar{Λ}{\ensuremath{\Lambda}}
\newunicodechar{Μ}{\ensuremath{\mu}}
\newunicodechar{Τ}{\ensuremath{\tau}}
\newunicodechar{Φ}{\ensuremath{\Phi}}
\newunicodechar{Ψ}{\ensuremath{\Psi}}
\newunicodechar{Ω}{\ensuremath{\Omega}}
\newunicodechar{↦}{\ensuremath{\mapsto}}
\newunicodechar{∈}{\ensuremath{\in}}
\newunicodechar{∉}{\ensuremath{\notin}}
\newunicodechar{∧}{\ensuremath{\wedge}}
\newunicodechar{⇔}{\ensuremath{\Leftrightarrow}}
\newunicodechar{⟺}{\ensuremath{\Longleftrightarrow}}
\newunicodechar{⇒}{\ensuremath{\Rightarrow}}
\newunicodechar{⟹}{\ensuremath{\Longrightarrow}}
\newunicodechar{∘}{\ensuremath{\circ}}
\newunicodechar{∪}{\ensuremath{\cup}}
\newunicodechar{∩}{\ensuremath{\cap}}
\newunicodechar{≡}{\ensuremath{\equiv}}
\newunicodechar{⊂}{\ensuremath{\subset}}
\newunicodechar{⊥}{\ensuremath{\perp}}
\newunicodechar{∃}{\ensuremath{\exists}}
\newunicodechar{∀}{\ensuremath{\forall}}
\newunicodechar{∛}{\ensuremath{\sqrt[3]{\,}}}
\newunicodechar{∜}{\ensuremath{\sqrt[4]{\,}}}
\newunicodechar{⃗}{\ensuremath{{}^{\to}}}
\newunicodechar{ⁿ}{\ifmmode^{n}\else\textsuperscript{\ensuremath{n}}\fi}
\newunicodechar{ˣ}{\ifmmode^{x}\else\textsuperscript{\ensuremath{x}}\fi}
\newunicodechar{ᵇ}{\ifmmode^{b}\else\textsuperscript{\ensuremath{b}}\fi}
\newunicodechar{ʳ}{\ifmmode^{r}\else\textsuperscript{\ensuremath{r}}\fi}
\newunicodechar{ᵘ}{\ifmmode^{u}\else\textsuperscript{\ensuremath{u}}\fi}
\newunicodechar{ʸ}{\ifmmode^{y}\else\textsuperscript{\ensuremath{y}}\fi}
\newunicodechar{ᵃ}{\ifmmode^{a}\else\textsuperscript{\ensuremath{a}}\fi}
\newunicodechar{ᵉ}{\ifmmode^{e}\else\textsuperscript{\ensuremath{e}}\fi}
\newunicodechar{ᵏ}{\ifmmode^{k}\else\textsuperscript{\ensuremath{k}}\fi}
\newunicodechar{ᵖ}{\ifmmode^{p}\else\textsuperscript{\ensuremath{p}}\fi}
\newunicodechar{ᴺ}{\ifmmode^{N}\else\textsuperscript{\ensuremath{N}}\fi}
\newunicodechar{ᶜ}{\ifmmode^{c}\else\textsuperscript{\ensuremath{c}}\fi}
\newunicodechar{ʰ}{\ifmmode^{h}\else\textsuperscript{\ensuremath{h}}\fi}
\newunicodechar{ᵠ}{\ifmmode^{q}\else\textsuperscript{\ensuremath{q}}\fi}
\newunicodechar{ₙ}{\ifmmode_{n}\else\textsubscript{\ensuremath{n}}\fi}
\newunicodechar{ₐ}{\ifmmode_{a}\else\textsubscript{\ensuremath{a}}\fi}
\newunicodechar{ᵢ}{\ifmmode_{i}\else\textsubscript{\ensuremath{i}}\fi}
\newunicodechar{ᵣ}{\ifmmode_{r}\else\textsubscript{\ensuremath{r}}\fi}
\newunicodechar{ₖ}{\ifmmode_{k}\else\textsubscript{\ensuremath{k}}\fi}
\newunicodechar{ᵦ}{\ifmmode_{\beta}\else\textsubscript{\ensuremath{\beta}}\fi}
\newunicodechar{ₓ}{\ifmmode_{x}\else\textsubscript{\ensuremath{x}}\fi}
\newfontfamily\popdisplay{ArchivoBlack-Regular.ttf}[Path=../../../fonts/]
\newfontfamily\poplogo{SpaceGrotesk-Bold.ttf}[Path=../../../fonts/]
\newfontfamily\popsemi{PlusJakartaSans-SemiBold.ttf}[Path=../../../fonts/]
\newfontfamily\popxbold{PlusJakartaSans-ExtraBold.ttf}[Path=../../../fonts/]
\newfontfamily\popxboldls{PlusJakartaSans-ExtraBold.ttf}[Path=../../../fonts/,LetterSpace=3.0]

\setlist[enumerate]{leftmargin=1.7em,itemsep=5pt,topsep=4pt}
\setlist[itemize]{leftmargin=1.7em,itemsep=4pt,topsep=4pt,label={\color{poporange}\textbullet}}
\setlength{\parindent}{0pt}
\setlength{\parskip}{5pt}
\renewcommand{\arraystretch}{1.25}

% pgfplots : style Pop par défaut
\pgfplotsset{
  every axis/.append style={axis line style={popink,line width=1pt},tick style={popink},
    grid style={popline,line width=0.5pt},label style={font=\small},tick label style={font=\footnotesize}},
  popcurve/.style={poporange,line width=1.8pt,no markers,smooth},
}
% Même style disponible dans un \draw TikZ simple (hors pgfplots)
\tikzset{popcurve/.style={poporange,line width=1.8pt}}
\tikzset{popgrid/.style={popline,very thin}}
\colorlet{popcurve}{poporange} % le nom du style est parfois utilisé comme couleur

% ============================================================
% Pill / squiggle
% ============================================================
\newcommand{\poppill}[3]{%
  \tikz[baseline=-0.55ex]\node[draw=popink,line width=1.2pt,rounded corners=2.6pt,%
    fill=#2,inner xsep=6.5pt,inner ysep=3pt]{%
    \popxboldls\fontsize{7}{7}\selectfont\color{#3}\MakeUppercase{#1}};%
}
\newcommand{\popsquiggle}[1]{%
  \tikz[baseline=-0.4ex]{
    \draw[#1,line width=2.6pt,line cap=round]
      plot [smooth] coordinates {(0,0.09) (0.34,0.01) (0.68,0.21) (1.02,0.01) (1.36,0.10)};
  }%
}
% Mot-clé surligné (marqueur orange sous le texte)
\newcommand{\cle}[1]{{\popxbold\color{popdark}#1}}

% ============================================================
% En-tête / pied
% ============================================================
\pagestyle{fancy}
\fancyhf{}
\renewcommand{\headrulewidth}{0pt}
\renewcommand{\footrulewidth}{0pt}
\lhead{\raisebox{-0.15ex}{\poplogo\fontsize{13}{13}\selectfont\color{popink}OnBuch\color{poporange}+}}
\rhead{\poppill{\DOCMATIERE\ \textbullet\ \DOCNIVEAU\ \textbullet\ Leçon \DOCLECON}{white}{popink}}
\lfoot{\popxbold\fontsize{7.6}{7.6}\selectfont\color{popmuted}\MakeUppercase{Cours signé OnBuch+}\ \textbullet\ {\popsemi\DOCTITRE}}
\rfoot{\tikz[baseline=-0.6ex]\node[rounded corners=8.5pt,fill=popink,minimum height=17pt,minimum width=17pt,inner xsep=5pt,inner sep=0]{\popxbold\fontsize{8}{8}\selectfont\color{popcream}\thepage\,/\,\pageref*{LastPage}};}

% ============================================================
% Titres
% ============================================================
\newcommand{\chaptitre}[2]{%
  \begin{tcolorbox}[enhanced,colback=poporange,colframe=popink,boxrule=2.6pt,arc=11pt,
      drop shadow=popink,left=15pt,right=15pt,top=12pt,bottom=11pt,before skip=3pt,after skip=18pt]
    \poppill{#2}{popink}{popcream}\par\vspace{9pt}
    {\raggedright\hyphenpenalty=10000\exhyphenpenalty=10000\popdisplay\fontsize{22}{24}\selectfont\color{popcream}\MakeUppercase{#1}\par}
  \end{tcolorbox}
}
% Section numérotée : pastille orange avec le numéro + titre Archivo
\newcounter{popsec}
\newcommand{\coursec}[1]{%
  \stepcounter{popsec}\setcounter{popsub}{0}%
  \par\vspace{16pt}\noindent
  \begin{minipage}{\linewidth}
  \tikz[baseline=-0.7ex]\node[draw=popink,line width=1.6pt,fill=poporange,rounded corners=4pt,minimum size=22pt,inner sep=0]{\popdisplay\fontsize{12}{12}\selectfont\color{popcream}\thepopsec};%
  \hspace{8pt}{\popdisplay\fontsize{15}{17}\selectfont\color{popink}\MakeUppercase{#1}}\par
  \vspace{4pt}\noindent{\color{poporange}\rule{36pt}{3.6pt}}
  \end{minipage}\par\nopagebreak\vspace{8pt}
}
\newcounter{popsub}
\newcommand{\courssub}[1]{%
  \stepcounter{popsub}%
  \par\vspace{9pt}\noindent{\popxbold\fontsize{11.5}{13}\selectfont\color{popink}{\color{poporange}\thepopsec.\thepopsub}\hspace{6pt}#1}\par\nopagebreak\vspace{4pt}
}

% ============================================================
% Boîtes pédagogiques — même recette : fond blanc, bordure épaisse colorée,
% ombre dure encre, badge plein. Titre optionnel [#1].
% ============================================================
\tcbset{popbase/.style={breakable,enhanced,colback=white,boxrule=2.3pt,arc=9pt,
  shadow={3pt}{-3pt}{0pt}{popink},left=14pt,right=14pt,top=11pt,bottom=11pt,before skip=11pt,after skip=15pt,
  fontupper=\popsemi\fontsize{10.6}{14.5}\selectfont\color{popink},
  fontlower=\popsemi\fontsize{10.6}{14.5}\selectfont\color{popink}}}
% Coche dessinée (le glyphe ✓ n'existe pas dans les polices du document)
\newcommand{\popcheck}[1]{\tikz[baseline=-0.5ex]\draw[#1,line width=1.6pt,line cap=round,line join=round](-0.13,0.0)--(-0.04,-0.09)--(0.13,0.1);}
\providecommand{\coche}{\popcheck{popgreen}}
\providecommand{\resultat}[1]{\par\smallskip\noindent\fcolorbox{popgreen}{popgreenL}{\parbox{\dimexpr\linewidth-2\fboxsep-2\fboxrule\relax}{\textbf{Résultat.} #1}}\par\smallskip}
\newcommand{\popboxhead}[3]{\poppill{#1}{#2}{white}\ifx\relax#3\relax\else\hspace{6pt}{\popxbold\fontsize{10.6}{12}\selectfont\color{popink}#3}\fi\par\vspace{7pt}}

\newtcolorbox{aretenir}[1][]{popbase,colframe=popgreen,before upper={\popboxhead{À retenir}{popgreen}{#1}}}
\newtcolorbox{exemplebox}[1][]{popbase,colframe=poporange,before upper={\popboxhead{Exemple}{poporange}{#1}}}
\newtcolorbox{definition}[1][]{popbase,colframe=popblue,before upper={\popboxhead{Définition}{popblue}{#1}}}
\newtcolorbox{propriete}[1][]{popbase,colframe=poppurple,before upper={\popboxhead{Propriété}{poppurple}{#1}}}
\newtcolorbox{methode}[1][]{popbase,colframe=popgold,before upper={\popboxhead{Méthode}{popgold}{#1}}}
\newtcolorbox{attention}[1][]{popbase,colframe=poppink,before upper={\popboxhead{Attention}{poppink}{#1}}}
\newtcolorbox{experience}[1][]{popbase,colframe=popblue,colback=popblueL,before upper={\popboxhead{Expérience}{popblue}{#1}}}
% « Le savais-tu ? » : carte violette pleine (contexte, culture, Cameroun)
\newtcolorbox{savaistu}[1][]{popbase,colback=poppurple,colframe=popink,
  fontupper=\popsemi\fontsize{10.4}{14.2}\selectfont\color{white},
  before upper={\poppill{Le savais-tu\,?}{white}{poppurple}\ifx\relax#1\relax\else\hspace{6pt}{\popxbold\color{white}#1}\fi\par\vspace{7pt}}}
% Exercice résolu : énoncé en haut, \tcblower puis la solution sur fond crème
\newcounter{popexo}\newcounter{popexr}
\newtcolorbox{exoresolu}[1][]{popbase,colframe=poporange,colbacklower=popcream,
  segmentation style={popink,line width=1.2pt,dashed},
  before upper={\stepcounter{popexr}\popboxhead{Exercice résolu \thepopexr}{poporange}{#1}},
  before lower={\poppill{Solution}{popink}{popcream}\par\vspace{6pt}}}
% Exercice (non résolu en place) — la correction est dans la partie Corrigés
\newtcolorbox{exercice}[1][]{popbase,colframe=popink,
  before upper={\stepcounter{popexo}\popboxhead{Exercice \thepopexo}{popink}{#1}}}
% Corrigé
\newtcolorbox{corrige}[1][]{popbase,colframe=popgreen,colback=popgreenL,
  before upper={\popboxhead{Corrigé}{popgreen}{#1}}}
% Objectifs / prérequis / bilan
\newtcolorbox{objectifs}{popbase,colframe=popink,colback=poporangeL,
  before upper={\popboxhead{Objectifs de la leçon}{popink}{}}}
\newtcolorbox{prerequis}{popbase,colframe=popink,colback=popblueL,
  before upper={\popboxhead{Prérequis}{popblue}{}}}
\newtcolorbox{bilan}{popbase,colframe=popink,colback=white,boxrule=2.8pt,
  before upper={\popboxhead{Fiche bilan}{poporange}{L'essentiel à connaître pour l'examen}}}
% Figure encadrée avec légende
\newtcolorbox{popfigure}[1][]{enhanced,breakable=false,colback=white,colframe=popline,boxrule=1.4pt,arc=8pt,
  left=10pt,right=10pt,top=10pt,bottom=8pt,before skip=10pt,after skip=12pt,halign=center,
  lower separated=false,halign lower=center,
  fontlower=\popsemi\fontsize{9}{11.5}\selectfont\color{popmuted},
  #1}
\newcommand{\legende}[1]{\par\vspace{6pt}{\popsemi\fontsize{9}{11.5}\selectfont\color{popmuted}#1}}

% ============================================================
% Bloc de code source — \begin{lstlisting}[language=Php]...\end{lstlisting}
% (ou language=C, HTML, JavaScript, SQL ; option omise = pseudo-code brut).
% Même recette visuelle que les figures (\popfigure) : cadre encre épais,
% fond clair (popcream), légende possible juste après via \legende{...}
% comme pour une figure (listings ne sait pas arrondir ses coins : on garde
% un cadre rectangulaire, cohérent avec les autres tableaux du document).
% CONTRAT : les modèles ne doivent utiliser QUE \begin{lstlisting}[...] tel
% quel (pas de \begin{tcblisting}, pas de \verb pour un bloc multi-lignes).
% ============================================================
\lstdefinestyle{poplst}{
  backgroundcolor=\color{popcream},
  basicstyle=\popsemi\fontsize{8.6}{11}\selectfont\color{popink},
  keywordstyle=\bfseries\color{popblue},
  commentstyle=\itshape\color{popmuted},
  stringstyle=\color{popgreen},
  numberstyle=\tiny\color{popmuted},
  identifierstyle=\color{popink},
  frame=l,
  rulecolor=\color{popink},
  framerule=1.6pt,
  framesep=7pt,
  framexleftmargin=7pt,
  framexrightmargin=7pt,
  xleftmargin=2pt,
  xrightmargin=2pt,
  breaklines=true,
  breakatwhitespace=false,
  showstringspaces=false,
  showspaces=false,
  showtabs=false,
  tabsize=2,
  columns=flexible,
  keepspaces=true,
  numbers=none,
  aboveskip=10pt,
  belowskip=10pt,
  captionpos=b,
}
\lstset{style=poplst, language=PHP}
% La distribution TeX embarquée ne fournit pas le fichier de langue
% JavaScript de listings (lstlang*.dat) : on la redéfinit nous-mêmes
% pour que \begin{lstlisting}[language=JavaScript] fonctionne.
\lstdefinelanguage{JavaScript}{
  keywords={break,case,catch,class,const,continue,debugger,default,delete,do,
    else,export,extends,finally,for,function,if,import,in,instanceof,let,new,
    return,static,super,switch,this,throw,try,typeof,var,void,while,with,yield,
    async,await,of,get,set},
  keywordstyle=\bfseries\color{popblue},
  ndkeywords={true,false,null,undefined,NaN,Infinity},
  ndkeywordstyle=\bfseries\color{poporange},
  identifierstyle=\color{popink},
  sensitive=true,
  comment=[l]{//},
  morecomment=[s]{/*}{*/},
  string=[b]",
  morestring=[b]',
  morestring=[b]`,
}
% Même limitation pour CSS et les fichiers de configuration (apacheconf,
% nginx, ini...) : ils ne sont pas fournis. CSS et un style générique de
% type "config" (commentaires # ou ;, pas de mots-clés) couvrent les besoins.
\lstdefinelanguage{CSS}{
  keywords={color,background,background-color,margin,padding,border,
    border-radius,font,font-size,font-weight,font-family,width,height,
    display,position,top,left,right,bottom,float,clear,text-align,
    line-height,list-style,cursor,overflow,box-shadow,transition,
    flex,grid,align-items,justify-content,z-index},
  keywordstyle=\bfseries\color{popblue},
  identifierstyle=\color{popink},
  sensitive=false,
  comment=[s]{/*}{*/},
  string=[b]",
  morestring=[b]',
}
\lstdefinelanguage{apacheconf}{
  sensitive=false,
  morecomment=[l]{\#},
}
\lstdefinelanguage{Apache}{
  sensitive=false,
  morecomment=[l]{\#},
}
\lstdefinelanguage{text}{sensitive=false}
\lstdefinelanguage{powershell}{
  keywords={New-Object,New-SmbShare,Get-Acl,Set-Acl,Get-ChildItem,Set-Content,
    Get-Content,Write-Host,ForEach-Object,Where-Object,Select-Object,
    Test-Path,Remove-Item,Copy-Item,Move-Item,Start-Process,Stop-Process,
    If,Else,ElseIf,ForEach,While,Function,Return,Try,Catch},
  keywordstyle=\bfseries\color{popblue},
  identifierstyle=\color{popink},
  sensitive=false,
  comment=[l]{\#},
  string=[b]",
  morestring=[b]',
}
\lstdefinelanguage{ini}{
  sensitive=false,
  morecomment=[l]{;},
  morecomment=[l]{\#},
}

% ============================================================
% Signature OnBuch+
% ============================================================
\newcommand{\poplogoinline}[1]{{\poplogo\fontsize{#1}{#1}\selectfont\color{popink}OnBuch\color{poporange}+}}
\newcommand{\signatureonbuch}{%
  \begin{tcolorbox}[enhanced,colback=popink,colframe=popink,boxrule=2.2pt,arc=10pt,
      drop shadow=poporange,left=14pt,right=14pt,top=10pt,bottom=10pt,before skip=0pt,after skip=0pt]
    \tikz[baseline=-0.7ex]\node[circle,fill=popgreen,draw=popcream,line width=1.3pt,minimum size=22pt,inner sep=0]{\popcheck{white}};%
    \hspace{9pt}\begin{minipage}[c]{0.62\linewidth}
      {\popxbold\fontsize{12}{13}\selectfont\color{popcream}Signé\ \textbullet\ {\poplogo OnBuch\color{poporange}+}}\par\vspace{2pt}
      {\raggedright\popsemi\fontsize{8}{10.5}\selectfont\color{popline}Ludovic A.\\\MakeUppercase{Conforme au programme officiel MINESEC}\par}
    \end{minipage}\hfill
    {\poplogo\fontsize{22}{22}\selectfont\color{popcream}OnBuch\color{poporange}+}
  \end{tcolorbox}%
}

% ============================================================
% Page de garde
% #1 titre · #2 matière · #3 niveau · #4 module · #5 n° leçon · #6 durée ·
% #7 description / objectifs (texte ou itemize)
% ============================================================
\newcommand{\pagegarde}[7]{%
  \thispagestyle{empty}
  \newgeometry{a4paper,margin=1.7cm,top=1.6cm,bottom=1.4cm}%
  \begin{tikzpicture}[remember picture,overlay]
    \fill[poporange!18] ([xshift=-3.2cm,yshift=-3.6cm]current page.north east) circle (4.6cm);
    \fill[poppurple!14] ([xshift=2.4cm,yshift=7.6cm]current page.south west) circle (3.8cm);
  \end{tikzpicture}%
  {\poplogo\fontsize{30}{30}\selectfont\color{popink}OnBuch\color{poporange}+}\hfill
  \raisebox{0.6ex}{\poppill{Cours signé \textbullet\ Premium}{popink}{popcream}}\par
  \vspace{1pt}\popsquiggle{poppurple}\par
  \vspace{8pt}
  \begin{tcolorbox}[enhanced,colback=poporange,colframe=popink,boxrule=2.8pt,arc=13pt,
      drop shadow=popink,left=18pt,right=18pt,top=12pt,bottom=13pt,after skip=10pt]
    \poppill{#2 \textbullet\ #3}{popink}{popcream}\hspace{5pt}\poppill{#4}{white}{popink}\par\vspace{8pt}
    {\popdisplay\fontsize{14}{14}\selectfont\color{popink}LEÇON #5}\par\vspace{4pt}
    {\raggedright\hyphenpenalty=10000\exhyphenpenalty=10000\popdisplay\fontsize{25}{27}\selectfont\color{popcream}\MakeUppercase{#1}\par}
    \vspace{8pt}
    {\popxbold\fontsize{9.5}{9.5}\selectfont\color{popink}\MakeUppercase{Durée conseillée : #6}}
  \end{tcolorbox}
  \begin{tcolorbox}[enhanced,colback=white,colframe=popink,boxrule=2.2pt,arc=10pt,
      drop shadow=popink,left=16pt,right=16pt,top=10pt,bottom=10pt,after skip=8pt]
    \poppill{Dans cette leçon}{poppurple}{white}\par\vspace{5pt}
    \popsemi\fontsize{9.3}{12.6}\selectfont\color{popink}#7
  \end{tcolorbox}
  \noindent\begin{minipage}[t]{0.487\linewidth}
    \begin{tcolorbox}[enhanced,colback=popgreen,colframe=popink,boxrule=2.2pt,arc=10pt,
        drop shadow=popink,left=11pt,right=11pt,top=10pt,bottom=10pt,height=3.1cm,valign=center]
      \tcbox[colback=white,colframe=popink,boxrule=1.3pt,arc=5pt,boxsep=0pt,nobeforeafter,
        left=3pt,right=3pt,top=3pt,bottom=3pt,valign=center]{\qrcode[height=1.9cm]{https://play.google.com/store/apps/details?id=com.luvvix.onbuch&hl=fr}}%
      \hspace{8pt}\begin{minipage}[c]{0.5\linewidth}
        {\popdisplay\fontsize{11}{12.5}\selectfont\color{white}\MakeUppercase{Télécharge OnBuch}}\par\vspace{4pt}
        {\raggedright\popsemi\fontsize{8.4}{11}\selectfont\color{white}Gratuit \textbullet\ Google Play\\Tuteur IA, annales, concours, résultats\par}
      \end{minipage}
    \end{tcolorbox}
  \end{minipage}\hfill
  \begin{minipage}[t]{0.487\linewidth}
    \begin{tcolorbox}[enhanced,colback=poppurple,colframe=popink,boxrule=2.2pt,arc=10pt,
        drop shadow=popink,left=11pt,right=11pt,top=10pt,bottom=10pt,height=3.1cm,valign=center]
      \tikz[baseline=-0.7ex]\node[circle,fill=white,draw=popink,line width=1.3pt,minimum size=1.6cm,inner sep=0]{\popdisplay\fontsize{24}{24}\selectfont\color{poppurple}+};%
      \hspace{8pt}\begin{minipage}[c]{0.6\linewidth}
        {\popdisplay\fontsize{11}{12.5}\selectfont\color{white}\MakeUppercase{Passe à OnBuch+}}\par\vspace{4pt}
        {\raggedright\popsemi\fontsize{8.4}{11}\selectfont\color{white}Tous les cours signés, Tuteur IA illimité, fiches PDF — onglet OnBuch+ de l'app\par}
      \end{minipage}
    \end{tcolorbox}
  \end{minipage}
  \vspace{8pt}
  \popcontenu
  \vfill
  \signatureonbuch
  \restoregeometry
  \newpage
}

% Rangée « Ce cours contient » (page de garde)
\newcommand{\poptile}[3]{% #1 couleur #2 gros texte #3 légende
  \begin{tcolorbox}[enhanced,colback=#1,colframe=popink,boxrule=2pt,arc=9pt,shadow={2.5pt}{-2.5pt}{0pt}{popink},
      left=6pt,right=6pt,top=9pt,bottom=9pt,halign=center,valign=center,height=1.75cm,width=\linewidth,nobeforeafter]
    {\popdisplay\fontsize{12.5}{13}\selectfont\color{white}\MakeUppercase{#2}}\par\vspace{3pt}
    {\centering\popsemi\fontsize{7.8}{9.5}\selectfont\color{white}#3\par}
  \end{tcolorbox}}
\newcommand{\popcontenu}{%
  \par\noindent{\popxboldls\fontsize{7.5}{7.5}\selectfont\color{popmuted}CE COURS SIGNÉ CONTIENT}\par\vspace{7pt}
  \noindent\begin{minipage}[t]{0.235\linewidth}\poptile{popblue}{Cours}{complet, illustré, pas à pas}\end{minipage}\hfill
  \begin{minipage}[t]{0.235\linewidth}\poptile{poporange}{Méthodes}{et exercices résolus}\end{minipage}\hfill
  \begin{minipage}[t]{0.235\linewidth}\poptile{poppink}{Exercices}{type examen + corrigés}\end{minipage}\hfill
  \begin{minipage}[t]{0.235\linewidth}\poptile{popgreen}{Fiche bilan}{l'essentiel à retenir}\end{minipage}\par}

% Sommaire
\newtcolorbox{sommaire}{enhanced,colback=white,colframe=popink,boxrule=2.2pt,arc=10pt,
  drop shadow=popink,left=18pt,right=18pt,top=13pt,bottom=13pt,before skip=6pt,after skip=20pt,
  before upper={\poppill{Sommaire}{poppurple}{white}\par\vspace{9pt}\popsemi\fontsize{10.6}{14.5}\selectfont\color{popink}}}

% Signature de fin de cours — poussée tout en bas de la dernière page
\newcommand{\findecours}{%
  \par\vfill\nopagebreak
  \begin{center}\popsquiggle{poporange}\end{center}\vspace{4pt}
  \signatureonbuch
}
\AtBeginDocument{\def\llbracket{\mathopen{[\mkern-3.5mu[}}\def\rrbracket{\mathclose{]\mkern-3.5mu]}}} % intervalles d'entiers (glyphe ⟦ absent de la police)
% Glyphes absents de Fira Math : redéfinis à partir de glyphes disponibles
\AtBeginDocument{%
  \def\setminus{\mathbin{\backslash}}\let\smallsetminus\setminus
  \def\vdots{\mathord{\vbox{\baselineskip3.2pt\lineskiplimit0pt\kern4pt\hbox{.}\hbox{.}\hbox{.}}}}%
  \def\ddots{\mathinner{\mkern1mu\raise6.5pt\hbox{.}\mkern2mu\raise3.6pt\hbox{.}\mkern2mu\raise0.7pt\hbox{.}\mkern1mu}}%
  \def\triangle{\mathord{\scalebox{0.9}{$\symup{\Delta}$}}}%
  \def\nabla{\mathord{\raisebox{\depth}{\scalebox{1}[-1]{$\symup{\Delta}$}}}}%
  \def\checkmark{\popcheck{popdark}}%
  \def\star{\mathbin{\ast}}\def\bigstar{\mathord{\ast}}%
  \def\diamond{\mathbin{\scalebox{0.6}{$\Diamond$}}}%
  \def\bigcup{\mathop{\mathchoice{\raisebox{-0.3em}{\scalebox{1.7}{$\cup$}}}{\scalebox{1.25}{$\cup$}}{\cup}{\cup}}\displaylimits}%
  \def\bigcap{\mathop{\mathchoice{\raisebox{-0.3em}{\scalebox{1.7}{$\cap$}}}{\scalebox{1.25}{$\cap$}}{\cap}{\cap}}\displaylimits}%
  \def\lesseqqgtr{\mathrel{\lessgtr}}\def\bigoplus{\mathop{\oplus}\displaylimits}%
}
\providecommand{\ssi}{si et seulement si} % abréviation fréquente
\AtBeginDocument{\providecommand{\wideparen}{\overparen}} % arc de cercle

%% FIGFIX-BEGIN v4 — correctifs globaux des figures (docs/onbuch-generation/tools/figfix)
\usepackage{adjustbox}\usepackage{etoolbox}
% toute figure TikZ/pgfplots plus large que la ligne est réduite à la largeur disponible
\BeforeBeginEnvironment{tikzpicture}{\begin{adjustbox}{max width=\linewidth}}
\AfterEndEnvironment{tikzpicture}{\end{adjustbox}}
% légendes pgfplots lisibles par-dessus les courbes
\pgfplotsset{every axis/.append style={legend style={fill=white,fill opacity=0.92,text opacity=1,draw=black!15,inner sep=3pt}}}
% étiquettes d'arêtes (syntaxe edge["..."]) avec fond blanc
\tikzset{every edge quotes/.append style={fill=white,inner sep=1.2pt}}
%% FIGFIX-END
\begin{document}

\pagegarde{Utilisation des structures de contrôle}{Informatique}{3ème}{Informatique – classe de 3ème}{N14}{3 séances (fiche de progression harmonisée)}{Cette leçon introduit l'écriture d'algorithmes simples et l'utilisation des structures alternatives (si... alors... sinon) et itératives (tant que, pour) pour résoudre des problèmes concrets. Elle présente également l'algorigramme, représentation graphique normalisée d'un algorithme.
\vspace{4pt}
\begin{multicols}{2}\small
\begin{itemize}
\item À la fin de cette leçon, je sais définir un algorithme et identifier ses constituants (variables, instructions, entrées, sorties)
\item À la fin de cette leçon, je sais écrire un algorithme simple en pseudo-code pour résoudre un problème de la vie courante
\item À la fin de cette leçon, je sais utiliser une structure alternative simple (si... alors) et complète (si... alors... sinon)
\item À la fin de cette leçon, je sais utiliser une structure itérative « tant que » et vérifier qu'elle se termine
\item À la fin de cette leçon, je sais utiliser une structure itérative « pour » pour répéter un nombre connu de fois
\item À la fin de cette leçon, je sais représenter un algorithme par un algorigramme en respectant les symboles normalisés
\end{itemize}\end{multicols}}

\begin{sommaire}
\begin{multicols}{2}
\begin{enumerate}
\item Notion d'algorithme et pseudo-code
\item Les structures alternatives
\item Les structures itératives : la boucle « tant que »
\item Les structures itératives : la boucle « pour »
\item L'algorigramme : représentation graphique
\item Résolution de problèmes concrets et justification
\item Activité d'intégration
\item Méthodes et exercices résolus
\item Exercices
\item Corrigés des exercices
\item Fiche bilan
\end{enumerate}
\end{multicols}
\end{sommaire}

\begin{objectifs}
\begin{itemize}
\item À la fin de cette leçon, je sais définir un algorithme et identifier ses constituants (variables, instructions, entrées, sorties)
\item À la fin de cette leçon, je sais écrire un algorithme simple en pseudo-code pour résoudre un problème de la vie courante
\item À la fin de cette leçon, je sais utiliser une structure alternative simple (si... alors) et complète (si... alors... sinon)
\item À la fin de cette leçon, je sais utiliser une structure itérative « tant que » et vérifier qu'elle se termine
\item À la fin de cette leçon, je sais utiliser une structure itérative « pour » pour répéter un nombre connu de fois
\item À la fin de cette leçon, je sais représenter un algorithme par un algorigramme en respectant les symboles normalisés
\item À la fin de cette leçon, je sais dérouler (tracer) un algorithme à la main pour vérifier son fonctionnement
\item À la fin de cette leçon, je sais rédiger et justifier la démarche suivie pour résoudre un problème concret à l'aide d'un algorithme
\end{itemize}
\end{objectifs}

\begin{prerequis}
\begin{itemize}
\item Notion de variable et d'affectation (début d'année de 3ème)
\item Opérations arithmétiques et comparaisons (<, >, =, ≤, ≥, ≠)
\item Expressions logiques simples : conditions vraies ou fausses
\item Résolution de problèmes numériques simples (calcul de moyenne, de prix, de périmètre)
\end{itemize}
\end{prerequis}

\newpage

\coursec{Notion d'algorithme et pseudo-code}

\textbf{Situation d'accroche.} Maman Aïcha vend des beignets au marché de Mokolo, à Yaoundé. Chaque matin, sans même y réfléchir, elle répète les mêmes gestes : \emph{tant qu'il reste de la pâte}, elle forme un beignet et le fait frire ; \emph{si} un client achète pour plus de 500 FCFA, elle offre un beignet en plus ; sinon, elle dit simplement merci. Le soir, elle compte sa recette : elle additionne les ventes, soustrait le prix de la farine et de l'huile, et obtient son bénéfice.

Sans le savoir, Maman Aïcha exécute un \cle{algorithme} : une suite d'instructions précises, dans un ordre bien défini, avec des conditions (« si le client achète pour plus de 500 FCFA... ») et des répétitions (« tant qu'il reste de la pâte... »).

\textbf{Problématique.} Comment écrire de telles instructions de façon précise et sans ambiguïté, pour qu'une machine (ou une personne) les exécute sans erreur ? C'est tout l'objet de cette leçon. Dans cette première section, nous découvrons ce qu'est un algorithme, comment on le décrit en \cle{pseudo-code}, et comment vérifier qu'il fonctionne grâce à la \cle{trace d'exécution}.

\courssub{Qu'est-ce qu'un algorithme ?}

\textbf{Intuition.} Imagine que tu dois expliquer à un petit frère comment préparer un thé : « mets de l'eau dans la bouilloire, fais chauffer, mets le thé dans la tasse, verse l'eau chaude, attends 3 minutes ». Tu viens de donner un algorithme ! Chaque étape est claire, les étapes sont dans le bon ordre (on ne verse pas l'eau avant de l'avoir chauffée), et au bout d'un nombre fini d'étapes, le résultat est obtenu : le thé est prêt.

En informatique, c'est exactement la même idée : avant de demander à un ordinateur de faire un calcul, on décrit d'abord, en langage clair, la suite des opérations à effectuer.

\begin{definition}[Algorithme]
Un \cle{algorithme} est une \textbf{suite finie et ordonnée d'instructions} permettant de résoudre un problème ou d'obtenir un résultat à partir de données fournies.
\end{definition}

Décomposons cette définition :
\begin{itemize}
  \item \textbf{suite finie} : l'algorithme s'arrête toujours après un nombre fini d'étapes ;
  \item \textbf{ordonnée} : les instructions s'exécutent dans un ordre précis, qu'on ne peut pas changer n'importe comment ;
  \item \textbf{instructions} : chaque étape est un ordre clair et réalisable (calculer, lire, afficher, comparer...) ;
  \item \textbf{résoudre un problème} : l'algorithme part de \emph{données} et produit un \emph{résultat}.
\end{itemize}

\begin{savaistu}[D'où vient le mot « algorithme » ?]
Le mot « algorithme » vient du nom du mathématicien perse \textbf{Al-Khwarizmi}, qui vécut au IX\textsuperscript{e} siècle à Bagdad. Dans ses ouvrages, il décrivait des méthodes de calcul étape par étape (par exemple pour résoudre des équations). Son nom, latinisé en \emph{Algoritmi}, a donné le mot « algorithme ». C'est aussi de son livre sur le calcul avec les chiffres indiens que vient le mot « algèbre » (\emph{al-jabr}). Plus de mille ans plus tard, chaque application de ton téléphone exécute des algorithmes inspirés de cette idée !
\end{savaistu}

\courssub{Entrées, traitements, sorties : la « machine à transformer »}

Tout algorithme peut être vu comme une « machine à transformer » : on lui donne des \textbf{entrées}, il effectue des \textbf{traitements}, et il produit des \textbf{sorties}.

\begin{definition}[Entrées, traitements, sorties]
\begin{itemize}
  \item Les \cle{entrées} sont les données fournies à l'algorithme avant son exécution (par exemple : les deux notes d'un élève).
  \item Les \cle{traitements} sont les opérations effectuées sur ces données (par exemple : additionner les notes, diviser par 2).
  \item Les \cle{sorties} sont les résultats produits et affichés (par exemple : la moyenne de l'élève).
\end{itemize}
\end{definition}

Un bon algorithme possède trois qualités essentielles :

\begin{propriete}[Qualités d'un bon algorithme]
\begin{enumerate}
  \item \textbf{Finitude} : l'algorithme doit se terminer après un nombre fini d'étapes, quelles que soient les entrées.
  \item \textbf{Clarté (non-ambiguïté)} : chaque instruction doit avoir un seul sens possible. « Ajoute un peu de sel » n'est pas une instruction claire ; « ajoute 5 g de sel » l'est.
  \item \textbf{Efficacité} : l'algorithme doit produire le bon résultat, et de préférence sans étapes inutiles.
\end{enumerate}
\end{propriete}

\begin{exemplebox}[Des algorithmes dans la vie courante]
\begin{itemize}
  \item \textbf{Recette de cuisine} : entrées = les ingrédients ; traitements = laver, couper, mélanger, cuire ; sortie = le plat (par exemple un eru bien préparé).
  \item \textbf{Itinéraire} : entrées = le point de départ (Carrefour Wada) et la destination (le collège) ; traitements = la suite des directions ; sortie = l'arrivée au collège.
  \item \textbf{Retrait d'argent au guichet mobile (OM/MoMo)} : entrées = ton numéro, le montant, ton code secret ; traitements = vérifier le code, vérifier que le solde est suffisant, débiter le compte ; sorties = l'argent remis et le message de confirmation. Remarque l'ordre : on vérifie le code \emph{avant} de donner l'argent !
\end{itemize}
\end{exemplebox}

\courssub{Les variables : les « boîtes » de l'algorithme}

\textbf{Intuition.} Pour calculer une moyenne, l'algorithme doit \emph{mémoriser} les deux notes, puis la somme, puis la moyenne. Où stocke-t-on ces valeurs ? Dans des \textbf{variables}.

\begin{definition}[Variable]
Une \cle{variable} est une zone de la mémoire, repérée par un \textbf{nom} (appelé identifiant), qui contient une \textbf{valeur} pouvant changer au cours de l'exécution de l'algorithme. On peut imaginer une variable comme une boîte étiquetée : le nom est écrit sur la boîte, la valeur est rangée dedans.
\end{definition}

Par exemple, la variable nommée \texttt{note1} peut contenir la valeur 12. Plus tard, on pourra ranger 14 dans cette même boîte : l'ancienne valeur est alors remplacée.

Quelques règles pour nommer les variables :
\begin{itemize}
  \item le nom commence par une lettre et ne contient ni espace ni accent : on écrit \texttt{note1}, \texttt{moyenne}, \texttt{prixTotal} ;
  \item on choisit des noms qui ont un sens : \texttt{moyenne} est bien plus clair que \texttt{x} ou \texttt{truc}.
\end{itemize}

\courssub{Le pseudo-code : un langage pour écrire les algorithmes}

\textbf{Intuition.} Un ordinateur ne comprend que des langages de programmation très stricts. Mais avant de programmer, il faut d'abord \emph{penser} la solution. On utilise alors le \cle{pseudo-code} : un langage simplifié, en français, qui décrit l'algorithme sans se soucier des détails techniques d'un vrai langage. Le pseudo-code est compréhensible par tout le monde et pourra ensuite être traduit dans n'importe quel langage de programmation.

\begin{methode}[Structure générale d'un algorithme en pseudo-code]
Tout algorithme écrit en pseudo-code respecte la structure suivante, en quatre parties :
\begin{enumerate}
  \item \textbf{La ligne \texttt{Algorithme}} suivie du nom de l'algorithme ;
  \item \textbf{La déclaration des variables} : on annonce toutes les « boîtes » dont on aura besoin, avec leur type (nombre, texte...) ;
  \item \textbf{Le corps}, compris entre \texttt{Début} et \texttt{Fin} : c'est là qu'on écrit les instructions, dans l'ordre où elles seront exécutées ;
  \item \textbf{La ligne \texttt{Fin}} qui marque l'arrêt de l'algorithme.
\end{enumerate}
\end{methode}

\begin{propriete}[Syntaxe générale]
\begin{lstlisting}
Algorithme nom_de_l_algorithme
Variables
    nom1 : type
    nom2 : type
Debut
    instruction 1
    instruction 2
    ...
Fin
\end{lstlisting}
Les types usuels en classe de 3ème sont : \texttt{Entier} (nombre sans virgule), \texttt{Réel} (nombre à virgule), \texttt{Chaîne} (texte).
\end{propriete}

\courssub{Les instructions de base : lire, écrire, affecter}

Trois instructions suffisent pour écrire nos premiers algorithmes.

\textbf{1. L'instruction \texttt{Lire} : la saisie des données.} Elle permet à l'utilisateur d'entrer une valeur au clavier, qui sera rangée dans une variable. C'est l'instruction des \emph{entrées}.
\begin{center}
\texttt{Lire note1} \quad signifie : « demande une valeur à l'utilisateur et range-la dans la variable \texttt{note1} ».
\end{center}

\textbf{2. L'instruction \texttt{Écrire} : l'affichage des résultats.} Elle affiche à l'écran un message ou la valeur d'une variable. C'est l'instruction des \emph{sorties}.
\begin{center}
\texttt{Écrire "La moyenne est : ", moyenne} \quad affiche le texte puis la valeur contenue dans \texttt{moyenne}.
\end{center}

\textbf{3. L'affectation : ranger une valeur dans une variable.} On utilise le symbole $\leftarrow$ (une flèche vers la gauche) :
\begin{center}
\texttt{moyenne $\leftarrow$ (note1 + note2) / 2}
\end{center}
Cette instruction se lit : « calcule la valeur de \texttt{(note1 + note2) / 2}, puis range le résultat dans la variable \texttt{moyenne} ». On évalue d'abord ce qui est \emph{à droite} de la flèche, puis on range le résultat dans la variable \emph{à gauche}.

\begin{attention}[Affectation ou égalité mathématique ?]
L'erreur la plus fréquente des débutants est de confondre l'\textbf{affectation} \texttt{x $\leftarrow$ 5} avec l'\textbf{égalité mathématique} $x = 5$.
\begin{itemize}
  \item En mathématiques, $x = 5$ est une \emph{affirmation} : « $x$ vaut 5 ». On peut écrire $x = 5$ ou $5 = x$, c'est pareil.
  \item En algorithmique, \texttt{x $\leftarrow$ 5} est un \emph{ordre} : « range la valeur 5 dans la boîte x ». L'ordre compte : \texttt{5 $\leftarrow$ x} n'a aucun sens, car on ne peut pas ranger une valeur dans le nombre 5 !
\end{itemize}
Conséquence importante : l'instruction \texttt{x $\leftarrow$ x + 1} est parfaitement correcte en algorithmique (on prend l'ancienne valeur de \texttt{x}, on ajoute 1, et on range le résultat dans \texttt{x}), alors que $x = x + 1$ est impossible en mathématiques !
\end{attention}

\courssub{Premier algorithme complet : la moyenne de deux notes}

\textbf{Problème.} Le professeur d'informatique du collège veut un algorithme qui demande les deux notes d'un élève (sur 20), calcule la moyenne et l'affiche.

\textbf{Analyse du problème} (toujours avant d'écrire !) :
\begin{itemize}
  \item \textbf{Entrées} : deux notes, que nous appellerons \texttt{note1} et \texttt{note2} ;
  \item \textbf{Traitement} : calculer $\text{moyenne} = \dfrac{\text{note1} + \text{note2}}{2}$ ;
  \item \textbf{Sortie} : afficher la moyenne.
\end{itemize}

\begin{exemplebox}[Algorithme « Moyenne de deux notes »]
\begin{lstlisting}
Algorithme MoyenneDeuxNotes
Variables
    note1, note2 : Reel      // les deux notes saisies
    moyenne : Reel           // le resultat calcule
Debut
    Ecrire "Entrez la premiere note : "
    Lire note1                          // saisie de la 1re note
    Ecrire "Entrez la deuxieme note : "
    Lire note2                          // saisie de la 2e note
    moyenne <- (note1 + note2) / 2      // calcul de la moyenne
    Ecrire "La moyenne est : ", moyenne // affichage du resultat
Fin
\end{lstlisting}
\end{exemplebox}

Lisons cet algorithme ligne par ligne : on affiche d'abord un message pour guider l'utilisateur, on lit la première note, puis la deuxième ; ensuite on calcule la moyenne par affectation ; enfin on affiche le résultat. L'ordre est essentiel : on ne peut pas calculer la moyenne avant d'avoir lu les deux notes !

\begin{popfigure}
\begin{center}
\begin{tikzpicture}[
  box/.style={draw=popblue, thick, rounded corners, fill=popblueL, minimum width=3.2cm, minimum height=1.6cm, align=center, font=\small},
  arr/.style={-{Stealth[length=3mm]}, thick, popdark}]
  \node[box, align=center] (in) at (0,0){\textbf{Entrées}\\ note1 = 12\\ note2 = 15};
  \node[box, fill=poporangeL, draw=poporange, align=center] (tr) at (5.2,0){\textbf{Traitement}\\ moyenne $\leftarrow$ (12 + 15) / 2};
  \node[box, fill=popgreenL, draw=popgreen, align=center] (out) at (10.4,0){\textbf{Sortie}\\ Afficher :\\ « La moyenne est : 13,5 »};
  \draw[arr] (in) -- (tr);
  \draw[arr] (tr) -- (out);
\end{tikzpicture}
\end{center}
\legende{Schéma « Entrées $\rightarrow$ Traitement $\rightarrow$ Sorties » pour l'algorithme de la moyenne.}
\end{popfigure}

\courssub{La trace d'exécution : vérifier son algorithme pas à pas}

\textbf{Intuition.} Comment être sûr qu'un algorithme est correct \emph{avant} de le faire exécuter par un ordinateur ? On joue soi-même le rôle de la machine : on exécute les instructions une par une, à la main, et on note dans un tableau la valeur de chaque variable après chaque étape. Ce tableau s'appelle la \cle{trace d'exécution}. C'est un outil formidable pour détecter les erreurs.

\begin{exoresolu}[Trace complète de l'algorithme de la moyenne]
Énoncé. On exécute l'algorithme \texttt{MoyenneDeuxNotes} avec les valeurs saisies \texttt{note1 = 12} et \texttt{note2 = 15}. Donner la trace d'exécution complète et le résultat affiché.
\tcblower
Solution détaillée. On suit les instructions dans l'ordre et on note le contenu de chaque variable après chaque étape (le tiret « -- » signifie que la variable n'a pas encore de valeur) :

\begin{center}
\begin{tabularx}{\linewidth}{|c|X|c|c|c|}
\hline
\rowcolor{popblueL}
\textbf{Étape} & \textbf{Instruction exécutée} & \textbf{note1} & \textbf{note2} & \textbf{moyenne} \\
\hline
1 & \texttt{Lire note1} (l'utilisateur tape 12) & 12 & -- & -- \\
\hline
2 & \texttt{Lire note2} (l'utilisateur tape 15) & 12 & 15 & -- \\
\hline
3 & \texttt{moyenne $\leftarrow$ (12 + 15) / 2} & 12 & 15 & 13,5 \\
\hline
4 & \texttt{Écrire "La moyenne est : ", moyenne} & 12 & 15 & 13,5 \\
\hline
\end{tabularx}
\end{center}

Détail du calcul de l'étape 3 :
\[ \text{moyenne} = \frac{\text{note1} + \text{note2}}{2} = \frac{12 + 15}{2} = \frac{27}{2} = 13{,}5 \]

\textbf{Résultat affiché à l'écran :} \texttt{La moyenne est : 13,5}

L'algorithme se termine (finitude respectée), chaque instruction était claire, et le résultat est correct : notre algorithme est validé.
\end{exoresolu}

\begin{attention}[Les pièges classiques du débutant]
\begin{itemize}
  \item \textbf{Utiliser une variable non déclarée} : toute variable doit être annoncée dans la partie \texttt{Variables} avant d'être utilisée.
  \item \textbf{Calculer avant de lire} : écrire \texttt{moyenne $\leftarrow$ (note1 + note2) / 2} \emph{avant} les instructions \texttt{Lire} donne un résultat faux, car les boîtes sont encore vides.
  \item \textbf{Oublier les parenthèses} : \texttt{(note1 + note2) / 2} et \texttt{note1 + note2 / 2} ne donnent pas le même résultat ! Sans parenthèses, la division est faite en priorité : $12 + 15/2 = 19{,}5$ au lieu de $13{,}5$.
  \item \textbf{Oublier d'afficher le résultat} : un calcul qui n'est jamais affiché ne sert à rien pour l'utilisateur.
\end{itemize}
\end{attention}

\begin{exercice}[À toi de jouer !]
\begin{enumerate}
  \item Maman Aïcha vend chaque beignet à 100 FCFA. Écris en pseudo-code un algorithme \texttt{Recette} qui demande le nombre de beignets vendus dans la journée, calcule la recette et l'affiche.
  \item Donne la trace d'exécution de ton algorithme pour 230 beignets vendus.
  \item Que se passe-t-il si l'on écrit \texttt{recette $\leftarrow$ nombre * 100} avant d'avoir lu la variable \texttt{nombre} ? Explique.
\end{enumerate}
\end{exercice}

\begin{corrige}[Exercice « À toi de jouer ! »]
\textbf{1.} Algorithme \texttt{Recette} :
\begin{lstlisting}
Algorithme Recette
Variables
    nombre : Entier      // nombre de beignets vendus
    recette : Entier     // recette de la journee en FCFA
Debut
    Ecrire "Entrez le nombre de beignets vendus : "
    Lire nombre
    recette <- nombre * 100        // chaque beignet coute 100 FCFA
    Ecrire "La recette du jour est : ", recette, " FCFA"
Fin
\end{lstlisting}

\textbf{2.} Trace d'exécution pour \texttt{nombre = 230} :
\begin{center}
\begin{tabularx}{\linewidth}{|c|X|c|c|}
\hline
\rowcolor{popblueL}
\textbf{Étape} & \textbf{Instruction exécutée} & \textbf{nombre} & \textbf{recette} \\
\hline
1 & \texttt{Lire nombre} (l'utilisateur tape 230) & 230 & -- \\
\hline
2 & \texttt{recette $\leftarrow$ 230 * 100} & 230 & 23\,000 \\
\hline
3 & \texttt{Écrire "La recette du jour est : ", recette, " FCFA"} & 230 & 23\,000 \\
\hline
\end{tabularx}
\end{center}
Résultat affiché : \texttt{La recette du jour est : 23000 FCFA}.

\textbf{3.} Si l'on écrit \texttt{recette $\leftarrow$ nombre * 100} avant le \texttt{Lire nombre}, la variable \texttt{nombre} ne contient encore aucune valeur : le calcul est impossible (ou donne un résultat imprévisible). C'est une violation de l'ordre logique des instructions : on doit toujours \emph{lire les entrées avant de les utiliser dans un traitement}.
\end{corrige}

\begin{aretenir}[L'essentiel de la section]
\begin{itemize}
  \item Un \cle{algorithme} est une suite \textbf{finie} et \textbf{ordonnée} d'instructions qui résout un problème : il part d'\textbf{entrées}, effectue des \textbf{traitements} et produit des \textbf{sorties}.
  \item Une \cle{variable} est une boîte nommée qui contient une valeur modifiable.
  \item En \cle{pseudo-code}, un algorithme s'écrit avec la structure \texttt{Algorithme / Variables / Début / Fin}.
  \item Les trois instructions de base sont : \texttt{Lire} (saisie), \texttt{Écrire} (affichage) et l'\textbf{affectation} $\leftarrow$ (ranger une valeur calculée dans une variable).
  \item L'affectation \texttt{x $\leftarrow$ 5} n'est pas une égalité mathématique : c'est un ordre, et le sens de la flèche compte.
  \item La \cle{trace d'exécution} est un tableau qui suit les valeurs des variables ligne par ligne : c'est l'outil pour vérifier qu'un algorithme est juste.
\end{itemize}
\end{aretenir}

Dans la prochaine section, nous verrons comment un algorithme peut \emph{prendre des décisions} grâce aux structures alternatives (« si... alors... sinon »), comme Maman Aïcha qui offre un beignet seulement si l'achat dépasse 500 FCFA.

\coursec{Les structures alternatives}

Dans la section précédente, nous avons découvert ce qu'est un \cle{algorithme} : une suite finie et ordonnée d'instructions, écrite en \cle{pseudo-code}, qui permet de résoudre un problème. Jusqu'ici, nos algorithmes exécutaient leurs instructions \emph{les unes après les autres}, toujours dans le même ordre : on parle d'\emph{exécution séquentielle}. Mais la vie n'est pas une ligne droite ! Le proviseur n'affiche pas le même message pour tous les élèves : il affiche « Admis » pour l'un, « Refusé » pour l'autre. La caissière de la boutique n'applique pas toujours la même remise. Bref : un bon algorithme doit savoir \emph{prendre des décisions}. C'est exactement le rôle des \cle{structures alternatives}, que nous allons étudier dans cette section.

\courssub{La notion de condition : une question qui n'a que deux réponses}

Avant de décider, il faut pouvoir \emph{tester}. Imagine la scène suivante : à la fin du BEPC blanc, le professeur principal regarde la moyenne de chaque élève. Pour chacun, il se pose une seule question : « la moyenne est-elle supérieure ou égale à 10 ? ». À cette question, il n'existe que deux réponses possibles : \emph{oui} ou \emph{non}. Une telle question s'appelle une \cle{condition}.

\begin{definition}[Condition]
Une \cle{condition} est une expression qui ne peut prendre que deux valeurs : \textbf{VRAI} ou \textbf{FAUX}. On dit aussi que c'est une \emph{expression booléenne}.
\end{definition}

\begin{exemplebox}[Des conditions de la vie courante]
\begin{itemize}
\item « note $\geq$ 10 » : VRAI si l'élève a 12, FAUX s'il a 8 ;
\item « solde $\geq$ montant du retrait » : VRAI si le compte contient assez d'argent ;
\item « âge $<$ 18 » : VRAI pour un élève de 3ème de 14 ans ;
\item « mot de passe $=$ "1234" » : VRAI seulement si la saisie est exacte.
\end{itemize}
\end{exemplebox}

Pour écrire des conditions, on utilise les \cle{opérateurs de comparaison}. Ils comparent deux valeurs (deux nombres, un nombre et une variable, etc.) et le résultat est VRAI ou FAUX.

\begin{center}
\begin{tabularx}{\linewidth}{|l|X|X|}
\hline
\rowcolor{popblueL} \textbf{Opérateur} & \textbf{Signification} & \textbf{Exemple (avec note $= 12$)} \\
\hline
$=$ & est égal à & note $= 12$ est VRAI \\
\hline
$\neq$ & est différent de & note $\neq 10$ est VRAI \\
\hline
$<$ & est strictement inférieur à & note $< 10$ est FAUX \\
\hline
$>$ & est strictement supérieur à & note $> 10$ est VRAI \\
\hline
$\leq$ & est inférieur ou égal à & note $\leq 12$ est VRAI \\
\hline
$\geq$ & est supérieur ou égal à & note $\geq 10$ est VRAI \\
\hline
\end{tabularx}
\end{center}

Il arrive qu'une seule comparaison ne suffise pas. Par exemple, la mention « Bien » au BEPC exige une moyenne \emph{à la fois} supérieure ou égale à 14 \emph{et} strictement inférieure à 16. On combine alors plusieurs conditions grâce aux \cle{connecteurs logiques} :

\begin{propriete}[Connecteurs logiques ET et OU]
\begin{itemize}
\item \textbf{condition1 ET condition2} est VRAI seulement si \emph{les deux} conditions sont vraies en même temps. Exemple : (moyenne $\geq$ 14) ET (moyenne $<$ 16).
\item \textbf{condition1 OU condition2} est VRAI dès qu'\emph{au moins une} des deux conditions est vraie. Exemple : (jour $=$ "samedi") OU (jour $=$ "dimanche").
\end{itemize}
\end{propriete}

\courssub{La structure alternative simple : Si \ldots Alors \ldots FinSi}

Revenons à notre boutique de quartier à Yaoundé. Le commerçant a décidé : « si un client achète pour plus de \num{10000} FCFA, je lui offre un bon de réduction ». Remarque bien : si le client achète pour \num{5000} FCFA, le commerçant ne fait \emph{rien de spécial}. Il n'y a qu'\emph{un seul cas} où une action est déclenchée. C'est la situation typique de l'\emph{alternative simple}.

\begin{definition}[Structure alternative simple]
La \cle{structure alternative simple} exécute un bloc d'instructions \emph{uniquement si} la condition est vraie. Si la condition est fausse, on ne fait rien et on passe directement à la suite de l'algorithme.
\end{definition}

\begin{propriete}[Syntaxe de l'alternative simple]
\begin{lstlisting}
Si condition Alors
    instructions
FinSi
\end{lstlisting}
La condition est évaluée : si elle est VRAIE, les \texttt{instructions} sont exécutées ; si elle est FAUSSE, elles sont ignorées. Dans les deux cas, l'algorithme continue après le \texttt{FinSi}.
\end{propriete}

\begin{exemplebox}[Féliciter les bons élèves]
\begin{lstlisting}
Algorithme Felicitations
Variable note : Reel
Debut
    Ecrire("Entrez la note : ")
    Lire(note)
    Si note >= 16 Alors
        Ecrire("Felicitations, excellent travail !")
    FinSi
    Ecrire("Fin du programme.")
Fin
\end{lstlisting}
Si l'élève saisit 17, le message de félicitations s'affiche. S'il saisit 12, le programme affiche directement « Fin du programme. » sans autre message.
\end{exemplebox}

\courssub{La structure alternative complète : Si \ldots Alors \ldots Sinon \ldots FinSi}

Très souvent, on ne veut pas « faire ou ne rien faire », mais choisir \emph{entre deux actions}. Le proviseur ne se contente pas d'afficher « Admis » : il doit afficher « Refusé » pour les autres. Il y a \emph{deux cas} et ils s'excluent mutuellement. C'est le rôle de l'\emph{alternative complète}.

\begin{definition}[Structure alternative complète]
La \cle{structure alternative complète} exécute un premier bloc d'instructions si la condition est vraie, et un \emph{autre} bloc si elle est fausse. L'un des deux blocs est \emph{toujours} exécuté, mais jamais les deux.
\end{definition}

\begin{propriete}[Syntaxe de l'alternative complète]
\begin{lstlisting}
Si condition Alors
    instructions1
Sinon
    instructions2
FinSi
\end{lstlisting}
Si la condition est VRAIE, seul \texttt{instructions1} est exécuté ; si elle est FAUSSE, seul \texttt{instructions2} est exécuté. Ensuite, l'algorithme continue après le \texttt{FinSi}.
\end{propriete}

\begin{exemplebox}[Admis ou Refusé]
\begin{lstlisting}
Algorithme Resultat
Variable moyenne : Reel
Debut
    Ecrire("Entrez la moyenne : ")
    Lire(moyenne)
    Si moyenne >= 10 Alors
        Ecrire("Admis")
    Sinon
        Ecrire("Refuse")
    FinSi
Fin
\end{lstlisting}
Avec une moyenne de 11,5 le programme affiche « Admis » ; avec 9,5 il affiche « Refusé ».
\end{exemplebox}

La figure suivante représente cette décision sous la forme d'un \emph{arbre de décision} : on part de la condition, puis on suit la branche « VRAI » ou la branche « FAUX ».

\begin{popfigure}
\begin{center}
\begin{tikzpicture}[
  cond/.style={diamond, draw=popblue, fill=popblueL, thick, aspect=2.2, align=center, inner sep=1pt},
  act/.style={rectangle, rounded corners, draw=popgreen, fill=popgreenL, thick, align=center, minimum width=2.6cm, minimum height=0.9cm},
  act2/.style={rectangle, rounded corners, draw=poporange, fill=poporangeL, thick, align=center, minimum width=2.6cm, minimum height=0.9cm},
  fleche/.style={-{Stealth[length=3mm]}, thick, popdark}]
\node[cond] (c) at (0,0) {moyenne $\geq$ 10 ?};
\node[act, align=center] (a) at (-3.4,-2.2){Afficher\\ « Admis »};
\node[act2, align=center] (r) at (3.4,-2.2){Afficher\\ « Refusé »};
\draw[fleche] (c) -- (a) node[midway, above left, popgreen] {\textbf{VRAI}};
\draw[fleche] (c) -- (r) node[midway, above right, poporange] {\textbf{FAUX}};
\end{tikzpicture}
\end{center}
\legende{Arbre de décision de l'alternative complète : la condition « moyenne $\geq$ 10 » oriente l'exécution vers l'une des deux branches, jamais les deux.}
\end{popfigure}

\begin{methode}[Choisir entre alternative simple et alternative complète]
\begin{enumerate}
\item Lis bien l'énoncé et repère la \emph{décision} à prendre (la condition).
\item Pose-toi la question : « si la condition est fausse, doit-on faire quelque chose de précis ? »
\item Si \textbf{non} (on ne fait rien de spécial), utilise l'alternative \emph{simple} : \texttt{Si \ldots Alors \ldots FinSi}.
\item Si \textbf{oui} (il y a une action pour le cas contraire), utilise l'alternative \emph{complète} : \texttt{Si \ldots Alors \ldots Sinon \ldots FinSi}.
\item Vérifie enfin que chaque \texttt{Si} possède bien son \texttt{FinSi}.
\end{enumerate}
\end{methode}

\begin{exoresolu}[La remise de 5 \%]
Énoncé. Une boutique de Douala accorde une remise de 5 \% sur tout achat dont le montant dépasse \num{10000} FCFA. Écris un algorithme qui lit le montant de l'achat et affiche le prix à payer. Fais ensuite la \emph{trace} (le déroulement pas à pas) pour un achat de \num{8000} FCFA puis pour un achat de \num{15000} FCFA.
\tcblower
Solution. Si le montant dépasse \num{10000} FCFA, on calcule la remise (5 \% du montant, c'est-à-dire montant $\times 0{,}05$) et on la retire. Sinon, le prix reste inchangé : c'est une alternative complète.
\begin{lstlisting}
Algorithme PrixAPayer
Variable montant, prix : Reel
Debut
    Ecrire("Entrez le montant de l'achat : ")
    Lire(montant)
    Si montant > 10000 Alors
        prix <- montant - montant * 0.05
    Sinon
        prix <- montant
    FinSi
    Ecrire("Prix a payer : ", prix, " FCFA")
Fin
\end{lstlisting}
\textbf{Trace pour montant $= \num{8000}$} (tableau de trace) :
\begin{center}
\begin{tabular}{|c|c|}
\hline
\textbf{Instruction} & \textbf{Valeur / Action} \\
\hline
Lire(montant) & 8000 \\
\hline
Si montant > 10000 & FAUX (8000 > 10000 est faux) \\
\hline
Exécuter Sinon & prix $\leftarrow$ 8000 \\
\hline
Ecrire(prix) & Affiche « Prix a payer : 8000 FCFA » \\
\hline
\end{tabular}
\end{center}

\textbf{Trace pour montant $= \num{15000}$} :
\begin{center}
\begin{tabular}{|c|c|}
\hline
\textbf{Instruction} & \textbf{Valeur / Action} \\
\hline
Lire(montant) & 15000 \\
\hline
Si montant > 10000 & VRAI (15000 > 10000 est vrai) \\
\hline
Exécuter Alors & prix $\leftarrow$ 15000 - 15000 $\times$ 0,05 = 14250 \\
\hline
Ecrire(prix) & Affiche « Prix a payer : 14250 FCFA » \\
\hline
\end{tabular}
\end{center}
\end{exoresolu}

\courssub{Les alternatives imbriquées : plusieurs cas possibles}

Que faire quand il y a \emph{plus de deux} cas ? Au BEPC, on distingue par exemple quatre situations : Ajourné, mention Passable, mention Bien, mention Très Bien. Une seule alternative ne suffit plus. La solution consiste à placer une alternative \emph{dans} le \texttt{Sinon} d'une autre : on parle d'\cle{alternatives imbriquées}.

\begin{propriete}[Syntaxe des alternatives imbriquées]
\begin{lstlisting}
Si condition1 Alors
    instructions1
Sinon
    Si condition2 Alors
        instructions2
    Sinon
        instructions3
    FinSi
FinSi
\end{lstlisting}
On teste les conditions \emph{dans l'ordre} : dès qu'une condition est vraie, son bloc est exécuté et tous les autres sont ignorés.
\end{propriete}

\begin{exemplebox}[Mentions au BEPC]
\begin{lstlisting}
Algorithme Mention
Variable moyenne : Reel
Debut
    Ecrire("Entrez la moyenne : ")
    Lire(moyenne)
    Si moyenne < 10 Alors
        Ecrire("Ajourne")
    Sinon
        Si moyenne < 14 Alors
            Ecrire("Mention Passable")
        Sinon
            Si moyenne < 16 Alors
                Ecrire("Mention Bien")
            Sinon
                Ecrire("Mention Tres Bien")
            FinSi
        FinSi
    FinSi
Fin
\end{lstlisting}
Remarque l'astuce : grâce à l'ordre des tests, la condition « moyenne $< 14$ » n'est examinée que si l'on sait déjà que moyenne $\geq 10$. Inutile donc d'écrire « moyenne $\geq$ 10 ET moyenne $<$ 14 », même si ce serait correct.
\end{exemplebox}

\begin{popfigure}
\begin{center}
\begin{tikzpicture}[
  cond/.style={diamond, draw=popblue, fill=popblueL, thick, aspect=2.4, align=center, inner sep=1pt},
  res/.style={rectangle, rounded corners, draw=poppurple, fill=poppurpleL, thick, align=center, minimum width=2.4cm, minimum height=0.8cm},
  fleche/.style={-{Stealth[length=3mm]}, thick, popdark}]
\node[cond] (c1) at (0,0) {moyenne $<$ 10 ?};
\node[res] (a1) at (-4.2,0) {Ajourné};
\node[cond] (c2) at (1.6,-2.1) {moyenne $<$ 14 ?};
\node[res] (a2) at (-2.8,-2.1) {Passable};
\node[cond] (c3) at (3.2,-4.2) {moyenne $<$ 16 ?};
\node[res] (a3) at (-1.2,-4.2) {Bien};
\node[res] (a4) at (6.4,-4.2) {Très Bien};
\draw[fleche] (c1) -- (a1) node[midway, above, poppurple] {\textbf{VRAI}};
\draw[fleche] (c1) -- (c2) node[midway, above right, popdark] {\textbf{FAUX}};
\draw[fleche] (c2) -- (a2) node[midway, above, poppurple] {\textbf{VRAI}};
\draw[fleche] (c2) -- (c3) node[midway, above right, popdark] {\textbf{FAUX}};
\draw[fleche] (c3) -- (a3) node[midway, above, poppurple] {\textbf{VRAI}};
\draw[fleche] (c3) -- (a4) node[midway, above, poppurple] {\textbf{FAUX}};
\end{tikzpicture}
\end{center}
\legende{Arbre de décision à trois niveaux pour les mentions : chaque losange est un test, et la branche « FAUX » mène au test suivant.}
\end{popfigure}

\courssub{Représentation par algorigramme : symboles normalisés}

Le programme officiel demande de savoir \emph{représenter un algorithme par un algorigramme}. Contrairement à l'arbre de décision qui ne montre que la logique des tests, l'algorigramme utilise des symboles normalisés (norme ISO 5807) pour décrire \emph{tout} l'algorithme, du début à la fin.

\begin{definition}[Symboles de base de l'algorigramme]
\begin{itemize}
\item \textbf{Ovale} : Début et Fin de l'algorithme.
\item \textbf{Parallélogramme} : Entrée (Saisir/Lire) et Sortie (Afficher/Écrire).
\item \textbf{Rectangle} : Traitement / Affectation (ex: \texttt{prix $\leftarrow$ montant * 0,95}).
\item \textbf{Losange} : Test / Condition (Structure alternative). Il a \emph{deux} sorties : « Oui » (VRAI) et « Non » (FAUX).
\item \textbf{Flèches} : Sens d'exécution (généralement de haut en bas).
\end{itemize}
\end{definition}

\begin{methode}[Dessiner l'algorigramme d'une alternative complète]
\begin{enumerate}
\item \textbf{Début} (Ovale).
\item \textbf{Saisir} les données d'entrée (Parallélogramme).
\item \textbf{Tester} la condition (Losange). Écrire la condition à l'intérieur.
\item Dessiner la branche \textbf{Oui (VRAI)} : y placer les instructions du \texttt{Alors} (Rectangles ou Parallélogrammes).
\item Dessiner la branche \textbf{Non (FAUX)} : y placer les instructions du \texttt{Sinon}.
\item \textbf{Rejoindre} les deux branches après le \texttt{FinSi}.
\item \textbf{Fin} (Ovale).
\end{enumerate}
\end{methode}

\begin{exemplebox}[Algorigramme complet : Algorithme « Resultat » (Admis/Refusé)]
\begin{popfigure}
\begin{center}
\begin{tikzpicture}[
  startstop/.style={ellipse, draw=popdark, fill=popcream, thick, minimum width=2cm, minimum height=1cm},
  io/.style={trapezium, draw=popblue, fill=popblueL, thick, minimum width=2.5cm, minimum height=0.9cm, trapezium left angle=70, trapezium right angle=110},
  proc/.style={rectangle, draw=popgreen, fill=popgreenL, thick, minimum width=2.5cm, minimum height=0.9cm},
  dec/.style={diamond, draw=poporange, fill=poporangeL, thick, aspect=2, align=center, inner sep=1pt},
  arrow/.style={-{Stealth[length=3mm]}, thick, popdark}]
\node[startstop] (debut) at (0,0) {Début};
\node[io] (saisie) at (0,-1.8) {Saisir moyenne};
\node[dec] (test) at (0,-3.6) {moyenne $\geq$ 10 ?};
\node[proc] (admis) at (-3.5,-5.4) {Afficher « Admis »};
\node[proc] (refuse) at (3.5,-5.4) {Afficher « Refusé »};
\node[coordinate] (joint) at (0,-7.2) {};
\node[startstop] (fin) at (0,-9) {Fin};
\draw[arrow] (debut) -- (saisie);
\draw[arrow] (saisie) -- (test);
\draw[arrow] (test) -- node[left, pos=0.5] {\textbf{Oui}} (admis);
\draw[arrow] (test) -- node[right, pos=0.5] {\textbf{Non}} (refuse);
\draw[arrow] (admis) -- (joint);
\draw[arrow] (refuse) -- (joint);
\draw[arrow] (joint) -- (fin);
\end{tikzpicture}
\end{center}
\legende{Algorigramme normalisé de l'algorithme « Resultat ». Notez la jonction des deux branches après le test.}
\end{popfigure}
\end{exemplebox}

\begin{attention}[Deux erreurs très fréquentes au BEPC]
\begin{itemize}
\item \textbf{Oublier le \texttt{FinSi}} : chaque \texttt{Si} doit être fermé par un \texttt{FinSi}. Avec trois alternatives imbriquées, il faut \emph{trois} \texttt{FinSi}. Compte-les toujours à la fin de l'écriture !
\item \textbf{Se tromper d'opérateur (piège des bornes)} : si l'énoncé dit « admis à partir de 10 », un élève qui a exactement 10 \emph{doit} être admis. Il faut donc écrire \texttt{Si moyenne >= 10} et non \texttt{Si moyenne > 10}, qui exclurait la note 10. De même, « dépasse \num{10000} » signifie \texttt{> 10000}, tandis que « à partir de \num{10000} » signifie \texttt{>= 10000}. Lis l'énoncé avec la plus grande attention.
\item \textbf{Mauvaise indentation} : dans les alternatives imbriquées, décale le code vers la droite à chaque niveau de \texttt{Si}. Cela rend la lecture visuelle des \texttt{FinSi} correspondants beaucoup plus facile.
\end{itemize}
\end{attention}

\begin{savaistu}[Les alternatives dans ton téléphone]
Les applications de transfert d'argent comme Orange Money ou MTN Mobile Money, très utilisées au Cameroun, exécutent des alternatives à chaque transaction. Quand tu envoies de l'argent, le système teste d'abord : « solde $\geq$ montant + frais ? ». Si c'est VRAI, le transfert est effectué ; \emph{Sinon}, tu reçois le message « Solde insuffisant ». Des millions de décisions de ce type sont ainsi prises automatiquement chaque jour !
\end{savaistu}

\begin{exercice}[À toi de jouer]
\begin{enumerate}
\item Écris un algorithme qui lit l'âge d'une personne et affiche « Mineur » si l'âge est inférieur à 18 ans, et « Majeur » sinon. Dessine l'algorigramme correspondant.
\item Un cybercafé facture l'heure de connexion à 300 FCFA, mais offre une heure gratuite à tout client qui achète plus de 5 heures. Écris l'algorithme qui lit le nombre d'heures achetées et affiche le montant à payer.
\item Reprends l'algorithme des mentions (section « Alternatives imbriquées ») et fais sa trace (sous forme de tableau) pour les moyennes 8, 12, 15 et 17.
\item \textbf{Défi :} Écris un algorithme qui lit trois nombres \texttt{a}, \texttt{b}, \texttt{c} et affiche le plus grand des trois. (Indication : utilise des alternatives imbriquées ou une variable \texttt{max}).
\end{enumerate}
\end{exercice}

\begin{corrige}[Exercice (éléments de correction)]
\begin{enumerate}
\item \textbf{Algorithme :} \texttt{Si age < 18 Alors Ecrire("Mineur") Sinon Ecrire("Majeur") FinSi}. \textbf{Algorigramme :} Début $\rightarrow$ Saisir age $\rightarrow$ Losange (age < 18 ?) $\rightarrow$ Oui : Afficher « Mineur » / Non : Afficher « Majeur » $\rightarrow$ Jonction $\rightarrow$ Fin.
\item \texttt{Algorithme Cyber\\ Variable heures, payer : Entier\\ Debut\\ \quad Ecrire("Heures achetees : ")\\ \quad Lire(heures)\\ \quad Si heures > 5 Alors\\ \quad \quad payer <- (heures - 1) * 300\\ \quad Sinon\\ \quad \quad payer <- heures * 300\\ \quad FinSi\\ \quad Ecrire("Montant : ", payer)\\ Fin}
\item \textbf{Traces :}
\begin{itemize}
\item 8 : Test1 (8<10) VRAI $\rightarrow$ « Ajourné ».
\item 12 : Test1 FAUX $\rightarrow$ Test2 (12<14) VRAI $\rightarrow$ « Mention Passable ».
\item 15 : Test1 FAUX $\rightarrow$ Test2 FAUX $\rightarrow$ Test3 (15<16) VRAI $\rightarrow$ « Mention Bien ».
\item 17 : Test1 FAUX $\rightarrow$ Test2 FAUX $\rightarrow$ Test3 FAUX $\rightarrow$ « Mention Tres Bien ».
\end{itemize}
\item \textbf{Algorithme Max3 :} \texttt{Si a >= b Alors\\ \quad Si a >= c Alors max <- a Sinon max <- c FinSi\\ Sinon\\ \quad Si b >= c Alors max <- b Sinon max <- c FinSi\\ FinSi\\ Ecrire(max)}.
\end{enumerate}
\end{corrige}

\begin{aretenir}[L'essentiel de la section]
\begin{itemize}
\item Une \cle{condition} est une expression qui vaut VRAI ou FAUX ; elle s'écrit avec les opérateurs $=$, $\neq$, $<$, $>$, $\leq$, $\geq$ et peut combiner plusieurs tests avec ET et OU.
\item L'\cle{alternative simple} (\texttt{Si \ldots Alors \ldots FinSi}) n'agit que si la condition est vraie.
\item L'\cle{alternative complète} (\texttt{Si \ldots Alors \ldots Sinon \ldots FinSi}) choisit entre deux actions : l'une des deux est toujours exécutée.
\item Pour plus de deux cas, on \emph{imbrique} les alternatives : chaque \texttt{Si} doit avoir son \texttt{FinSi}.
\item L'\cle{algorigramme} utilise des symboles normalisés : Ovale (Début/Fin), Parallélogramme (Entrée/Sortie), Rectangle (Traitement), Losange (Test).
\item Attention aux bornes : « à partir de 10 » se traduit par $\geq 10$, pas par $> 10$.
\end{itemize}
\end{aretenir}

\coursec{Les structures itératives : la boucle « tant que »}

Dans la section précédente, nous avons découvert les structures alternatives : grâce à \texttt{Si ... Alors ... Sinon ... FinSi}, un algorithme peut \emph{choisir} entre plusieurs chemins selon une condition. Mais un choix, même bien fait, ne s'exécute qu'\textbf{une seule fois}. Or, dans la vie courante, beaucoup de tâches consistent à \emph{répéter} la même action : compter les élèves d'une classe un par un, ajouter de l'argent à une épargne chaque semaine, afficher les dix premiers nombres entiers... Écrire dix fois la même instruction serait fastidieux, et mille fois, impossible. C'est précisément le problème que résolvent les \cle{structures itératives}, aussi appelées \cle{boucles}. Dans cette section, nous étudions la plus intuitive d'entre elles : la boucle \texttt{TantQue}.

\courssub{Notion de boucle : répéter tant qu'une condition est vraie}

\begin{experience}[Une situation pour comprendre : l'épargne d'Amina]
Amina, élève de 3ème, veut acheter un livre de préparation au BEPC qui coûte \SI{50000}{FCFA}. Chaque semaine, elle dépose \SI{5000}{FCFA} dans sa cagnotte. Observe sa démarche :
\begin{enumerate}
  \item Elle regarde le contenu de sa cagnotte : a-t-elle atteint \SI{50000}{FCFA} ?
  \item Si \textbf{non}, elle ajoute \SI{5000}{FCFA} et recommence à l'étape 1.
  \item Si \textbf{oui}, elle s'arrête : elle peut acheter son livre.
\end{enumerate}
\textbf{Observation :} Amina répète la même action (ajouter \SI{5000}{FCFA}) \emph{tant que} le total est inférieur à \SI{50000}{FCFA}. Elle ne sait pas à l'avance combien de fois elle répétera l'action en la comptant mentalement : elle s'arrête quand la \emph{condition} devient fausse. C'est exactement le fonctionnement de la boucle \texttt{TantQue}.
\end{experience}

\begin{definition}[Structure itérative et boucle TantQue]
Une \cle{structure itérative} (ou \cle{boucle}) est une structure de contrôle qui permet d'\textbf{exécuter plusieurs fois} un même bloc d'instructions.

La boucle \texttt{TantQue} exécute un bloc d'instructions \textbf{tant qu'une condition est vraie} :
\begin{itemize}
  \item on \textbf{teste} d'abord la condition ;
  \item si la condition est \textbf{vraie}, on exécute le bloc d'instructions (appelé \emph{corps de la boucle}), puis on \textbf{revient tester} la condition ;
  \item si la condition est \textbf{fausse}, on \textbf{sort} de la boucle et on continue l'algorithme après le \texttt{FinTantQue}.
\end{itemize}
Chaque exécution du corps de la boucle s'appelle une \cle{itération}.
\end{definition}

\begin{propriete}[Syntaxe de la boucle TantQue]
En pseudo-code, la boucle \texttt{TantQue} s'écrit :
\begin{lstlisting}
TantQue condition Faire
    instructions
FinTantQue
\end{lstlisting}
\begin{itemize}
  \item \texttt{condition} est une expression logique (vraie ou fausse), par exemple \texttt{total < 50000} ;
  \item le test a lieu \textbf{avant} chaque itération, y compris la première : si la condition est fausse dès le départ, le corps de la boucle n'est \textbf{jamais} exécuté ;
  \item les instructions du corps sont indentées (décalées vers la droite) pour bien montrer qu'elles appartiennent à la boucle.
\end{itemize}
\end{propriete}

\begin{popfigure}
\begin{center}
\begin{tikzpicture}[
  node distance=0.7cm,
  bloc/.style={rectangle, rounded corners=3pt, draw=popblue, fill=popblueL, thick, minimum width=4.6cm, minimum height=0.9cm, align=center, font=\small},
  test/.style={diamond, draw=poporange, fill=poporangeL, thick, aspect=2.4, align=center, font=\small, inner sep=1pt},
  fleche/.style={-{Stealth[length=2.5mm]}, thick, popdark}]
  \node[bloc, align=center] (init){Initialisation\\ (donner une valeur de départ)};
  \node[test, below=of init] (cond) {Condition ?};
  \node[bloc, below=of cond, align=center] (corps){Corps de la boucle\\ (traitement + mise à jour)};
  \node[bloc, below=of corps, fill=popgreenL, draw=popgreen] (suite) {Suite de l'algorithme};
  \draw[fleche] (init) -- (cond);
  \draw[fleche] (cond) -- node[right, font=\small\itshape]{Vraie} (corps);
  \draw[fleche] (corps.west) -- ++(-1.6,0) |- node[pos=0.25, left, font=\small\itshape]{retour au test} (cond.west);
  \draw[fleche] (cond.east) -- ++(1.6,0) node[pos=0.5, above, font=\small\itshape]{Fausse} |- (suite.east);
\end{tikzpicture}
\end{center}
\legende{Algorigramme de la boucle \texttt{TantQue} : on teste la condition ; si elle est vraie, on exécute le corps puis on revient au test ; si elle est fausse, on sort de la boucle.}
\end{popfigure}

\courssub{Les trois ingrédients d'une boucle correcte}

Pour qu'une boucle \texttt{TantQue} fonctionne bien, trois éléments doivent toujours être présents. Leur oubli est la cause de presque toutes les erreurs.

\begin{methode}[Construire une boucle TantQue correcte]
\begin{enumerate}
  \item \textbf{Initialiser} : avant la boucle, donner une valeur de départ à la variable utilisée dans la condition (exemple : \texttt{total $\leftarrow$ 0}). Sans initialisation, la condition porte sur une valeur inconnue.
  \item \textbf{Tester} : écrire la condition qui doit être vraie pour \emph{continuer} (exemple : \texttt{TantQue total < 50000}).
  \item \textbf{Traiter et mettre à jour} : dans le corps de la boucle, effectuer le travail voulu \emph{et} modifier la variable de la condition (exemple : \texttt{total $\leftarrow$ total + 5000}), afin que la condition puisse devenir fausse.
\end{enumerate}
On retient le réflexe : \cle{initialiser}, \cle{tester}, \cle{traiter}, \cle{mettre à jour}.
\end{methode}

\begin{attention}[Erreur fréquente : la boucle infinie]
Si l'on \textbf{oublie la mise à jour} de la variable dans le corps de la boucle, la condition reste vraie pour toujours : la boucle ne s'arrête jamais. On dit qu'on a créé une \cle{boucle infinie}. L'ordinateur exécute alors les mêmes instructions sans fin, et le programme « plante ».
\begin{lstlisting}
total <- 0
TantQue total < 50000 Faire
    Afficher("Pas encore assez...")
    // ERREUR : total n'est jamais modifié !
FinTantQue
\end{lstlisting}
Ici, \texttt{total} vaut 0 pour l'éternité : la condition \texttt{total < 50000} reste vraie indéfiniment. Vérifie toujours que quelque chose \emph{change} dans la boucle.
\end{attention}

\begin{propriete}[Arrêt d'une boucle TantQue (admis)]
Une boucle \texttt{TantQue} s'arrête si et seulement si sa \textbf{condition finit par devenir fausse}. Pour cela, il faut que les instructions du corps modifient au moins une variable de la condition, dans le sens qui rapproche de l'arrêt. Avant d'utiliser une boucle, on doit toujours se poser la question : \emph{« qu'est-ce qui garantit que la condition deviendra fausse ? »}. Sur l'exemple de l'épargne, \texttt{total} augmente de \SI{5000}{FCFA} à chaque tour : il dépassera forcément \SI{50000}{FCFA}, donc la boucle s'arrête.
\end{propriete}

\courssub{Exemple 1 : compter de 1 à 10}

\begin{exemplebox}[Afficher les nombres de 1 à 10]
On veut afficher les nombres 1, 2, 3, ..., 10. Sans boucle, il faudrait dix instructions \texttt{Afficher}. Avec une boucle \texttt{TantQue}, trois lignes suffisent :
\begin{lstlisting}
Algorithme Compter
Variable
    i : Entier
Debut
    i <- 1                  // Initialisation : on commence à 1
    TantQue i <= 10 Faire   // Condition : on continue tant que i ne dépasse pas 10
        Afficher(i)         // Traitement : afficher le nombre courant
        i <- i + 1          // Mise à jour : passer au nombre suivant
    FinTantQue
Fin
\end{lstlisting}
\textbf{Déroulement :} \texttt{i} vaut 1, la condition \texttt{i <= 10} est vraie, on affiche 1 puis \texttt{i} devient 2 ; on reteste, on affiche 2, etc. Quand \texttt{i} vaut 11, la condition devient fausse : on sort de la boucle. La boucle a effectué exactement \textbf{10 itérations}.
\end{exemplebox}

\courssub{Exemple 2 : l'épargne hebdomadaire d'Amina}

Reprenons la situation de départ et traduisons-la en algorithme complet.

\begin{exemplebox}[Combien de semaines pour atteindre 50 000 FCFA ?]
\begin{lstlisting}
Algorithme Epargne
Variable
    total, semaine : Entier
Debut
    total <- 0                    // Initialisation : cagnotte vide
    semaine <- 0                  // Aucune semaine écoulée
    TantQue total < 50000 Faire   // On continue tant que l'objectif n'est pas atteint
        total <- total + 5000     // Mise à jour : dépôt hebdomadaire
        semaine <- semaine + 1    // Une semaine de plus
    FinTantQue
    Afficher("Objectif atteint en ", semaine, " semaines.")
Fin
\end{lstlisting}
\end{exemplebox}

Pour vérifier qu'un algorithme est correct, on réalise une \cle{trace d'exécution} : on joue soi-même le rôle de l'ordinateur et on note, dans un tableau, la valeur de chaque variable à la fin de chaque itération.

\begin{center}
\rowcolors{1}{popcream}{white}
\begin{tabularx}{\linewidth}{|c|c|c|X|}
\hline
\rowcolor{popblueL} \textbf{Semaine} & \textbf{Montant ajouté (FCFA)} & \textbf{Total (FCFA)} & \textbf{Condition \texttt{total < 50000}} \\
\hline
(départ) & --- & 0 & Vraie : on entre dans la boucle \\
\hline
1 & 5 000 & 5 000 & Vraie \\
\hline
2 & 5 000 & 10 000 & Vraie \\
\hline
3 & 5 000 & 15 000 & Vraie \\
\hline
4 & 5 000 & 20 000 & Vraie \\
\hline
5 & 5 000 & 25 000 & Vraie \\
\hline
6 & 5 000 & 30 000 & Vraie \\
\hline
7 & 5 000 & 35 000 & Vraie \\
\hline
8 & 5 000 & 40 000 & Vraie \\
\hline
9 & 5 000 & 45 000 & Vraie \\
\hline
10 & 5 000 & 50 000 & \textbf{Fausse : on sort de la boucle} \\
\hline
\end{tabularx}
\end{center}

Au bout de la 10\textsuperscript{e} semaine, \texttt{total} vaut \SI{50000}{FCFA} : la condition \texttt{total < 50000} devient fausse, la boucle s'arrête et l'algorithme affiche : « Objectif atteint en 10 semaines. » On retrouve bien le calcul $50\,000 \div 5\,000 = 10$.

\begin{exoresolu}[Lecture de SMS : compter les messages non lus]
\textbf{Énoncé.} Le téléphone de la maman de Rodrigue affiche ses SMS un par un. Écris un algorithme qui, partant de 25 messages non lus, en lit 5 par jour et affiche au bout de combien de jours il ne restera plus de message non lu. Donne la trace d'exécution.
\tcblower
\textbf{Solution.} On initialise \texttt{reste} à 25 et \texttt{jour} à 0. Tant qu'il reste des messages, on en lit 5 et on compte un jour.
\begin{lstlisting}
Algorithme LectureSMS
Variable
    reste, jour : Entier
Debut
    reste <- 25
    jour <- 0
    TantQue reste > 0 Faire
        reste <- reste - 5    // Mise à jour : 5 messages lus
        jour <- jour + 1
    FinTantQue
    Afficher("Tous les messages sont lus en ", jour, " jours.")
Fin
\end{lstlisting}
Trace d'exécution : \texttt{reste} prend les valeurs 20, 15, 10, 5 puis 0, et \texttt{jour} prend 1, 2, 3, 4, 5. Quand \texttt{reste} vaut 0, la condition \texttt{reste > 0} est fausse : on sort. L'algorithme affiche « Tous les messages sont lus en 5 jours. » La mise à jour (\texttt{reste} diminue de 5) garantit l'arrêt : pas de boucle infinie.
\end{exoresolu}

\begin{aretenir}[L'essentiel sur la boucle TantQue]
\begin{itemize}
  \item Une boucle permet de \textbf{répéter} des instructions ; la boucle \texttt{TantQue} répète \textbf{tant qu'une condition est vraie}, en testant la condition \emph{avant} chaque itération.
  \item Syntaxe : \texttt{TantQue condition Faire ... FinTantQue}.
  \item Une boucle correcte exige toujours : une \textbf{initialisation} avant la boucle, un \textbf{traitement} et une \textbf{mise à jour} dans la boucle.
  \item Sans mise à jour, la condition reste vraie : c'est la \textbf{boucle infinie}, l'erreur à éviter absolument.
  \item La \textbf{trace d'exécution} (tableau des valeurs des variables) est l'outil pour vérifier qu'une boucle fait bien ce qu'on attend.
\end{itemize}
\end{aretenir}

\begin{exercice}[À toi de jouer]
\begin{enumerate}
  \item Écris un algorithme qui affiche les nombres pairs de 2 à 20 à l'aide d'une boucle \texttt{TantQue}.
  \item Un commerçant de Mokolo dépose \SI{10000}{FCFA} par jour dans sa caisse. Écris un algorithme qui calcule en combien de jours il atteindra \SI{150000}{FCFA}, puis donne sa trace d'exécution.
  \item L'algorithme suivant contient une erreur. Trouve-la et corrige-la.
\begin{lstlisting}
i <- 1
TantQue i <= 5 Faire
    Afficher(i)
FinTantQue
\end{lstlisting}
\end{enumerate}
\end{exercice}

\begin{corrige}[Exercice : À toi de jouer]
\begin{enumerate}
  \item On initialise \texttt{i} à 2 et on avance de 2 en 2 :
\begin{lstlisting}
i <- 2
TantQue i <= 20 Faire
    Afficher(i)
    i <- i + 2
FinTantQue
\end{lstlisting}
  \item \texttt{caisse <- 0}, \texttt{jour <- 0} ; tant que \texttt{caisse < 150000}, on ajoute \SI{10000}{FCFA} et un jour. La trace donne 10 000, 20 000, ..., 150 000 : l'objectif est atteint en \textbf{15 jours} ($150\,000 \div 10\,000 = 15$).
  \item Il manque la \textbf{mise à jour} de \texttt{i} dans la boucle : la condition \texttt{i <= 5} reste vraie pour toujours, c'est une boucle infinie. Correction : ajouter \texttt{i <- i + 1} après \texttt{Afficher(i)}.
\end{enumerate}
\end{corrige}

\begin{savaistu}[Et dans un vrai programme ?]
La boucle \texttt{TantQue} existe dans tous les langages de programmation : elle s'écrit \texttt{while} en anglais (mot qui signifie « tant que »). C'est grâce à des boucles qu'un distributeur automatique de la SGC ou d'Afriland recompte ton solde après chaque opération, et que le logiciel de la cantine scolaire parcourt la liste de tous les élèves inscrits, un par un, sans jamais en oublier.
\end{savaistu}

\coursec{Les structures itératives : la boucle « pour »}

Dans la section précédente, nous avons découvert la boucle \cle{TantQue} : elle répète des instructions \emph{tant qu'une condition reste vraie}, et c'est le programmeur qui doit gérer lui-même l'initialisation et la progression de la variable de contrôle. Cette boucle est très souple, mais elle demande de la vigilance (risque de boucle infinie si l'on oublie de faire progresser la condition). Or, dans de très nombreuses situations de la vie courante, on sait \textbf{à l'avance combien de fois} on veut répéter une action : afficher les notes des 25 élèves d'une classe de 3ème, calculer la paie des 12 mois de l'année, afficher la table de multiplication de 7 (de $1 \times 7$ jusqu'à $10 \times 7$). Pour ces cas, l'algorithmique nous offre une boucle plus simple et plus sûre : la boucle \cle{Pour}.

\courssub{Quand utilise-t-on la boucle « Pour » ?}

Imaginons la cantine du collège : la cuisinière doit distribuer un repas à chacun des 100 élèves inscrits. Elle ne se demande pas « est-ce que je continue ? » à chaque assiette : elle sait qu'elle fera \emph{exactement} 100 distributions, ni plus, ni moins. Le nombre de répétitions est \textbf{connu à l'avance}.

\begin{definition}[La boucle Pour]
La boucle \cle{Pour} (dite aussi boucle « à compteur ») est une structure itérative qui répète un bloc d'instructions un \textbf{nombre de fois connu à l'avance}. Elle utilise une variable appelée \cle{compteur}, qui prend automatiquement, tour après tour, toutes les valeurs entières comprises entre une \textbf{valeur de départ} et une \textbf{valeur d'arrivée}.
\end{definition}

\begin{definition}[Le compteur]
Le \cle{compteur} est la variable qui pilote la boucle \cle{Pour}. Il est \textbf{initialisé automatiquement} à la valeur de départ, puis \textbf{augmenté automatiquement} d'une quantité fixe appelée le \cle{pas} (par défaut, le pas vaut $1$) à la fin de chaque tour. La boucle s'arrête lorsque le compteur dépasse la valeur d'arrivée.
\end{definition}

La grande différence avec le \texttt{TantQue} est donc là : avec le \texttt{Pour}, le programmeur n'a \textbf{ni à initialiser, ni à incrémenter} le compteur lui-même. La boucle s'en charge. C'est ce qui la rend plus sûre : on ne peut pas oublier l'incrémentation, donc pas de boucle infinie par oubli.

\courssub{Syntaxe de la boucle « Pour »}

\begin{propriete}[Syntaxe générale de la boucle Pour]
En pseudo-code, la boucle Pour s'écrit :
\begin{lstlisting}
Pour compteur de valeur_debut a valeur_fin Faire
    instruction 1
    instruction 2
    ...
FinPour
\end{lstlisting}
\begin{itemize}
  \item \texttt{compteur} : une variable de type entier ;
  \item \texttt{valeur\_debut} et \texttt{valeur\_fin} : deux expressions entières, avec valeur\_debut $\leq$ valeur\_fin ;
  \item à chaque tour, le compteur augmente automatiquement de $1$ (c'est le \cle{pas} par défaut) ;
  \item la boucle s'arrête dès que le compteur dépasse \texttt{valeur\_fin}.
\end{itemize}
\emph{Note :} on peut préciser un pas différent de $1$ avec la clause optionnelle \texttt{Pas valeur\_pas} (ex. \texttt{Pour i de 2 a 20 Pas 2 Faire}).
\end{propriete}

\textbf{Le pas différent de 1 (complément).} Dans certains langages et dans certaines écritures de pseudo-code, on peut préciser un \cle{pas} différent de $1$ : par exemple \texttt{Pour i de 2 a 20 Pas 2} fait prendre à \texttt{i} les valeurs $2, 4, 6, \dots, 20$ (les nombres pairs). On peut même utiliser un pas négatif pour compter à rebours : \texttt{Pour i de 10 a 1 Pas -1} produit $10, 9, 8, \dots, 1$. Au niveau de la classe de 3ème, on retient surtout le cas du pas égal à $1$, mais il est utile de savoir que cette possibilité existe.

\courssub{Premier exemple : la somme des entiers de 1 à 100}

\textbf{Problème.} On veut calculer $S = 1 + 2 + 3 + \dots + 100$. Le nombre d'additions à effectuer est parfaitement connu : 100. C'est donc le cas d'école de la boucle \texttt{Pour}.

\textbf{Idée.} On utilise une variable \texttt{somme}, initialisée à $0$, qui va « accumuler » les entiers un par un : au tour où le compteur \texttt{i} vaut $1$, on ajoute $1$ ; au tour où il vaut $2$, on ajoute $2$ ; et ainsi de suite jusqu'à $100$.

\begin{exemplebox}[Algorithme : somme des entiers de 1 à 100]
\begin{lstlisting}
Algorithme Somme1a100
Variables
    i, somme : Entier
Debut
    somme <- 0            // l'accumulateur part de zero
    Pour i de 1 a 100 Faire
        somme <- somme + i   // on ajoute la valeur du compteur
    FinPour
    Ecrire("La somme vaut : ", somme)
Fin
\end{lstlisting}
\end{exemplebox}

Déroulons mentalement les premiers tours : avant la boucle, \texttt{somme} vaut $0$. Au premier tour, \texttt{i} vaut $1$ et \texttt{somme} devient $0 + 1 = 1$. Au deuxième tour, \texttt{i} vaut $2$ et \texttt{somme} devient $1 + 2 = 3$. Au troisième tour, \texttt{somme} devient $3 + 3 = 6$. On voit que la variable \texttt{somme} contient à chaque instant la somme de tous les entiers déjà rencontrés : c'est pourquoi on l'appelle un \cle{accumulateur}.

\begin{exoresolu}[Trace de la somme $1 + 2 + \dots + 10$ et vérification]
Énoncé : pour bien comprendre le fonctionnement de l'accumulateur, on exécute « à la main » l'algorithme précédent en le limitant à 10 (au lieu de 100). Compléter le tableau de trace donnant, à chaque tour, la valeur de \texttt{i} et celle de \texttt{somme} après l'instruction \texttt{somme <- somme + i}. Quelle valeur finale obtient-on ? Vérifier avec la formule $\dfrac{n(n+1)}{2}$.
\tcblower
Solution détaillée : on part de \texttt{somme} $= 0$ et on ajoute successivement $1, 2, 3, \dots, 10$ :
\begin{center}
\begin{tabularx}{\linewidth}{|c|*{10}{>{\centering\arraybackslash}X|}}
\hline
\rowcolor{popblueL} Tour (valeur de i) & 1 & 2 & 3 & 4 & 5 & 6 & 7 & 8 & 9 & 10 \\
\hline
somme après le tour & 1 & 3 & 6 & 10 & 15 & 21 & 28 & 36 & 45 & 55 \\
\hline
\end{tabularx}
\end{center}
La valeur finale affichée est donc \texttt{somme} $= 55$. Vérification par la formule de la somme des $n$ premiers entiers :
\[ S = \frac{n(n+1)}{2} = \frac{10 \times 11}{2} = \frac{110}{2} = 55. \]
Les deux résultats concordent : l'algorithme est correct. Pour $n = 100$, la même formule donne $S = \dfrac{100 \times 101}{2} = \dfrac{10100}{2} = 5050$, ce qui confirme le résultat affiché par l'algorithme.
\end{exoresolu}

\begin{savaistu}[Le jeune Gauss et la somme de 1 à 100]
Carl Friedrich Gauss (1777--1855), l'un des plus grands mathématiciens de l'histoire, aurait ébloui son instituteur dès l'âge de 8 ans. Pour occuper la classe, le maître avait demandé de calculer $1 + 2 + 3 + \dots + 100$, pensant que cela prendrait longtemps. Le jeune Gauss aurait répondu presque instantanément : 5050 ! Son astuce ? Regrouper les nombres par paires : $1 + 100 = 101$, $2 + 99 = 101$, $3 + 98 = 101$, etc. Il y a 50 paires, chacune valant 101, donc $50 \times 101 = 5050$. C'est exactement la formule $\dfrac{n(n+1)}{2}$. Ton ordinateur, lui, fait les 100 additions une par une avec la boucle \texttt{Pour}... mais il les fait en une fraction de seconde !
\end{savaistu}

\courssub{Deuxième exemple : la table de multiplication de 7}

\textbf{Problème.} Un élève de 3ème veut afficher la table de multiplication de 7, de $1 \times 7$ jusqu'à $10 \times 7$. Le nombre de lignes à afficher est connu : 10. Encore un cas parfait pour la boucle \texttt{Pour}.

\begin{exemplebox}[Algorithme : table de multiplication de 7]
\begin{lstlisting}
Algorithme TableDe7
Variables
    i, produit : Entier
Debut
    Pour i de 1 a 10 Faire
        produit <- 7 * i
        Ecrire("7 x ", i, " = ", produit)
    FinPour
Fin
\end{lstlisting}
\end{exemplebox}

À chaque tour, le compteur \texttt{i} vaut successivement $1, 2, 3, \dots, 10$, et l'instruction \texttt{Ecrire} affiche la ligne correspondante : « 7 x 1 = 7 », puis « 7 x 2 = 14 », ..., jusqu'à « 7 x 10 = 70 ». Remarque que le compteur n'est pas seulement un « numéro de tour » : on l'\textbf{utilise dans les calculs} du corps de la boucle. C'est très fréquent et très pratique.

\courssub{Équivalence entre « Pour » et « TantQue »}

On peut se poser la question : la boucle \texttt{Pour} est-elle vraiment une nouveauté, ou bien peut-on la « traduire » avec ce que l'on connaît déjà ? La réponse est importante :

\begin{propriete}[Équivalence des deux boucles]
Toute boucle \cle{Pour} peut être réécrite avec une boucle \cle{TantQue}. Il suffit de faire « à la main » ce que le \texttt{Pour} fait automatiquement :
\begin{enumerate}
  \item \textbf{initialiser} le compteur à la valeur de départ, \emph{avant} la boucle ;
  \item écrire la condition \texttt{compteur <= valeur\_fin} dans le \texttt{TantQue} ;
  \item \textbf{incrémenter} le compteur (\texttt{compteur <- compteur + 1}) à la \emph{fin} du corps de la boucle.
\end{enumerate}
La réciproque est fausse : un \texttt{TantQue} dont on ne connaît pas le nombre de tours à l'avance (par exemple « tant que le mot de passe saisi est faux ») ne peut pas toujours s'écrire avec un \texttt{Pour}.
\end{propriete}

\begin{exemplebox}[Réécriture guidée : la table de 7 avec un TantQue]
Reprenons l'algorithme de la table de 7 et appliquons la méthode de traduction, étape par étape :
\begin{lstlisting}
Algorithme TableDe7Bis
Variables
    i, produit : Entier
Debut
    i <- 1                    // etape 1 : initialisation du compteur
    TantQue i <= 10 Faire     // etape 2 : condition de continuation
        produit <- 7 * i
        Ecrire("7 x ", i, " = ", produit)
        i <- i + 1            // etape 3 : incrementation (a ne pas oublier !)
    FinTantQue
Fin
\end{lstlisting}
Ce programme affiche exactement les mêmes 10 lignes que la version avec \texttt{Pour}. On retrouve, écrites explicitement, les trois opérations que la boucle \texttt{Pour} réalisait automatiquement.
\end{exemplebox}

Le schéma suivant résume visuellement cette correspondance :

\begin{popfigure}
\begin{center}
\begin{tikzpicture}[
  bloc/.style={rectangle, draw=popblue, thick, fill=popblueL, rounded corners=2pt, align=center, minimum height=0.85cm, text width=4.6cm, font=\small},
  bloc2/.style={rectangle, draw=poporange, thick, fill=poporangeL, rounded corners=2pt, align=center, minimum height=0.85cm, text width=4.6cm, font=\small},
  lab/.style={font=\small\bfseries, text=popdark},
  fleche/.style={-{Stealth[length=2.5mm]}, thick, popdark}
]
% Colonne TantQue (gauche)
\node[lab] at (0,4.6) {Boucle TantQue (manuel)};
\node[bloc, align=center] (t1) at (0,3.6){i $\leftarrow$ 1\\ \emph{(initialisation explicite)}};
\node[bloc, align=center] (t2) at (0,2.2){Condition : i $\leq$ 10 ?\\ \emph{(testée par le programmeur)}};
\node[bloc, align=center] (t3) at (0,0.8){Corps de la boucle\\ (instructions)};
\node[bloc, align=center] (t4) at (0,-0.6){i $\leftarrow$ i + 1\\ \emph{(incrémentation explicite)}};
\draw[fleche] (t1) -- (t2);
\draw[fleche] (t2) -- (t3);
\draw[fleche] (t3) -- (t4);
\draw[fleche] (t4.west) -- ++(-1.1,0) |- (t2.west);
% Colonne Pour (droite)
\node[lab] at (8.6,4.6) {Boucle Pour (automatique)};
\node[bloc2, align=center] (p1) at (8.6,2.9){Pour i de 1 a 10\\ \emph{(initialisation, test et\\ compteur gérés tout seuls)}};
\node[bloc2, align=center] (p2) at (8.6,0.8){Corps de la boucle\\ (instructions)};
\draw[fleche] (p1) -- (p2);
\draw[fleche] (p2.east) -- ++(1.1,0) |- (p1.east) node[midway, right, font=\small\itshape, text=popmuted, align=left]{compteur\\ automatique};
% Equivalence
\node[font=\Large\bfseries, text=popgreen] at (4.3,2.2) {$\equiv$};
\end{tikzpicture}
\end{center}
\legende{Comparaison TantQue / Pour : la boucle Pour gère automatiquement l'initialisation, le test et l'incrémentation du compteur, que le TantQue laisse à la charge du programmeur.}
\end{popfigure}

\courssub{Comment choisir entre « Pour » et « TantQue » ?}

\begin{methode}[Choisir la bonne boucle]
Face à un problème qui demande de répéter des instructions, pose-toi une seule question : \textbf{« Est-ce que je connais le nombre de répétitions avant de commencer ? »}
\begin{enumerate}
  \item \textbf{Oui, le nombre de tours est connu} (les 25 élèves de la classe, les 12 mois de l'année, les entiers de 1 à 100) $\Rightarrow$ utilise la boucle \cle{Pour}. Elle est plus courte à écrire et plus sûre (pas de risque d'oublier l'incrémentation).
  \item \textbf{Non, l'arrêt dépend d'une condition} qu'on ne peut pas prévoir (tant que le mot de passe est faux, tant que la saisie est invalide, tant que le stock n'est pas épuisé) $\Rightarrow$ utilise la boucle \cle{TantQue}.
  \item En cas de doute, le \texttt{TantQue} convient toujours (puisque tout \texttt{Pour} peut s'écrire avec un \texttt{TantQue}), mais le \texttt{Pour} rend l'algorithme plus lisible quand le nombre de tours est fixe.
\end{enumerate}
\end{methode}

\begin{center}
\begin{tabularx}{\linewidth}{|l|X|X|}
\hline
\rowcolor{popblueL} \textbf{Critère} & \textbf{Boucle Pour} & \textbf{Boucle TantQue} \\
\hline
Nombre de répétitions & Connu à l'avance & Inconnu : dépend d'une condition \\
\hline
Compteur & Initialisé et incrémenté automatiquement & Géré manuellement par le programmeur \\
\hline
Risque de boucle infinie & Quasi nul & Réel si l'on oublie de faire évoluer la condition \\
\hline
Exemple typique & Somme de 1 à 100, table de 7 & Saisie d'un mot de passe correct \\
\hline
\end{tabularx}
\end{center}

\begin{attention}[Ne jamais modifier le compteur à l'intérieur d'une boucle Pour]
Une erreur fréquente consiste à écrire, dans le corps d'une boucle \texttt{Pour}, une instruction qui modifie le compteur, par exemple :
\begin{lstlisting}
Pour i de 1 a 10 Faire
    i <- i + 1        // INTERDIT : le compteur est deja gere par le Pour !
    Ecrire(i)
FinPour
\end{lstlisting}
Le compteur d'une boucle \texttt{Pour} est géré \textbf{automatiquement} : le modifier à la main rend le comportement imprévisible (tours sautés, boucle qui ne s'arrête pas comme prévu, résultats faux). Règle d'or : dans le corps d'un \texttt{Pour}, on \textbf{lit} la valeur du compteur (pour calculer ou afficher), mais on ne lui affecte \textbf{jamais} de nouvelle valeur. Si tu sens le besoin de contrôler toi-même la progression, c'est le signe qu'il faut utiliser un \texttt{TantQue}.
\end{attention}

\begin{exercice}[À toi de jouer]
\begin{enumerate}
  \item Écris un algorithme qui calcule et affiche la somme des entiers pairs de 2 à 50 (utilise un compteur qui progresse de 2 en 2, ou un calcul adapté).
  \item Le directeur d'un cybercafé veut afficher, pour chacun des 15 postes de la salle, le message « Poste n° X : contrôle effectué ». Rédige l'algorithme avec la boucle la plus adaptée et justifie ton choix.
  \item Réécris l'algorithme de la question 2 avec une boucle \texttt{TantQue}, en appliquant la méthode de traduction vue dans cette section.
\end{enumerate}
\end{exercice}

\begin{corrige}[Corrigé de l'exercice]
\begin{enumerate}
  \item On initialise \texttt{somme <- 0}, puis \texttt{Pour i de 2 a 50 Pas 2 Faire somme <- somme + i FinPour}. Le résultat attendu est $2 + 4 + \dots + 50 = 650$.
  \item Le nombre de passages est connu (15 postes) : on choisit donc la boucle \texttt{Pour}. Algorithme : \texttt{Pour i de 1 a 15 Faire Ecrire("Poste n° ", i, " : contrôle effectué") FinPour}.
  \item On initialise \texttt{i <- 1}, la condition est \texttt{TantQue i <= 15}, et on n'oublie pas \texttt{i <- i + 1} à la fin du corps de la boucle.
\end{enumerate}
\end{corrige}

\begin{aretenir}[L'essentiel de la section]
\begin{itemize}
  \item La boucle \cle{Pour} s'utilise lorsque le \textbf{nombre de répétitions est connu à l'avance} ; sa syntaxe est : \texttt{Pour compteur de début à fin Faire ... FinPour}.
  \item Le \cle{compteur} est initialisé et augmenté \textbf{automatiquement} (pas de $1$ par défaut) ; on ne le modifie jamais dans le corps de la boucle.
  \item Toute boucle \texttt{Pour} peut se réécrire avec un \texttt{TantQue} (initialisation avant, condition, incrémentation à la fin), mais l'inverse n'est pas toujours possible.
  \item Pour accumuler un résultat tour après tour (somme, total), on utilise un \cle{accumulateur} initialisé à $0$ : c'est ainsi qu'on obtient $1 + 2 + \dots + 100 = 5050$, résultat confirmé par la formule $\dfrac{n(n+1)}{2}$.
\end{itemize}
\end{aretenir}

\coursec{L'algorigramme : représentation graphique}

Dans les sections précédentes, tu as appris à écrire des algorithmes en pseudo-code : des instructions écrites en français, avec des structures alternatives (\texttt{Si...Alors...Sinon}) et des structures itératives (\texttt{Tant que}, \texttt{Pour}). Ce texte est précis, mais il a un défaut : pour comprendre le déroulement d'un algorithme long, il faut lire chaque ligne, une par une. Un dessin, lui, se comprend d'un coup d'œil. C'est exactement le rôle de l'\cle{algorigramme} : transformer le texte d'un algorithme en un schéma clair, où l'on suit le cheminement du programme comme on suit un plan de métro.

\courssub{Qu'est-ce qu'un algorigramme ?}

Imagine que tu dois expliquer à un camarade comment ton téléphone décide d'afficher « Admis » ou « Refusé » sur un bulletin. Tu peux lui réciter le pseudo-code... ou lui dessiner un schéma : « on entre la moyenne, on arrive à une question : la moyenne est-elle supérieure ou égale à 10 ? Si oui, on va par là ; si non, par ici. » Ce schéma, c'est un algorigramme.

\begin{definition}[Algorigramme]
Un \cle{algorigramme} (ou \emph{organigramme}) est la \textbf{représentation graphique d'un algorithme} à l'aide de \textbf{symboles normalisés} reliés par des \textbf{flèches}. Chaque symbole correspond à un type d'opération (début, lecture, traitement, test, affichage), et les flèches indiquent l'\textbf{ordre d'exécution} des opérations.
\end{definition}

L'algorigramme et le pseudo-code décrivent donc \emph{le même algorithme}, sous deux formes différentes : l'un est un dessin, l'autre est un texte. Savoir passer de l'un à l'autre est une compétence essentielle, que nous travaillerons à la fin de cette section.

\courssub{Les symboles normalisés}

Pour que tout le monde comprenne un algorigramme, les symboles sont \textbf{normalisés} : ils ont une forme et un rôle précis, identiques partout dans le monde. En voici les quatre symboles de base à connaître, plus les flèches de liaison.

\begin{popfigure}
\begin{center}
\begin{tikzpicture}[font=\small]
  \node[draw=popblue, fill=popblueL, ellipse, minimum width=2.6cm, minimum height=1cm, thick] (oval) at (0,0) {Début / Fin};
  \node[draw=popgreen, fill=popgreenL, trapezium, trapezium left angle=70, trapezium right angle=110, minimum width=3.2cm, minimum height=1cm, thick] (para) at (5,0) {Lire / Écrire};
  \node[draw=poporange, fill=poporangeL, rectangle, minimum width=3.2cm, minimum height=1cm, thick] (rect) at (10,0) {Traitement};
  \node[draw=poppurple, fill=poppurpleL, diamond, aspect=2.4, minimum width=3.4cm, thick, inner sep=1pt] (los) at (2.5,-2.6) {Test (condition)};
  \draw[-{Stealth[length=3mm]}, very thick, popdark] (7,-2.6) -- (9.5,-2.6) node[midway, above]{Flèche de liaison};
\end{tikzpicture}
\end{center}
\legende{Planche des symboles normalisés d'un algorigramme.}
\end{popfigure}

\begin{aretenir}[Rôle de chaque symbole]
\begin{itemize}
  \item L'\cle{ovale} marque le \textbf{Début} ou la \textbf{Fin} de l'algorithme. Il y a toujours \emph{un seul} Début et au moins une Fin.
  \item Le \cle{parallélogramme} représente les \textbf{entrées et sorties} : \texttt{Lire} (l'utilisateur saisit une donnée) et \texttt{Écrire} (l'algorithme affiche un résultat).
  \item Le \cle{rectangle} représente un \textbf{traitement} : un calcul ou une affectation (par exemple $S \leftarrow S + i$).
  \item Le \cle{losange} représente un \textbf{test} (une condition). Il possède \textbf{exactement deux sorties} : une branche \emph{Oui} et une branche \emph{Non}.
  \item Les \cle{flèches} relient les symboles et indiquent le sens d'exécution.
\end{itemize}
\end{aretenir}

\begin{attention}[Erreur fréquente : le test n'est pas un rectangle]
L'erreur la plus classique consiste à dessiner un rectangle pour une condition comme « moyenne $\geq 10$ ». C'est faux : un rectangle a \emph{une seule} sortie, car un traitement ne pose aucune question. Un \textbf{test est toujours un losange}, avec \textbf{exactement deux flèches sortantes} étiquetées \emph{Oui} et \emph{Non}. De même, n'utilise jamais un losange pour un calcul : un calcul ne se pose pas de question, il s'exécute, donc c'est un rectangle.
\end{attention}

\courssub{Les flèches et le sens de lecture}

Un algorigramme se lit \textbf{de haut en bas} : le symbole \emph{Début} est placé en haut, et l'exécution descend de flèche en flèche jusqu'au symbole \emph{Fin}. Il n'y a qu'une seule exception à cette règle : les \textbf{boucles}, où une flèche \emph{remonte} vers un symbole situé plus haut pour répéter des instructions. On parle alors de \emph{retour en arrière}.

\begin{propriete}[Règles de construction des flèches]
\begin{itemize}
  \item Chaque symbole (sauf \emph{Fin}) possède une flèche sortante ; chaque symbole (sauf \emph{Début}) possède une flèche entrante.
  \item Un losange possède exactement deux flèches sortantes, étiquetées \emph{Oui} et \emph{Non}.
  \item Deux flèches ne se croisent pas si l'on peut l'éviter ; le schéma doit rester lisible.
  \item Le sens général de lecture est du haut vers le bas ; seule une boucle autorise une flèche qui remonte.
\end{itemize}
\end{propriete}

\courssub{Représenter une alternative et une boucle}

\textbf{L'alternative} (\texttt{Si...Alors...Sinon}) se traduit par un losange d'où partent deux branches : la branche \emph{Oui} mène au traitement à faire quand la condition est vraie, la branche \emph{Non} mène à l'autre traitement. Les deux branches \textbf{se rejoignent ensuite} pour continuer la suite de l'algorithme.

\textbf{La boucle} (\texttt{Tant que} ou \texttt{Pour}) se traduit par un losange dont la branche \emph{Oui} effectue un traitement puis \textbf{revient en arrière} vers le losange : tant que la condition est vraie, on recommence. La branche \emph{Non} quitte la boucle et poursuit l'algorithme.

\begin{methode}[Construire un algorigramme en quatre étapes]
Pour dessiner l'algorigramme d'un problème, procède toujours dans le même ordre :
\begin{enumerate}
  \item \textbf{Identifier les entrées} : quelles données l'utilisateur doit-il saisir ? Ce seront des parallélogrammes \texttt{Lire}.
  \item \textbf{Identifier les traitements} : quels calculs et affectations sont nécessaires ? Ce seront des rectangles.
  \item \textbf{Identifier les tests} : où l'algorithme doit-il prendre une décision ou répéter des actions ? Ce seront des losanges à deux branches.
  \item \textbf{Identifier les sorties} : quels résultats affiche-t-on ? Ce seront des parallélogrammes \texttt{Écrire}. Relie ensuite tous les symboles par des flèches, du \emph{Début} vers la \emph{Fin}.
\end{enumerate}
\end{methode}

\courssub{Exemple complet : l'algorithme Admis/Refusé}

Reprenons notre algorithme classique : on saisit la moyenne d'un élève ; si elle est supérieure ou égale à 10, on affiche « Admis », sinon « Refusé ». Voici d'abord le pseudo-code, puis sa traduction graphique.

\begin{lstlisting}
Algorithme AdmisRefuse
Variable moyenne : reel
Debut
    Lire(moyenne)
    Si moyenne >= 10 Alors
        Ecrire("Admis")
    Sinon
        Ecrire("Refuse")
    FinSi
Fin
\end{lstlisting}

\begin{exemplebox}[Conversion pseudo-code $\rightarrow$ algorigramme : Admis/Refusé]
Appliquons la méthode en quatre étapes. Entrée : la moyenne (parallélogramme \texttt{Lire}). Traitement : aucun calcul ici. Test : « moyenne $\geq 10$ ? » (losange à deux branches). Sorties : « Admis » sur la branche \emph{Oui}, « Refusé » sur la branche \emph{Non}. Les deux branches se rejoignent avant le \emph{Fin}. On obtient l'algorigramme ci-dessous.
\end{exemplebox}

\begin{popfigure}
\begin{center}
\begin{tikzpicture}[font=\small, node distance=0.9cm,
  startstop/.style={ellipse, draw=popblue, fill=popblueL, thick, minimum width=2.2cm, minimum height=0.9cm},
  io/.style={trapezium, trapezium left angle=70, trapezium right angle=110, draw=popgreen, fill=popgreenL, thick, minimum width=3.4cm, minimum height=0.9cm},
  test/.style={diamond, aspect=2.6, draw=poppurple, fill=poppurpleL, thick, inner sep=1pt},
  fleche/.style={-{Stealth[length=2.5mm]}, thick, popdark}]
  \node[startstop] (deb) {Début};
  \node[io, below=of deb] (lire) {Lire(moyenne)};
  \node[test, below=of lire] (cond) {moyenne $\geq 10$ ?};
  \node[io, below left=1.1cm and 0.6cm of cond] (admis) {Écrire(« Admis »)};
  \node[io, below right=1.1cm and 0.6cm of cond] (refus) {Écrire(« Refusé »)};
  \node[startstop, below=2.6cm of cond] (fin) {Fin};
  \draw[fleche] (deb) -- (lire);
  \draw[fleche] (lire) -- (cond);
  \draw[fleche] (cond) -| node[near start, above left]{Oui} (admis);
  \draw[fleche] (cond) -| node[near start, above right]{Non} (refus);
  \draw[fleche] (admis) |- (fin);
  \draw[fleche] (refus) |- (fin);
\end{tikzpicture}
\end{center}
\legende{Algorigramme de l'algorithme Admis/Refusé : le losange possède deux branches qui se rejoignent avant la fin.}
\end{popfigure}

Remarque comme le schéma se lit naturellement : on descend, on arrive à la question, on choisit un chemin, puis les deux chemins se rejoignent. C'est exactement le déroulement du pseudo-code, rendu visible.

\courssub{Exemple complet : la somme des entiers de 1 à n}

Passons à un algorithme avec boucle : calculer $S = 1 + 2 + 3 + \dots + n$, où $n$ est saisi par l'utilisateur. Le pseudo-code (avec une boucle \texttt{Pour}) est le suivant :

\begin{lstlisting}
Algorithme Somme1aN
Variables n, i, S : entiers
Debut
    Lire(n)
    S <- 0
    Pour i allant de 1 a n Faire
        S <- S + i
    FinPour
    Ecrire("La somme est : ", S)
Fin
\end{lstlisting}

Pour le traduire en algorigramme, on représente la boucle par un losange « $i \leq n$ ? » : tant que la réponse est \emph{Oui}, on exécute le traitement $S \leftarrow S + i$, on incrémente $i$, et une flèche \textbf{remonte} vers le losange. Quand la réponse devient \emph{Non}, on sort de la boucle et on affiche le résultat.

\begin{popfigure}
\begin{center}
\begin{tikzpicture}[font=\small, node distance=0.8cm,
  startstop/.style={ellipse, draw=popblue, fill=popblueL, thick, minimum width=2.2cm, minimum height=0.9cm},
  io/.style={trapezium, trapezium left angle=70, trapezium right angle=110, draw=popgreen, fill=popgreenL, thick, minimum width=3.4cm, minimum height=0.9cm},
  proc/.style={rectangle, draw=poporange, fill=poporangeL, thick, minimum width=3cm, minimum height=0.9cm},
  test/.style={diamond, aspect=2.6, draw=poppurple, fill=poppurpleL, thick, inner sep=1pt},
  fleche/.style={-{Stealth[length=2.5mm]}, thick, popdark}]
  \node[startstop] (deb) {Début};
  \node[io, below=of deb] (lire) {Lire(n)};
  \node[proc, below=of lire] (init) {$S \leftarrow 0$ ; $i \leftarrow 1$};
  \node[test, below=of init] (cond) {$i \leq n$ ?};
  \node[proc, right=1.6cm of cond] (corps) {$S \leftarrow S + i$};
  \node[proc, above=0.8cm of corps] (incr) {$i \leftarrow i + 1$};
  \node[io, below=of cond] (ecrire) {Écrire(S)};
  \node[startstop, below=of ecrire] (fin) {Fin};
  \draw[fleche] (deb) -- (lire);
  \draw[fleche] (lire) -- (init);
  \draw[fleche] (init) -- (cond);
  \draw[fleche] (cond) -- node[above]{Oui} (corps);
  \draw[fleche] (corps) -- (incr);
  \draw[fleche] (incr.west) -| ($(cond.north)+(0.6,0)$) -- (cond.north);
  \draw[fleche] (cond) -- node[right]{Non} (ecrire);
  \draw[fleche] (ecrire) -- (fin);
\end{tikzpicture}
\end{center}
\legende{Algorigramme de la somme $1 + 2 + \dots + n$ : la branche \emph{Oui} du losange remonte vers le test, formant la boucle.}
\end{popfigure}

\begin{exoresolu}[De l'algorigramme au pseudo-code]
On donne l'algorigramme suivant, décrit symbole par symbole : \emph{Début} ; \texttt{Lire(a)} ; \texttt{Lire(b)} ; losange « $a > b$ ? » ; branche \emph{Oui} : \texttt{Écrire(a)} ; branche \emph{Non} : \texttt{Écrire(b)} ; \emph{Fin}. Écris le pseudo-code correspondant et dis ce que fait cet algorithme.
\tcblower
\textbf{Solution.} On suit les flèches de haut en bas : deux lectures, puis un test à deux branches, puis la fin. Le pseudo-code est :
\begin{lstlisting}
Algorithme Maximum
Variables a, b : reels
Debut
    Lire(a)
    Lire(b)
    Si a > b Alors
        Ecrire(a)
    Sinon
        Ecrire(b)
    FinSi
Fin
\end{lstlisting}
Cet algorithme lit deux nombres et affiche \textbf{le plus grand des deux}. Vérification : si $a = 7$ et $b = 12$, le test « $7 > 12$ ? » est faux, on suit la branche \emph{Non} et on affiche $12$ : c'est bien le maximum.
\end{exoresolu}

\begin{attention}[Pièges fréquents dans les boucles dessinées]
\begin{itemize}
  \item \textbf{Oublier l'initialisation} : le rectangle $S \leftarrow 0$ ; $i \leftarrow 1$ doit être placé \emph{avant} le losange, jamais dans la boucle, sinon $S$ serait remis à zéro à chaque tour.
  \item \textbf{Oublier l'incrémentation} : sans le rectangle $i \leftarrow i + 1$, la condition $i \leq n$ reste vraie pour toujours : la flèche remonte indéfiniment, c'est une \emph{boucle infinie}.
  \item \textbf{Inverser Oui et Non} : la branche qui \emph{remonte} doit correspondre au cas où l'on continue (condition vraie), la branche qui \emph{descend} au cas où l'on s'arrête.
\end{itemize}
\end{attention}

\begin{savaistu}[Les algorigrammes en entreprise]
Dans les entreprises et les administrations — y compris au Cameroun, par exemple dans les services informatiques des banques ou de la CAMTEL — les développeurs dessinent souvent des algorigrammes \textbf{avant} d'écrire le moindre programme. Ce schéma sert de document de travail : il permet à toute l'équipe de discuter de la logique du futur logiciel, de repérer les oublis et les cas non prévus, puis de coder ensuite dans n'importe quel langage. Un bon dessin vaut parfois mille lignes de code !
\end{savaistu}

\begin{exercice}[À toi de jouer]
\begin{enumerate}
  \item Dessine l'algorigramme de l'algorithme qui lit l'âge d'un élève et affiche « Mineur » si l'âge est inférieur à 18, « Majeur » sinon.
  \item Dessine l'algorigramme de l'algorithme qui calcule et affiche la table de multiplication de 7 (de $7 \times 1$ à $7 \times 10$) à l'aide d'une boucle.
  \item Pour chacun de tes algorigrammes, écris le pseudo-code correspondant.
\end{enumerate}
\end{exercice}

\begin{corrige}[Exercice]
\textbf{1.} Symboles dans l'ordre : \emph{Début} ; parallélogramme \texttt{Lire(age)} ; losange « age $< 18$ ? » ; branche \emph{Oui} : parallélogramme \texttt{Écrire(« Mineur »)} ; branche \emph{Non} : parallélogramme \texttt{Écrire(« Majeur »)} ; les deux branches se rejoignent ; \emph{Fin}. Pseudo-code :
\begin{lstlisting}
Algorithme Majorite
Variable age : entier
Debut
    Lire(age)
    Si age < 18 Alors
        Ecrire("Mineur")
    Sinon
        Ecrire("Majeur")
    FinSi
Fin
\end{lstlisting}
\textbf{2.} L'algorigramme détaille le fonctionnement de la boucle \texttt{Pour} : \emph{Début} ; rectangle $i \leftarrow 1$ ; losange « $i \leq 10$ ? » ; branche \emph{Oui} : rectangle $p \leftarrow 7 \times i$, puis parallélogramme \texttt{Écrire(p)}, puis rectangle $i \leftarrow i + 1$, puis flèche de retour vers le losange ; branche \emph{Non} : \emph{Fin}. Pseudo-code :
\begin{lstlisting}
Algorithme TableDe7
Variables i, p : entiers
Debut
    Pour i allant de 1 a 10 Faire
        p <- 7 * i
        Ecrire("7 x ", i, " = ", p)
    FinPour
Fin
\end{lstlisting}
\textbf{3.} Les pseudo-codes ci-dessus correspondent exactement aux algorigrammes : chaque parallélogramme donne un \texttt{Lire} ou un \texttt{Écrire}, chaque rectangle une affectation, chaque losange un \texttt{Si} ou une boucle.
\end{corrige}

\begin{aretenir}[L'essentiel de la section]
\begin{itemize}
  \item Un \cle{algorigramme} est la représentation graphique d'un algorithme : symboles normalisés reliés par des flèches, lu de haut en bas.
  \item Ovale = Début/Fin ; parallélogramme = Lire/Écrire ; rectangle = traitement ; losange = test avec exactement deux sorties \emph{Oui}/\emph{Non}.
  \item Une alternative se dessine par un losange dont les deux branches se rejoignent ; une boucle par un losange dont la branche \emph{Oui} remonte vers le test.
  \item Algorigramme et pseudo-code sont deux présentations du même algorithme : il faut savoir passer de l'une à l'autre.
\end{itemize}
\end{aretenir}

\coursec{Résolution de problèmes concrets et justification}

Dans la section précédente, nous avons appris à \emph{représenter graphiquement} un algorithme grâce à l'algorigramme : rectangles pour les traitements, losanges pour les tests, flèches pour l'enchaînement des étapes. Nous savons donc maintenant écrire une séquence d'instructions, une alternative \cle{Si...Alors...Sinon}, une boucle \cle{Tant que} ou \cle{Pour}, et dessiner tout cela. Mais un algorithme ne s'écrit jamais pour le plaisir : il sert à \textbf{résoudre un problème concret}. C'est exactement ce que l'on te demandera au BEPC : face à une situation de la vie courante (une facture, une moyenne de notes, un petit jeu), sauras-tu proposer une démarche complète et \emph{justifiée} ? C'est l'objet de cette dernière section, qui est aussi la synthèse de toute la leçon.

\courssub{La démarche complète de résolution d'un problème}

Quand un élève reçoit un énoncé, sa première tentation est d'écrire immédiatement des instructions. C'est une erreur : on ne construit pas une maison sans plan. Un bon informaticien suit toujours une démarche en étapes, que l'on peut apprendre par cœur.

\begin{methode}[Résoudre un problème par un algorithme : les 6 étapes]
\begin{enumerate}
\item \textbf{Comprendre le problème.} Lire l'énoncé plusieurs fois, le reformuler avec ses propres mots, repérer ce qui est demandé.
\item \textbf{Identifier les entrées et les sorties.} Les \cle{entrées} sont les données fournies (exemple : la consommation en kWh) ; les \cle{sorties} sont les résultats attendus (exemple : le montant à payer).
\item \textbf{Choisir les structures de contrôle.} Y a-t-il un cas particulier ? $\rightarrow$ structure alternative. Y a-t-il une action répétée ? $\rightarrow$ structure itérative (\texttt{Tant que} si l'on ne connaît pas le nombre de répétitions, \texttt{Pour} si on le connaît).
\item \textbf{Écrire l'algorithme} en pseudo-code, avec un nom, des variables déclarées, un \texttt{Début} et une \texttt{Fin}.
\item \textbf{Tracer l'algorithme} : le faire « tourner » à la main sur au moins deux jeux de valeurs tests, en notant l'évolution des variables dans un tableau.
\item \textbf{Conclure et justifier} : répondre à la question posée et expliquer \emph{pourquoi} on a choisi telle structure plutôt qu'une autre.
\end{enumerate}
\end{methode}

\begin{aretenir}[La trace d'exécution]
La \cle{trace} (ou table d'exécution) est un tableau où chaque colonne correspond à une variable et chaque ligne à une étape de l'exécution. Elle permet de \textbf{vérifier} qu'un algorithme est correct \emph{avant} de le faire exécuter par une machine.
\end{aretenir}

\courssub{Problème 1 : la facture d'électricité}

\textbf{Situation.} Une compagnie d'électricité facture le kilowattheure (kWh) à 50 FCFA. Mais pour décourager le gaspillage, si la consommation mensuelle dépasse 110~kWh, le tarif passe à 75 FCFA par kWh pour \emph{toute} la consommation. Écrire un algorithme qui lit la consommation et affiche le montant à payer.

\textbf{Étapes 1 et 2 — comprendre, identifier.} Entrée : la consommation $C$ (en kWh). Sortie : le montant $M$ (en FCFA). Il y a deux cas : c'est donc une \textbf{alternative}.

\textbf{Étape 4 — l'algorithme.}

\begin{lstlisting}
Algorithme FactureElectricite
Variables
    C, M : Reel
Debut
    Ecrire("Entrez la consommation en kWh : ")
    Lire(C)
    Si C > 110 Alors
        M <- C * 75      // tarif majoré
    Sinon
        M <- C * 50      // tarif simple
    FinSi
    Ecrire("Montant a payer : ", M, " FCFA")
Fin
\end{lstlisting}

\begin{exemplebox}[Trace complète pour une consommation de 150 kWh]
Faisons « tourner » l'algorithme avec $C = 150$ :

\begin{center}
\begin{tabularx}{\linewidth}{|l|X|X|X|}
\hline
\rowcolor{popblueL} \textbf{Étape} & \textbf{C} & \textbf{Test $C>110$} & \textbf{M} \\
\hline
Lecture & 150 & --- & --- \\
\hline
Test & 150 & Vrai & --- \\
\hline
Calcul & 150 & --- & $150 \times 75 = 11\,250$ \\
\hline
Affichage & \multicolumn{3}{X|}{Montant à payer : 11\,250 FCFA} \\
\hline
\end{tabularx}
\end{center}

Vérifions avec un second jeu de valeurs, $C = 80$ : le test $80 > 110$ est faux, donc $M = 80 \times 50 = 4\,000$ FCFA. L'algorithme se comporte correctement dans les deux cas.
\end{exemplebox}

Voici l'algorigramme correspondant, qui montre clairement la structure alternative :

\begin{popfigure}
\begin{center}
\begin{tikzpicture}[
  font=\small,
  startstop/.style={rounded rectangle, draw=popdark, fill=popgoldL, minimum width=2.6cm, minimum height=0.8cm},
  io/.style={trapezium, trapezium left angle=70, trapezium right angle=110, draw=popdark, fill=popblueL, minimum height=0.8cm},
  process/.style={rectangle, draw=popdark, fill=popgreenL, minimum width=3cm, minimum height=0.8cm},
  decision/.style={diamond, draw=popdark, fill=poppinkL, aspect=2.4, inner sep=1pt},
  fleche/.style={-{Stealth[length=2.5mm]}, thick, popdark}
]
\node[startstop] (debut) {Début};
\node[io, below=0.5cm of debut] (lire) {Lire C};
\node[decision, below=0.6cm of lire] (test) {C $>$ 110 ?};
\node[process, below left=1.1cm and 0.4cm of test] (maj) {M $\leftarrow$ C $\times$ 75};
\node[process, below right=1.1cm and 0.4cm of test] (sim) {M $\leftarrow$ C $\times$ 50};
\node[io, below=2.6cm of test] (aff) {Afficher M};
\node[startstop, below=0.5cm of aff] (fin) {Fin};
\draw[fleche] (debut) -- (lire);
\draw[fleche] (lire) -- (test);
\draw[fleche] (test.west) -| node[pos=0.25, above]{Oui} (maj.north);
\draw[fleche] (test.east) -| node[pos=0.25, above]{Non} (sim.north);
\draw[fleche] (maj.south) |- (aff.west);
\draw[fleche] (sim.south) |- (aff.east);
\draw[fleche] (aff) -- (fin);
\end{tikzpicture}
\end{center}
\legende{Algorigramme du calcul de la facture d'électricité : le losange matérialise l'alternative.}
\end{popfigure}

\textbf{Étape 6 — justification.} « J'ai utilisé une structure \texttt{Si...Alors...Sinon} car le calcul du montant dépend d'une condition (la consommation dépasse-t-elle 110~kWh ?). Les deux cas sont traités de façon exclusive : le tarif majoré s'applique à toute la consommation, pas seulement au dépassement. »

\courssub{Problème 2 : le jeu de devinette}

\textbf{Situation.} L'ordinateur « pense » à un nombre secret entre 1 et 100. Le joueur propose un nombre ; tant qu'il n'a pas trouvé, on lui répond « trop grand » ou « trop petit » et on lui redemande. Décrire l'idée de l'algorithme (sans le programmer sur machine).

\textbf{Choix de la structure.} Ici, on ne sait \emph{pas à l'avance} combien d'essais seront nécessaires : le joueur peut trouver du premier coup... ou au bout de dix essais. C'est donc une boucle \texttt{Tant que}, dont la condition d'arrêt est « le nombre proposé est égal au nombre secret ».

\begin{lstlisting}
Algorithme Devinette
Variables
    secret, proposition : Entier
Debut
    secret <- 42          // nombre choisi par l'ordinateur
    Lire(proposition)
    Tant que proposition <> secret Faire
        Si proposition > secret Alors
            Ecrire("Trop grand !")
        Sinon
            Ecrire("Trop petit !")
        FinSi
        Lire(proposition)   // on redemande
    FinTantQue
    Ecrire("Bravo, trouve !")
Fin
\end{lstlisting}

\begin{attention}[La boucle infinie]
Dans une boucle \texttt{Tant que}, il faut \textbf{toujours} qu'une instruction à l'intérieur de la boucle puisse rendre la condition fausse. Ici, c'est le \texttt{Lire(proposition)} placé \emph{dans} la boucle qui permet de progresser. Si on l'oublie, la condition ne change jamais : la boucle ne s'arrête plus, c'est une \cle{boucle infinie}.
\end{attention}

\textbf{Justification type.} « J'ai choisi un \texttt{Tant que} car le nombre de répétitions est inconnu à l'avance : on répète tant que la condition "nombre non trouvé" est vraie. Un \texttt{Pour} aurait été inadapté car il exige de connaître le nombre d'itérations. »

\courssub{Problème 3 : moyenne générale et mention}

\textbf{Situation.} Un professeur veut un algorithme qui lit la moyenne générale d'un élève (sur 20) et affiche la décision : « Ajourné » si la moyenne est inférieure à 10 ; sinon « Admis », avec la mention « Passable » si la moyenne est inférieure à 12, « Bien » si elle est inférieure à 14, et « Excellent » à partir de 14.

\textbf{Choix des structures.} Il y a d'abord une grande alternative (admis ou ajourné), puis, \emph{à l'intérieur} du cas « admis », d'autres tests : ce sont des \cle{alternatives imbriquées}.

\begin{lstlisting}
Algorithme Mention
Variables
    moy : Reel
Debut
    Ecrire("Entrez la moyenne : ")
    Lire(moy)
    Si moy < 10 Alors
        Ecrire("Ajourné")
    Sinon
        Ecrire("Admis")
        Si moy < 12 Alors
            Ecrire("Mention Passable")
        Sinon
            Si moy < 14 Alors
                Ecrire("Mention Bien")
            Sinon
                Ecrire("Mention Excellent")
            FinSi
        FinSi
    FinSi
Fin
\end{lstlisting}

\begin{exoresolu}[Trace de l'algorithme « Mention »]
Énoncé : tracer l'algorithme pour $moy = 13,5$ puis pour $moy = 8$.
\tcblower
Pour $moy = 13,5$ : le test $13,5 < 10$ est faux $\rightarrow$ affichage « Admis » ; le test $13,5 < 12$ est faux ; le test $13,5 < 14$ est vrai $\rightarrow$ affichage « Mention Bien ».

Pour $moy = 8$ : le test $8 < 10$ est vrai $\rightarrow$ affichage « Ajourné » ; les tests de mention ne sont \emph{jamais exécutés}, car ils sont dans la branche \texttt{Sinon}.

L'algorithme est correct : chaque cas possible donne le résultat attendu. Remarque : on teste d'abord \texttt{moy < 12} puis \texttt{moy < 14} dans cet ordre, car un élève arrivé au second test a déjà une moyenne supérieure ou égale à 12.
\end{exoresolu}

\courssub{Vérifier par la trace et rédiger la justification}

\begin{attention}[Erreur fréquente : ne jamais tester son algorithme]
Beaucoup d'élèves écrivent un algorithme et s'arrêtent là. Or un algorithme peut sembler correct et contenir une erreur (condition inversée, mauvaise variable, cas limite oublié). \textbf{Règle d'or : tester toujours sur au moins deux jeux de valeurs}, dont un cas « limite » (par exemple exactement 110 kWh pour la facture, ou une moyenne de 10,00). Si la trace ne donne pas le résultat attendu, l'algorithme doit être corrigé.
\end{attention}

La justification écrite est ce qui distingue un élève qui a \emph{compris} d'un élève qui a recopié. Elle répond à trois questions : \textbf{quoi} (quelles entrées, quelles sorties), \textbf{comment} (quelles structures de contrôle), \textbf{pourquoi} (pourquoi cette structure et pas une autre). Exemple de rédaction complète :

\begin{exemplebox}[Modèle de justification rédigée]
« Le problème comporte une donnée d'entrée (la consommation $C$) et un résultat (le montant $M$). Le calcul dépend d'une condition : $C > 110$. J'ai donc utilisé une structure alternative \texttt{Si...Alors...Sinon}, qui permet d'exécuter un traitement différent selon que la condition est vraie ou fausse. Aucune boucle n'est nécessaire car aucune action n'est répétée. La trace avec $C = 150$ donne 11\,250 FCFA et celle avec $C = 80$ donne 4\,000 FCFA, ce qui confirme que l'algorithme est correct. »
\end{exemplebox}

\begin{savaistu}[Le conseil du correcteur]
Au BEPC, un algorithme clair, bien indenté, accompagné d'une trace d'exécution et d'une justification, vaut souvent \emph{plus de points} qu'un algorithme juste mais présenté sans aucune explication. Le correcteur évalue ta démarche autant que ton résultat : montre-lui que tu sais \emph{pourquoi} tu écris chaque ligne !
\end{savaistu}

\courssub{Synthèse de la leçon}

Toute la leçon tient en une idée : un algorithme est une suite finie d'instructions qui combine trois « briques » — la \cle{séquence}, l'\cle{alternative} et l'\cle{itérative} — et que l'on peut représenter par un algorigramme.

\begin{popfigure}
\begin{center}
\begin{tikzpicture}[
  font=\small,
  root/.style={rounded rectangle, draw=popdark, fill=popgold, text=white, minimum width=3.2cm, minimum height=0.9cm},
  branche/.style={rectangle, rounded corners, draw=popdark, fill=popblueL, minimum width=2.6cm, minimum height=0.8cm},
  feuille/.style={rectangle, rounded corners, draw=popmuted, fill=popcream, minimum width=2.4cm, minimum height=0.7cm, font=\footnotesize},
  lien/.style={thick, popmuted}
]
\node[root] (algo) {Algorithme};
\node[branche, above left=1.2cm and 0.6cm of algo] (seq) {Séquence};
\node[branche, left=1.6cm of algo] (alt) {Alternative};
\node[branche, below left=1.2cm and 0.6cm of algo] (iter) {Itérative};
\node[branche, right=1.6cm of algo] (gram) {Algorigramme};
\node[feuille, above=0.35cm of seq] (seqd) {instructions dans l'ordre};
\node[feuille, above=0.35cm of alt] (altd) {Si...Alors...Sinon};
\node[feuille, below=0.35cm of iter] (tq) {Tant que (nb. inconnu)};
\node[feuille, below=0.35cm of tq] (pour) {Pour (nb. connu)};
\node[feuille, above=0.35cm of gram] (gramd) {rectangles, losanges, flèches};
\node[feuille, below=0.35cm of gram] (gramv) {vérifié par la trace};
\draw[lien] (algo) -- (seq);
\draw[lien] (algo) -- (alt);
\draw[lien] (algo) -- (iter);
\draw[lien] (algo) -- (gram);
\draw[lien] (seq) -- (seqd);
\draw[lien] (alt) -- (altd);
\draw[lien] (iter) -- (tq);
\draw[lien] (tq) -- (pour);
\draw[lien] (gram) -- (gramd);
\draw[lien] (gram) -- (gramv);
\end{tikzpicture}
\end{center}
\legende{Carte mentale de synthèse : les structures de contrôle et leur représentation.}
\end{popfigure}

\begin{aretenir}[L'essentiel]
\begin{itemize}
\item Résoudre un problème suit 6 étapes : comprendre, identifier entrées/sorties, choisir les structures, écrire, tracer, conclure en justifiant.
\item \texttt{Si...Alors...Sinon} quand le traitement dépend d'une condition ; alternatives imbriquées pour plusieurs cas.
\item \texttt{Tant que} quand le nombre de répétitions est inconnu ; \texttt{Pour} quand il est connu.
\item On vérifie toujours un algorithme par une trace sur au moins deux jeux de valeurs, et on justifie ses choix par écrit.
\end{itemize}
\end{aretenir}

\begin{exercice}[À toi de jouer]
Un cybercafé facture la connexion 200 FCFA l'heure ; au-delà de 3 heures, l'heure passe à 150 FCFA pour \emph{toutes} les heures. Écris l'algorithme complet, trace-le pour 2 heures puis pour 5 heures, dessine son algorigramme et rédige la justification de tes choix de structures.
\end{exercice}

\begin{corrige}[Exercice 1 : Cybercafé]
\textbf{1. Algorithme complet}
\begin{lstlisting}
Algorithme FactureCyber
Variables
    duree, montant : Reel
Debut
    Ecrire("Entrez la duree de connexion (en heures) : ")
    Lire(duree)
    Si duree > 3 Alors
        montant <- duree * 150
    Sinon
        montant <- duree * 200
    FinSi
    Ecrire("Montant a payer : ", montant, " FCFA")
Fin
\end{lstlisting}

\textbf{2. Traces d'exécution}
\begin{itemize}
\item \textbf{Cas 1 : $duree = 2$}
\begin{center}
\begin{tabularx}{\linewidth}{|l|X|X|X|}
\hline
\rowcolor{popblueL} \textbf{Étape} & \textbf{duree} & \textbf{Test $duree>3$} & \textbf{montant} \\
\hline
Lecture & 2 & --- & --- \\
\hline
Test & 2 & Faux & --- \\
\hline
Calcul & 2 & --- & $2 \times 200 = 400$ \\
\hline
Affichage & \multicolumn{3}{X|}{Montant à payer : 400 FCFA} \\
\hline
\end{tabularx}
\end{center}

\item \textbf{Cas 2 : $duree = 5$}
\begin{center}
\begin{tabularx}{\linewidth}{|l|X|X|X|}
\hline
\rowcolor{popblueL} \textbf{Étape} & \textbf{duree} & \textbf{Test $duree>3$} & \textbf{montant} \\
\hline
Lecture & 5 & --- & --- \\
\hline
Test & 5 & Vrai & --- \\
\hline
Calcul & 5 & --- & $5 \times 150 = 750$ \\
\hline
Affichage & \multicolumn{3}{X|}{Montant à payer : 750 FCFA} \\
\hline
\end{tabularx}
\end{center}
\end{itemize}

\textbf{3. Algorigramme}
\begin{popfigure}
\begin{center}
\begin{tikzpicture}[
  font=\small,
  startstop/.style={rounded rectangle, draw=popdark, fill=popgoldL, minimum width=2.6cm, minimum height=0.8cm},
  io/.style={trapezium, trapezium left angle=70, trapezium right angle=110, draw=popdark, fill=popblueL, minimum height=0.8cm},
  process/.style={rectangle, draw=popdark, fill=popgreenL, minimum width=3.2cm, minimum height=0.8cm},
  decision/.style={diamond, draw=popdark, fill=poppinkL, aspect=2.4, inner sep=1pt},
  fleche/.style={-{Stealth[length=2.5mm]}, thick, popdark}
]
\node[startstop] (debut) {Début};
\node[io, below=0.5cm of debut] (lire) {Lire duree};
\node[decision, below=0.6cm of lire] (test) {duree $>$ 3 ?};
\node[process, below left=1.1cm and 0.4cm of test] (maj) {montant $\leftarrow$ duree $\times$ 150};
\node[process, below right=1.1cm and 0.4cm of test] (sim) {montant $\leftarrow$ duree $\times$ 200};
\node[io, below=2.6cm of test] (aff) {Afficher montant};
\node[startstop, below=0.5cm of aff] (fin) {Fin};
\draw[fleche] (debut) -- (lire);
\draw[fleche] (lire) -- (test);
\draw[fleche] (test.west) -| node[pos=0.25, above]{Oui} (maj.north);
\draw[fleche] (test.east) -| node[pos=0.25, above]{Non} (sim.north);
\draw[fleche] (maj.south) |- (aff.west);
\draw[fleche] (sim.south) |- (aff.east);
\draw[fleche] (aff) -- (fin);
\end{tikzpicture}
\end{center}
\legende{Algorigramme de la facturation du cybercafé.}
\end{popfigure}

\textbf{4. Justification}
« Le problème comporte une entrée (la durée de connexion) et une sortie (le montant). Le tarif change selon une condition unique : la durée dépasse-t-elle 3 heures ? J'ai donc utilisé une structure alternative \texttt{Si...Alors...Sinon}. Aucune boucle n'est nécessaire car le calcul est direct et non répétitif. Les traces pour 2 h (400 FCFA) et 5 h (750 FCFA) valident l'algorithme, y compris le cas limite de 3 h exactement (tarif normal : 600 FCFA). »
\end{corrige}

\newpage

\coursec{Activité d'intégration}

\coursec{Activité d'intégration : la caisse de la coopérative scolaire}

\courssub{Objectifs de l'activité}

À l'issue de cette activité d'intégration, tu seras capable de :
\begin{itemize}
\item analyser une situation réelle de gestion et identifier les \cle{données}, les \cle{traitements} et les \cle{résultats} attendus ;
\item choisir et \textbf{justifier} la structure de contrôle adaptée : \cle{structure alternative} (Si \ldots Alors \ldots Sinon) pour la remise, \cle{structure itérative} (TantQue) pour la répétition des clients ;
\item écrire un algorithme complet en \cle{pseudo-code} avec des variables correctement déclarées et initialisées ;
\item représenter cet algorithme par un \cle{algorigramme} respectant les symboles conventionnels ;
\item \textbf{tracer} (exécuter à la main) ton algorithme sur un jeu de données pour vérifier qu'il fonctionne, puis rédiger une conclusion argumentée.
\end{itemize}

\courssub{Situation}

\textbf{La caisse de la coopérative scolaire du Collège Bilingue de Nkongsamba.}

La coopérative scolaire du collège vend des fournitures aux élèves : cahiers, stylos, règles, tenues de classe. Chaque matin, Aminatou, élève en classe de 3\textsuperscript{ème} et trésorière de la coopérative, tient la caisse. Pour fidéliser les « gros » acheteurs (souvent les parents qui achètent pour plusieurs enfants), le bureau de la coopérative a décidé la règle suivante :

\begin{quote}
\emph{Tout client dont le montant total des achats dépasse strictement 5\,000 FCFA bénéficie d'une remise de 10 \% sur ce montant. Les autres clients paient le montant exact.}
\end{quote}

À la fermeture de la caisse, Aminatou doit annoncer au président de la coopérative :
\begin{itemize}
\item la \textbf{recette totale} de la matinée (somme des montants effectivement payés, remises déduites) ;
\item le \textbf{nombre de clients} servis.
\end{itemize}

Le problème, c'est qu'on ne sait \textbf{jamais à l'avance combien de clients} vont se présenter : un jour il y en a 3, un autre jour 27. Aminatou saisit donc les montants les uns après les autres, et elle signale la fin de la journée en tapant le montant \textbf{0} (un client ne peut pas avoir un montant d'achat nul : 0 est une valeur « sentinelle » qui veut dire « c'est fini »).

Ce matin-là, voici les montants des achats des clients qui se sont présentés, dans l'ordre :

\begin{center}
\begin{tabularx}{\linewidth}{|l|X|X|X|X|X|}
\hline
\rowcolor{popblueL} \textbf{Client} & n°1 & n°2 & n°3 & n°4 & n°5 \\
\hline
\textbf{Montant des achats (FCFA)} & 3\,500 & 8\,000 & 5\,000 & 12\,000 & 1\,200 \\
\hline
\end{tabularx}
\end{center}

Après le cinquième client, plus personne ne se présente : Aminatou saisit 0.

Fatiguée de tout calculer de tête (et de se tromper parfois !), Aminatou te demande de l'aider : \textbf{écris l'algorithme du caissier} qui, pour n'importe quelle journée, calcule et affiche la recette totale et le nombre de clients.

\courssub{Consignes}

Travail en groupes de 3 ou 4 élèves. Durée : une séance.

\begin{enumerate}
\item \textbf{Analyse du problème.} Identifie et liste : (a) les données que l'algorithme doit lire ; (b) les résultats qu'il doit afficher ; (c) les variables nécessaires (donne à chacune un nom et un rôle précis).
\item \textbf{Choix des structures.} Pour chacune des deux questions suivantes, choisis la structure adaptée et \textbf{justifie ton choix} en une phrase :
\begin{enumerate}
\item Faut-il une structure alternative ? Si oui, quelle est la condition testée et que fait-on dans chaque cas ?
\item Faut-il une boucle TantQue ou une boucle Pour ? Justifie précisément en t'appuyant sur le nombre de répétitions.
\end{enumerate}
\item \textbf{Pseudo-code.} Rédige l'algorithme complet en pseudo-code, avec la déclaration des variables, leur initialisation, la boucle, le test, et les affichages finaux.
\item \textbf{Algorigramme.} Représente ton algorithme par un algorigramme en respectant les symboles : ovale (Début/Fin), parallélogramme (lecture/affichage), rectangle (traitement), losange (test).
\item \textbf{Trace d'exécution.} Exécute ton algorithme « à la main » sur le jeu des 5 montants du tableau ci-dessus. Présente un tableau donnant, à chaque étape : le montant lu, la remise éventuelle, le montant payé, la recette cumulée et le nombre de clients. Vérifie que les résultats finaux sont cohérents.
\item \textbf{Justification orale.} Prépare une réponse argumentée (3 ou 4 phrases) à la question du professeur : \emph{« Pourquoi avez-vous choisi la boucle TantQue plutôt que la boucle Pour ? »}
\item \textbf{Pour aller plus loin (bonus).} Modifie ton algorithme pour qu'il affiche aussi, à la fin, le \textbf{montant moyen} payé par client. Attention au cas où aucun client ne s'est présenté !
\end{enumerate}

\courssub{Ressources à mobiliser}

\begin{aretenir}[Ce dont tu as besoin pour réussir]
\begin{itemize}
\item \textbf{Structure alternative} : \texttt{Si condition Alors traitement1 Sinon traitement2 FinSi}. La condition ici est une comparaison : \texttt{montant > 5000}.
\item \textbf{Calcul d'une remise de 10 \%} : montant payé $=$ montant $-$ montant $\times \dfrac{10}{100}$, c'est-à-dire montant payé $=$ montant $\times 0{,}9$.
\item \textbf{Boucle TantQue} : \texttt{TantQue condition Faire ... FinTantQue}. On l'utilise quand on \textbf{ne connaît pas à l'avance le nombre de répétitions} ; la condition est testée \emph{avant} chaque tour.
\item \textbf{Boucle Pour} : réservée au cas où le nombre de tours est \textbf{connu à l'avance}.
\item \textbf{Accumulateur} : pour cumuler une recette, on initialise une variable à 0 puis on ajoute : \texttt{recette $\leftarrow$ recette + montantPaye}.
\item \textbf{Compteur} : pour compter les clients : \texttt{nbClients $\leftarrow$ nbClients + 1}.
\item \textbf{Lecture avant la boucle} : avec une valeur sentinelle, il faut lire une première fois \emph{avant} le TantQue, puis relire \emph{à la fin} du corps de la boucle.
\item \textbf{Symboles de l'algorigramme} : ovale = début/fin ; parallélogramme = entrée/sortie ; rectangle = traitement ; losange = test (deux flèches sortantes : Oui/Non).
\end{itemize}
\end{aretenir}

\courssub{Démarche et corrigé commenté}

\begin{exoresolu}[La caisse de la coopérative scolaire]
Énoncé : écrire en pseudo-code puis en algorigramme l'algorithme du caissier de la coopérative (remise de 10 \% au-delà de 5\,000 FCFA, arrêt sur la saisie de 0), le tracer sur les montants 3\,500 ; 8\,000 ; 5\,000 ; 12\,000 ; 1\,200, et justifier le choix de la boucle.

\tcblower

\textbf{Étape 1 : analyse du problème.}

\begin{itemize}
\item \emph{Données lues} : les montants des achats, un par un, jusqu'à la saisie de 0.
\item \emph{Résultats affichés} : la recette totale et le nombre de clients.
\item \emph{Variables} :
\begin{itemize}
\item \texttt{montant} : montant des achats du client en cours (réel) ;
\item \texttt{montantPaye} : montant effectivement payé après remise éventuelle (réel) ;
\item \texttt{recette} : accumulateur de la recette totale, initialisé à 0 (réel) ;
\item \texttt{nbClients} : compteur de clients, initialisé à 0 (entier).
\end{itemize}
\end{itemize}

\textbf{Étape 2 : choix et justification des structures.}

\begin{enumerate}
\item \emph{Structure alternative} : oui, il en faut une, car le traitement dépend du montant. La condition est \texttt{montant > 5000} : si elle est vraie, on applique la remise de 10 \% ; sinon, le client paie le montant exact. Attention : « dépasse 5\,000 » signifie \emph{strictement supérieur} ; un montant de exactement 5\,000 FCFA ne donne \textbf{pas} droit à la remise.
\item \emph{Boucle} : on choisit \textbf{TantQue}, car le nombre de clients n'est \textbf{pas connu à l'avance} : la répétition dépend d'une condition (\texttt{montant <> 0}) et non d'un compteur de tours fixé au départ. La boucle Pour est impossible ici, puisqu'elle exige de connaître le nombre de répétitions avant de commencer.
\end{enumerate}

\textbf{Étape 3 : le pseudo-code.}

\begin{lstlisting}
Algorithme CaisseCooperative
Variables
    montant, montantPaye, recette : Reel
    nbClients : Entier
Debut
    recette <- 0
    nbClients <- 0
    Ecrire("Entrez le montant des achats (0 pour terminer) : ")
    Lire(montant)
    TantQue montant <> 0 Faire
        Si montant > 5000 Alors
            montantPaye <- montant * 0.9   // remise de 10 %
        Sinon
            montantPaye <- montant         // pas de remise
        FinSi
        recette <- recette + montantPaye
        nbClients <- nbClients + 1
        Ecrire("Entrez le montant des achats (0 pour terminer) : ")
        Lire(montant)
    FinTantQue
    Ecrire("Recette totale : ", recette, " FCFA")
    Ecrire("Nombre de clients : ", nbClients)
Fin
\end{lstlisting}

Deux points délicats à remarquer : la \textbf{première lecture} a lieu \emph{avant} la boucle (sinon la condition du TantQue porterait sur une variable vide), et la \textbf{lecture suivante} est placée \emph{à la fin} du corps de la boucle (sinon on bouclerait indéfiniment sur le même montant).

\textbf{Étape 4 : l'algorigramme.}

\begin{popfigure}
\begin{center}
\begin{tikzpicture}[
  font=\small,
  startstop/.style={ellipse, draw=popdark, fill=poporangeL, minimum width=2.6cm, minimum height=0.8cm},
  io/.style={trapezium, trapezium left angle=70, trapezium right angle=110, draw=popdark, fill=popblueL, minimum height=0.8cm},
  process/.style={rectangle, draw=popdark, fill=popgreenL, minimum width=4.2cm, minimum height=0.8cm},
  decision/.style={diamond, draw=popdark, fill=poppinkL, aspect=2.4, inner sep=1pt},
  arrow/.style={-{Stealth}, thick, popdark},
  node distance=0.55cm and 1.1cm]
\node[startstop] (debut) {Début};
\node[process, below=of debut] (init) {recette $\leftarrow$ 0 ; nbClients $\leftarrow$ 0};
\node[io, below=of init] (lire1) {Lire montant};
\node[decision, below=of lire1] (test0) {montant $\neq$ 0 ?};
\node[decision, below=1.1cm of test0] (test1) {montant $>$ 5000 ?};
\node[process, right=1.6cm of test1] (remise) {montantPaye $\leftarrow$ montant $\times$ 0,9};
\node[process, left=1.6cm of test1] (pasremise) {montantPaye $\leftarrow$ montant};
\node[process, below=2.4cm of test1] (cumul) {recette $\leftarrow$ recette + montantPaye};
\node[process, below=of cumul] (compt) {nbClients $\leftarrow$ nbClients + 1};
\node[io, below=of compt] (lire2) {Lire montant};
\node[io, right=3.4cm of test0] (affich) {Afficher recette, nbClients};
\node[startstop, below=of affich] (fin) {Fin};
\draw[arrow] (debut) -- (init);
\draw[arrow] (init) -- (lire1);
\draw[arrow] (lire1) -- (test0);
\draw[arrow] (test0) -- node[right]{Oui} (test1);
\draw[arrow] (test0) -- node[above]{Non} (affich);
\draw[arrow] (test1) -- node[above]{Oui} (remise);
\draw[arrow] (test1) -- node[above]{Non} (pasremise);
\draw[arrow] (remise) |- (cumul);
\draw[arrow] (pasremise) |- (cumul);
\draw[arrow] (cumul) -- (compt);
\draw[arrow] (compt) -- (lire2);
\draw[arrow] (lire2.west) -- ++(-1.4,0) |- (test0);
\draw[arrow] (affich) -- (fin);
\end{tikzpicture}
\end{center}
\legende{Algorigramme de l'algorithme du caissier : la boucle TantQue apparaît comme un losange dont la branche « Oui » revient en arrière.}
\end{popfigure}

\textbf{Étape 5 : trace d'exécution sur le jeu de montants.}

\begin{center}
\begin{tabularx}{\linewidth}{|l|X|X|X|X|X|}
\hline
\rowcolor{popblueL} \textbf{Étape} & \textbf{Montant lu} & \textbf{Test $> 5\,000$ ?} & \textbf{Montant payé} & \textbf{Recette cumulée} & \textbf{nbClients} \\
\hline
Initialisation & --- & --- & --- & 0 & 0 \\
\hline
Client 1 & 3\,500 & Non & 3\,500 & 3\,500 & 1 \\
\hline
Client 2 & 8\,000 & Oui & 7\,200 & 10\,700 & 2 \\
\hline
Client 3 & 5\,000 & Non & 5\,000 & 15\,700 & 3 \\
\hline
Client 4 & 12\,000 & Oui & 10\,800 & 26\,500 & 4 \\
\hline
Client 5 & 1\,200 & Non & 1\,200 & 27\,700 & 5 \\
\hline
Fin & 0 & \multicolumn{4}{l|}{Sortie de boucle : affichage des résultats} \\
\hline
\end{tabularx}
\end{center}

Détail des remises : $8\,000 \times 0{,}9 = 7\,200$ FCFA et $12\,000 \times 0{,}9 = 10\,800$ FCFA. Le client 3, avec exactement 5\,000 FCFA, ne \emph{dépasse} pas 5\,000 : il paie donc 5\,000 FCFA, sans remise. L'algorithme affiche finalement :
\begin{center}
\textbf{Recette totale : 27\,700 FCFA \quad --- \quad Nombre de clients : 5}
\end{center}

\textbf{Étape 6 : justification orale du choix de la boucle.}

\emph{« Nous avons choisi la boucle TantQue parce que le nombre de clients n'est pas connu à l'avance : il varie chaque jour. La boucle Pour exige de fixer le nombre de tours avant de commencer, ce qui est impossible ici. Avec TantQue, la répétition est contrôlée par la condition \texttt{montant <> 0} : dès qu'Aminatou saisit la valeur sentinelle 0, la condition devient fausse et la boucle s'arrête. C'est donc la structure adaptée à une répétition dont le nombre de tours dépend d'un événement. »}

\textbf{Étape 7 (bonus) : montant moyen par client.}

Il suffit d'ajouter, après la boucle, un test pour éviter la division par zéro :
\begin{lstlisting}
Si nbClients > 0 Alors
    Ecrire("Montant moyen : ", recette / nbClients, " FCFA")
Sinon
    Ecrire("Aucun client ce jour.")
FinSi
\end{lstlisting}
Sur notre exemple : $27\,700 \div 5 = 5\,540$ FCFA en moyenne par client.
\end{exoresolu}

\begin{attention}[Erreurs fréquentes à éviter]
\begin{itemize}
\item \textbf{Oublier la lecture avant la boucle} : la condition du TantQue porterait sur une variable non initialisée.
\item \textbf{Oublier la lecture à la fin du corps de boucle} : le programme traite le même montant indéfiniment (boucle infinie).
\item \textbf{Écrire \texttt{montant >= 5000}} au lieu de \texttt{montant > 5000} : le client à exactement 5\,000 FCFA recevrait la remise à tort. Lis bien l'énoncé : « \emph{dépasse} » signifie strictement supérieur.
\item \textbf{Cumuler \texttt{montant} au lieu de \texttt{montantPaye}} : la recette serait fausse dès qu'une remise s'applique.
\item \textbf{Compter la sentinelle} : le 0 final ne doit ni être ajouté à la recette ni être compté comme un client ; c'est bien le test du TantQue qui l'écarte.
\end{itemize}
\end{attention}

\courssub{Bilan de l'activité}

Cette activité t'a permis de mettre en évidence plusieurs idées essentielles de la leçon :

\begin{itemize}
\item \textbf{Le choix d'une structure de contrôle se justifie} : on ne choisit pas une boucle « au hasard ». La boucle \cle{TantQue} s'impose quand le nombre de répétitions est inconnu et dépend d'une condition (ici, la saisie de la valeur sentinelle 0) ; la boucle \cle{Pour} est réservée aux répétitions dont le nombre de tours est connu à l'avance ; la structure \cle{alternative} traduit une décision (accorder ou non la remise).
\item \textbf{Un algorithme de gestion réel} repose souvent sur le trio \emph{lecture -- test -- cumul} : lire une donnée, décider du traitement, accumuler dans un \cle{accumulateur} et compter avec un \cle{compteur}.
\item \textbf{L'algorigramme} est un outil de communication : il rend la logique de l'algorithme visible (la boucle apparaît comme une flèche qui « remonte ») et permet de vérifier la démarche avant de la rédiger en pseudo-code.
\item \textbf{La trace d'exécution} est une méthode de vérification rigoureuse : en exécutant l'algorithme à la main sur un jeu de données connu, on détecte les erreurs (remise mal calculée, sentinelle comptée, boucle infinie) avant toute utilisation réelle.
\item Enfin, tu as rédigé et justifié une démarche complète, de l'analyse du problème jusqu'à la conclusion argumentée : c'est exactement la compétence attendue au BEPC.
\end{itemize}

\newpage

\coursec{Méthodes et exercices résolus}

\coursec{Utilisation des structures de contrôle}

\courssub{Méthode 1 : Écrire un algorithme simple en pseudo-code}

\begin{methode}[Rédiger un algorithme proprement]
Pour écrire un algorithme simple qui résout un problème de la vie courante :
\begin{enumerate}
\item \textbf{Lire l'énoncé} et identifier ce que l'on \emph{connaît} (les \cle{données}) et ce que l'on \emph{cherche} (le \cle{résultat}).
\item \textbf{Déclarer les variables} : donner un nom explicite et un \cle{type} (\texttt{Entier}, \texttt{Reel}, \cle{Chaîne de caractères}) à chaque donnée et à chaque résultat.
\item \textbf{Écrire l'en-tête} : le mot \texttt{Algorithme} suivi d'un nom significatif, puis la liste des variables.
\item \textbf{Saisir les données} avec l'instruction \texttt{Lire} (ou \texttt{Saisir}).
\item \textbf{Décrire les traitements} dans l'ordre logique, avec des \cle{affectations} du type \texttt{variable $\leftarrow$ expression}.
\item \textbf{Afficher le résultat} avec l'instruction \texttt{Écrire} (ou \texttt{Afficher}).
\item \textbf{Terminer} par le mot \texttt{Fin} et \textbf{vérifier} l'algorithme en le testant « à la main » avec des valeurs simples (trace d'exécution).
\end{enumerate}
\textbf{Réflexes à retenir :} une variable doit toujours être déclarée avant d'être utilisée ; l'affectation se lit de droite à gauche (\texttt{a $\leftarrow$ 5} signifie « a reçoit la valeur 5 ») ; chaque instruction occupe une ligne ; on écrit les mots-clés en français (\texttt{Debut}, \texttt{Fin}, \texttt{Si}, \texttt{Alors}, \texttt{TantQue}, \texttt{Pour}).
\end{methode}

\begin{exoresolu}[Le taxi de Ngoa-Ekellé]
\textbf{Énoncé.} Un taximan de Yaoundé facture une course de la façon suivante : une prise en charge fixe de 200 FCFA, plus 150 FCFA par kilomètre parcouru. Écrire un algorithme qui demande le nombre de kilomètres parcourus et affiche le prix total de la course.
\tcblower
\textbf{Solution détaillée.}

\textbf{Étape 1 — Analyse.} Donnée connue : le nombre de kilomètres, noté \texttt{km} (réel). Résultat cherché : le prix total, noté \texttt{prix} (réel). La formule est :
\[ \text{prix} = 200 + 150 \times \text{km} \]

\textbf{Étape 2 — Rédaction de l'algorithme.}
\begin{lstlisting}
Algorithme PrixCourse
Variables
    km, prix : Reel
Debut
    Ecrire("Entrez le nombre de kilometres : ")
    Lire(km)
    prix <- 200 + 150 * km
    Ecrire("Le prix de la course est : ", prix, " FCFA")
Fin
\end{lstlisting}

\textbf{Étape 3 — Vérification.} Testons avec $\text{km} = 4$ : $\text{prix} = 200 + 150 \times 4 = 200 + 600 = 800$ FCFA. L'algorithme affiche bien 800 FCFA, ce qui correspond au calcul manuel.

\textbf{Conclusion.} L'algorithme \texttt{PrixCourse} saisit la distance, applique la formule $\text{prix} = 200 + 150\,\text{km}$ et affiche le montant à payer : il résout correctement le problème posé.
\end{exoresolu}

\courssub{Méthode 2 : Utiliser une structure alternative (Si ... Alors ... Sinon)}

\begin{methode}[Choisir et rédiger un test]
Quand le traitement dépend d'une \cle{condition} :
\begin{enumerate}
\item \textbf{Identifier la condition} : c'est une question dont la réponse est \texttt{VRAI} ou \texttt{FAUX} (comparaison avec $=$, $<$, $>$, $\leq$, $\geq$, $\neq$). On peut combiner des conditions avec \texttt{Et}, \texttt{Ou}, \texttt{Non}.
\item \textbf{Choisir la forme adaptée} :
\begin{itemize}
\item \emph{Alternative simple} (un seul cas à traiter) : \texttt{Si condition Alors ... FinSi}
\item \emph{Alternative complète} (deux cas) : \texttt{Si condition Alors ... Sinon ... FinSi}
\item \emph{Cas multiples} : enchaîner avec \texttt{SinonSi}.
\end{itemize}
\item \textbf{Écrire les instructions} de chaque branche, en les décalant (indentation) pour la lisibilité.
\item \textbf{Fermer le test} avec \texttt{FinSi} : oublier ce mot est une erreur classique.
\item \textbf{Tester les deux branches} avec des valeurs qui rendent la condition vraie, puis fausse (dont les valeurs limites).
\end{enumerate}
\textbf{Attention aux bornes :} « au moins 10 » se traduit par \texttt{note >= 10} et non \texttt{note > 10} ; « au plus 10 » par \texttt{note <= 10}.
\end{methode}

\begin{exoresolu}[Décision du conseil de classe]
\textbf{Énoncé.} Dans un collège de Douala, un élève de 3ème est déclaré admis en classe supérieure si sa moyenne annuelle est supérieure ou égale à 10/20. Écrire un algorithme qui demande la moyenne d'un élève et affiche « Admis » ou « Ajourné ».
\tcblower
\textbf{Solution détaillée.}

\textbf{Étape 1 — Analyse.} Donnée : la moyenne \texttt{m} (réel, comprise entre 0 et 20). Condition : $\texttt{m} \geq 10$. Deux cas possibles, donc alternative \emph{complète}.

\textbf{Étape 2 — Rédaction.}
\begin{lstlisting}
Algorithme DecisionConseil
Variables
    m : Reel
Debut
    Ecrire("Entrez la moyenne annuelle : ")
    Lire(m)
    Si m >= 10 Alors
        Ecrire("Admis en classe superieure")
    Sinon
        Ecrire("Ajourne")
    FinSi
Fin
\end{lstlisting}

\textbf{Étape 3 — Vérification des deux branches.}
\begin{itemize}
\item Avec $\texttt{m} = 12,5$ : la condition $12,5 \geq 10$ est VRAIE, l'algorithme affiche « Admis ». Correct.
\item Avec $\texttt{m} = 10$ : la condition $10 \geq 10$ est VRAIE, l'algorithme affiche « Admis ». C'est bien le cas limite demandé (« supérieure \emph{ou égale} à 10 »).
\item Avec $\texttt{m} = 7$ : la condition est FAUSSE, l'algorithme affiche « Ajourné ». Correct.
\end{itemize}

\textbf{Conclusion.} L'alternative complète \texttt{Si ... Alors ... Sinon ... FinSi} permet de traiter les deux décisions possibles ; le test sur la valeur limite 10 confirme que l'opérateur $\geq$ était le bon choix.
\end{exoresolu}

\courssub{Méthode 3 : Utiliser la boucle « Tant que »}

\begin{methode}[Construire une boucle Tant que]
On utilise \texttt{TantQue} quand on \textbf{ne connaît pas à l'avance} le nombre de répétitions :
\begin{enumerate}
\item \textbf{Identifier la condition de continuation} : on répète \emph{tant que} la condition est vraie ; on s'arrête dès qu'elle devient fausse.
\item \textbf{Initialiser} avant la boucle les variables utilisées dans la condition (sinon le premier test est impossible).
\item \textbf{Écrire le corps de la boucle} : les instructions à répéter.
\item \textbf{Prévoir la modification} : dans le corps, au moins une instruction doit faire évoluer la condition (par exemple \texttt{Lire} ou une incrémentation), sinon la boucle est \emph{infinie}.
\item \textbf{Fermer} avec \texttt{FinTantQue} et \textbf{tracer l'exécution} dans un tableau (une colonne par variable, une ligne par tour de boucle).
\end{enumerate}
\textbf{Point de vigilance :} si la condition est fausse dès le départ, le corps de la boucle n'est \emph{jamais} exécuté : c'est normal et parfois souhaité.
\end{methode}

\begin{exoresolu}[L'épargne de la tontine]
\textbf{Énoncé.} Amina dépose chaque mois 5\,000 FCFA sur son compte épargne, qui contient initialement 12\,000 FCFA. Elle veut savoir au bout de combien de mois son épargne atteindra ou dépassera 50\,000 FCFA. Écrire un algorithme qui calcule et affiche ce nombre de mois.
\tcblower
\textbf{Solution détaillée.}

\textbf{Étape 1 — Analyse.} On ne connaît pas à l'avance le nombre de dépôts : c'est une boucle \texttt{TantQue}. Variables : \texttt{solde} (réel) initialisé à 12\,000, \texttt{mois} (entier) initialisé à 0. Condition de continuation : $\texttt{solde} < 50000$ (on continue tant que l'objectif n'est pas atteint).

\textbf{Étape 2 — Rédaction.}
\begin{lstlisting}
Algorithme EpargneTontine
Variables
    solde : Reel
    mois : Entier
Debut
    solde <- 12000
    mois <- 0
    TantQue solde < 50000 Faire
        solde <- solde + 5000
        mois <- mois + 1
    FinTantQue
    Ecrire("Objectif atteint au bout de ", mois, " mois")
Fin
\end{lstlisting}

\textbf{Étape 3 — Trace d'exécution.}
\begin{center}
\begin{tabularx}{\linewidth}{|c|X|X|l|}
\hline
\rowcolor{popblueL} Tour & solde (avant test) & mois & Condition $\texttt{solde}<50000$ \\
\hline
0 & 12\,000 & 0 & VRAI \\
\hline
1 & 17\,000 & 1 & VRAI \\
\hline
2 & 22\,000 & 2 & VRAI \\
\hline
$\cdots$ & $\cdots$ & $\cdots$ & VRAI \\
\hline
7 & 47\,000 & 7 & VRAI \\
\hline
8 & 52\,000 & 8 & FAUX : on sort \\
\hline
\end{tabularx}
\end{center}
Le solde augmente de 5\,000 par tour ; après 8 tours, $\text{solde} = 12\,000 + 8 \times 5\,000 = 52\,000 \geq 50\,000$ : la boucle s'arrête.

\textbf{Conclusion.} L'algorithme affiche « Objectif atteint au bout de 8 mois ». La boucle \texttt{TantQue} était bien adaptée car le nombre de répétitions était inconnu au départ.
\end{exoresolu}

\courssub{Méthode 4 : Utiliser la boucle « Pour »}

\begin{methode}[Construire une boucle Pour]
On utilise \texttt{Pour} quand le \textbf{nombre de répétitions est connu à l'avance} :
\begin{enumerate}
\item \textbf{Déterminer le compteur} : une variable entière (souvent \texttt{i}) qui va parcourir les valeurs.
\item \textbf{Déterminer les bornes} : valeur de début et valeur de fin (la borne de fin est \emph{incluse}).
\item \textbf{Écrire la syntaxe} : \texttt{Pour i de 1 a N Faire ... FinPour}. Le compteur augmente automatiquement de 1 à chaque tour : il ne faut \textbf{pas} l'incrémenter soi-même.
\item \textbf{Initialiser avant la boucle} les variables d'accumulation (somme $\leftarrow$ 0, produit $\leftarrow$ 1).
\item \textbf{Vérifier le nombre de tours} : de $a$ à $b$, il y a exactement $b - a + 1$ répétitions.
\end{enumerate}
\textbf{Réflexe :} pour calculer une somme $S = 1 + 2 + \cdots + N$, on écrit dans la boucle \texttt{S <- S + i}.
\end{methode}

\begin{exoresolu}[Total des ventes de la semaine]
\textbf{Énoncé.} La boutique de la maman de Junior est ouverte 6 jours par semaine. Chaque jour, on note la recette du jour. Écrire un algorithme qui demande les 6 recettes, puis calcule et affiche la recette totale de la semaine.
\tcblower
\textbf{Solution détaillée.}

\textbf{Étape 1 — Analyse.} Le nombre de saisies est connu : 6 jours, donc 6 répétitions $\Rightarrow$ boucle \texttt{Pour}. Variables : \texttt{recette} (réel) pour chaque saisie, \texttt{total} (réel) accumulateur initialisé à 0, \texttt{i} (entier) compteur.

\textbf{Étape 2 — Rédaction.}
\begin{lstlisting}
Algorithme RecetteSemaine
Variables
    recette, total : Reel
    i : Entier
Debut
    total <- 0
    Pour i de 1 a 6 Faire
        Ecrire("Recette du jour ", i, " : ")
        Lire(recette)
        total <- total + recette
    FinPour
    Ecrire("Recette totale de la semaine : ", total, " FCFA")
Fin
\end{lstlisting}

\textbf{Étape 3 — Vérification.} Nombre de tours : $6 - 1 + 1 = 6$, donc 6 saisies, comme prévu. Testons avec six recettes de 10\,000 FCFA chacune : après chaque tour, \texttt{total} vaut successivement 10\,000 ; 20\,000 ; 30\,000 ; 40\,000 ; 50\,000 ; 60\,000. L'algorithme affiche 60\,000 FCFA, ce qui est exact : $6 \times 10\,000 = 60\,000$.

\textbf{Conclusion.} La boucle \texttt{Pour} convient parfaitement car le nombre de jours (6) est fixe ; l'accumulateur \texttt{total}, initialisé à 0 avant la boucle, contient bien la somme des 6 recettes à la sortie.
\end{exoresolu}

\courssub{Méthode 5 : Représenter un algorithme par un algorigramme}

\begin{methode}[Dessiner un algorigramme correct]
Pour traduire un algorithme en schéma graphique normé :
\begin{enumerate}
\item \textbf{Respecter les symboles normalisés} :
\begin{itemize}
\item \emph{Ovale} : Début / Fin ;
\item \emph{Parallélogramme} : entrée (Lire) ou sortie (Écrire) ;
\item \emph{Rectangle} : traitement (affectation, calcul) ;
\item \emph{Losange} : test (condition), avec deux flèches sortantes étiquetées \emph{Oui} / \emph{Non}.
\end{itemize}
\item \textbf{Commencer} par un ovale « Début » et \textbf{terminer} par un ovale « Fin ».
\item \textbf{Relier les symboles par des flèches} orientées de haut en bas, dans l'ordre d'exécution.
\item \textbf{Pour une alternative} : le losange a deux sorties ; les deux branches se rejoignent (ou vont vers Fin) après leurs instructions.
\item \textbf{Pour une boucle} : la flèche « Oui » du losange descend vers le corps de la boucle, puis une flèche \emph{remonte} vers le losange ; la flèche « Non » continue vers la suite.
\item \textbf{Vérifier} que chaque losange a exactement deux sorties et que tout chemin mène à « Fin ».
\end{enumerate}
\end{methode}

\begin{exoresolu}[Algorigramme du plus grand de deux nombres]
\textbf{Énoncé.} Écrire l'algorithme qui lit deux nombres et affiche le plus grand, puis le représenter par un algorigramme.
\tcblower
\textbf{Solution détaillée.}

\textbf{Étape 1 — Algorithme.}
\begin{lstlisting}
Algorithme PlusGrand
Variables
    a, b : Reel
Debut
    Lire(a)
    Lire(b)
    Si a > b Alors
        Ecrire("Le plus grand est : ", a)
    Sinon
        Ecrire("Le plus grand est : ", b)
    FinSi
Fin
\end{lstlisting}

\textbf{Étape 2 — Algorigramme.}
\begin{popfigure}
\begin{center}
\begin{tikzpicture}[
  node distance=1.2cm,
  startstop/.style={ellipse, draw=popdark, fill=poporangeL, minimum width=2.5cm, minimum height=1cm, font=\small},
  io/.style={trapezium, trapezium left angle=70, trapezium right angle=110, draw=popdark, fill=popblueL, minimum width=3.5cm, minimum height=1cm, font=\small, trapezium stretches body},
  decision/.style={diamond, draw=popdark, fill=poppinkL, aspect=2.5, inner sep=2pt, font=\small},
  process/.style={rectangle, draw=popdark, fill=popgreenL, minimum width=3.5cm, minimum height=1cm, font=\small},
  fleche/.style={-{Stealth}, thick, popdark}]
\node[startstop] (debut) {Début};
\node[io, below=of debut] (lire) {Lire a, b};
\node[decision, below=of lire] (test) {$a > b$ ?};
\node[process, below left=1.2cm and 1.8cm of test] (ea) {Écrire a};
\node[process, below right=1.2cm and 1.8cm of test] (eb) {Écrire b};
\node[startstop, below=3.5cm of test] (fin) {Fin};
\draw[fleche] (debut) -- (lire);
\draw[fleche] (lire) -- (test);
\draw[fleche] (test) -- node[above left, pos=0.5]{Oui} (ea);
\draw[fleche] (test) -- node[above right, pos=0.5]{Non} (eb);
\draw[fleche] (ea) |- (fin);
\draw[fleche] (eb) |- (fin);
\end{tikzpicture}
\end{center}
\legende{Algorigramme de l'algorithme \texttt{PlusGrand} : ovale (début/fin), parallélogramme (lecture), losange (test à deux sorties), rectangles (affichages).}
\end{popfigure}

\textbf{Étape 3 — Vérification.} Le losange possède bien deux sorties étiquetées « Oui » et « Non » ; les deux chemins se rejoignent avant « Fin ». Test : avec $a = 7$ et $b = 3$, la condition $7 > 3$ est vraie, on affiche 7 ; avec $a = 2$ et $b = 9$, elle est fausse, on affiche 9. Les deux branches fonctionnent.

\textbf{Conclusion.} L'algorigramme traduit fidèlement l'algorithme : chaque symbole normalisé correspond à une instruction, et les flèches décrivent l'ordre d'exécution.
\end{exoresolu}

\courssub{Méthode 6 : Résoudre un problème concret et justifier sa démarche}

\begin{methode}[Rédiger et justifier une solution complète]
Face à un problème de la vie courante à résoudre par un algorithme (situation type BEPC) :
\begin{enumerate}
\item \textbf{Dégager les données et les résultats} dans une petite phrase d'analyse.
\item \textbf{Choisir les structures} et le \emph{justifier} : « j'utilise un \texttt{Si} car il y a un choix », « j'utilise un \texttt{Pour} car le nombre de répétitions est connu », « j'utilise un \texttt{TantQue} car on répète jusqu'à atteindre un objectif ».
\item \textbf{Rédiger l'algorithme} complet : en-tête, déclarations, \texttt{Debut} ... \texttt{Fin}.
\item \textbf{Valider par une trace d'exécution} : tableau donnant la valeur de chaque variable à chaque étape, avec des valeurs numériques de l'énoncé.
\item \textbf{Conclure par une phrase} qui répond exactement à la question posée, avec l'unité (FCFA, mois, élèves...).
\end{enumerate}
\textbf{Barème implicite :} la justification du choix des structures et la trace d'exécution rapportent souvent autant de points que l'algorithme lui-même.
\end{methode}

\begin{exoresolu}[La cantine du collège]
\textbf{Énoncé.} Le cantinier d'un collège de Bafoussam vend le ticket de repas à 500 FCFA. Un élève ne peut acheter un ticket que s'il présente sa carte et si le stock n'est pas épuisé. Le stock initial est de 100 tickets. Écrire un algorithme qui, pour chacun des élèves d'une file de $N$ élèves, demande s'il a sa carte (1 = oui, 0 = non), vend le ticket si possible, et affiche à la fin la recette totale et le nombre de tickets restants.
\tcblower
\textbf{Solution détaillée.}

\textbf{Étape 1 — Analyse.} Données : $N$ (nombre d'élèves, entier), la réponse \texttt{carte} (0 ou 1) pour chaque élève. Résultats : \texttt{recette} et \texttt{stock}. 

\textbf{Étape 2 — Choix justifié des structures.} Le nombre d'élèves $N$ est connu à l'avance : boucle \texttt{Pour}. À chaque élève, il y a une décision (vendre ou non) : alternative \texttt{Si}. La double condition « carte présentée \textbf{et} stock disponible » se traduit par \texttt{carte = 1 Et stock > 0}.

\textbf{Étape 3 — Algorithme.}
\begin{lstlisting}
Algorithme Cantine
Variables
    N, i, carte, stock : Entier
    recette : Reel
Debut
    stock <- 100
    recette <- 0
    Ecrire("Nombre d'eleves dans la file : ")
    Lire(N)
    Pour i de 1 a N Faire
        Ecrire("L'eleve ", i, " a-t-il sa carte ? (1=oui, 0=non) ")
        Lire(carte)
        Si carte = 1 Et stock > 0 Alors
            stock <- stock - 1
            recette <- recette + 500
            Ecrire("Ticket vendu")
        Sinon
            Ecrire("Vente refusee")
        FinSi
    FinPour
    Ecrire("Recette totale : ", recette, " FCFA")
    Ecrire("Tickets restants : ", stock)
Fin
\end{lstlisting}

\textbf{Étape 4 — Trace d'exécution} avec $N = 3$ et les réponses 1, 0, 1 (stock initial 100) :
\begin{center}
\begin{tabularx}{\linewidth}{|c|c|X|X|l|}
\hline
\rowcolor{popblueL} Tour $i$ & carte & stock & recette & Action \\
\hline
1 & 1 & 99 & 500 & Vente acceptée ($1=1$ et $100>0$) \\
\hline
2 & 0 & 99 & 500 & Vente refusée (pas de carte) \\
\hline
3 & 1 & 98 & 1\,000 & Vente acceptée \\
\hline
\end{tabularx}
\end{center}

\textbf{Étape 5 — Conclusion.} Pour cette file de 3 élèves, l'algorithme affiche une recette de 1\,000 FCFA et 98 tickets restants, ce qui correspond exactement aux deux ventes réalisées. La démarche est justifiée : boucle \texttt{Pour} car $N$ est connu, alternative avec condition double car la vente exige à la fois la carte et un stock non nul.
\end{exoresolu}

\begin{aretenir}[Synthèse : Les trois structures de contrôle]
\begin{itemize}
\item \textbf{Structure séquentielle} : instructions exécutées les unes après les autres (par défaut).
\item \textbf{Structure alternative (Choix)} : \texttt{Si condition Alors ... Sinon ... FinSi}. Exécute un bloc \textbf{ou} l'autre selon une condition.
\item \textbf{Structure itérative (Boucle)} :
\begin{itemize}
\item \texttt{Pour i de a à b Faire ... FinPour} : nombre de tours \textbf{connu} ($b-a+1$ tours). Compteur auto-incrémenté.
\item \texttt{TantQue condition Faire ... FinTantQue} : nombre de tours \textbf{inconnu}. Condition testée \textbf{avant} chaque tour. Initialisation obligatoire avant la boucle.
\end{itemize}
\end{itemize}
\textbf{Règle d'or :} Toute variable utilisée (compteur, accumulateur, condition) doit être \textbf{initialisée} avant d'entrer dans la structure.
\end{aretenir}

\begin{attention}[Les 5 erreurs qui coûtent des points à l'examen]
\begin{enumerate}
\item \textbf{Confondre « $>$ » et « $\geq$ ».} « Au moins 10/20 » se traduit par \texttt{m >= 10}. Avec \texttt{m > 10}, un élève ayant exactement 10 est rejeté : la condition d'application est fausse et la réponse au cas limite est incorrecte.
\item \textbf{Oublier d'initialiser les variables avant la boucle.} Un accumulateur non initialisé (\texttt{total} ou \texttt{solde}) a une valeur inconnue : tous les résultats suivants sont faux. Toujours écrire \texttt{total <- 0} (ou la valeur initiale de l'énoncé) \emph{avant} la boucle.
\item \textbf{Créer une boucle infinie avec « TantQue ».} Si aucune instruction du corps ne fait évoluer la condition (oubli de \texttt{mois <- mois + 1}), la boucle ne s'arrête jamais. Vérifier toujours que la condition finira par devenir fausse.
\item \textbf{Incrémenter soi-même le compteur d'une boucle « Pour ».} Dans \texttt{Pour i de 1 a N}, le compteur augmente automatiquement ; ajouter \texttt{i <- i + 1} dans le corps fausse le nombre de tours (il y a $b - a + 1$ tours, pas plus).
\item \textbf{Mal dessiner l'algorigramme.} Erreurs fréquentes : utiliser un rectangle pour un test (il faut un \emph{losange}), oublier les étiquettes « Oui/Non » sur les deux flèches sortantes, ou faire remonter la flèche d'une boucle au mauvais endroit. Chaque losange doit avoir exactement deux sorties et tout chemin doit aboutir à « Fin ».
\end{enumerate}
\end{attention}

\begin{exercice}[Entraînement 1 : Pair ou Impair]
Écrire un algorithme qui demande un entier positif à l'utilisateur et affiche « Pair » s'il est divisible par 2, « Impair » sinon. Représenter l'algorithme par un algorigramme.
\end{exercice}

\begin{exercice}[Entraînement 2 : Table de multiplication]
Écrire un algorithme qui demande un entier $n$ et affiche sa table de multiplication de 1 à 10 (ex: « 1 x 7 = 7 », « 2 x 7 = 14 », ...). Justifier le choix de la boucle.
\end{exercice}

\begin{exercice}[Entraînement 3 : Saisie sécurisée]
Écrire un algorithme qui demande à l'utilisateur de saisir une note comprise entre 0 et 20. Tant que la note n'est pas dans cet intervalle, il redemande la saisie. Une fois la note valide, il l'affiche. Justifier le choix de la boucle.
\end{exercice}

\newpage

\coursec{Exercices}

\courssub{Je vérifie mes connaissances}

\begin{exercice}[QCM : structures de contrôle]
Pour chaque question, entoure la (ou les) bonne(s) réponse(s).
\begin{enumerate}
\item Un algorithme est :
\begin{enumerate}
\item un langage de programmation ;
\item une suite finie et ordonnée d'instructions permettant de résoudre un problème ;
\item un composant de l'unité centrale.
\end{enumerate}
\item La structure alternative simple s'écrit en pseudo-code :
\begin{enumerate}
\item \texttt{Si ... Alors ... Sinon ... FinSi} ;
\item \texttt{TantQue ... Faire ... FinTantQue} ;
\item \texttt{Lire ... Ecrire}.
\end{enumerate}
\item La boucle \texttt{Pour} est utilisée lorsque :
\begin{enumerate}
\item le nombre de répétitions est connu à l'avance ;
\item la condition d'arrêt dépend d'un calcul imprévisible ;
\item on veut lire une donnée au clavier.
\end{enumerate}
\item Dans un algorigramme, un test (une décision) se représente par :
\begin{enumerate}
\item un rectangle ;
\item un losange ;
\item un ovale.
\end{enumerate}
\end{enumerate}
\end{exercice}

\begin{exercice}[Vrai ou faux ? Justifie]
Réponds par Vrai ou Faux, puis justifie ta réponse en une phrase.
\begin{enumerate}
\item Une boucle \texttt{TantQue} peut ne jamais s'exécuter si sa condition est fausse dès le départ.
\item Dans une boucle \texttt{TantQue}, il n'est pas nécessaire de modifier la variable de la condition à l'intérieur de la boucle.
\item Un algorigramme commence et se termine toujours par un ovale.
\item L'instruction \texttt{Sinon} est obligatoire dans toute structure alternative.
\end{enumerate}
\end{exercice}

\begin{exercice}[Définitions]
Définis brièvement les termes suivants et donne un exemple pour chacun :
\begin{enumerate}
\item algorithme ;
\item structure alternative ;
\item structure itérative (boucle) ;
\item algorigramme.
\end{enumerate}
\end{exercice}

\begin{exercice}[Trace d'exécution]
On donne l'algorithme suivant :
\begin{lstlisting}
Algorithme Mystere
Variable n, i, s : Entier
Debut
    Lire(n)
    s <- 0
    Pour i allant de 1 a n Faire
        s <- s + i
    FinPour
    Ecrire(s)
Fin
\end{lstlisting}
\begin{enumerate}
\item Recopie et complète la table d'exécution ci-dessous pour $n = 4$ :
\end{enumerate}
\begin{center}
\begin{tabularx}{\linewidth}{|l|X|X|}
\hline
\rowcolor{popblueL} \textbf{Étape} & \textbf{Valeur de i} & \textbf{Valeur de s} \\
\hline
Avant la boucle & -- & \\
\hline
1\ier{} tour & 1 & \\
\hline
2\ieme{} tour & 2 & \\
\hline
3\ieme{} tour & 3 & \\
\hline
4\ieme{} tour & 4 & \\
\hline
\end{tabularx}
\end{center}
\begin{enumerate}
\setcounter{enumi}{1}
\item Quelle valeur l'algorithme affiche-t-il ? Que calcule cet algorithme en général ?
\end{enumerate}
\end{exercice}

\courssub{Je m'entraîne}

\begin{exercice}[Mineur ou majeur]
\begin{enumerate}
\item Écris en pseudo-code un algorithme qui lit l'âge d'un élève et affiche \og Mineur \fg{} si l'âge est strictement inférieur à 18 ans, et \og Majeur \fg{} sinon.
\item Trace l'exécution de ton algorithme pour un élève de 15 ans, puis pour un élève de 19 ans (indique la valeur du test et le message affiché dans chaque cas).
\end{enumerate}
\end{exercice}

\begin{exercice}[Somme des nombres pairs]
\begin{enumerate}
\item Écris en pseudo-code un algorithme qui calcule et affiche la somme des nombres pairs de 2 à 50, en utilisant une boucle \texttt{Pour}.
\item Donne la trace partielle de l'exécution : valeurs de la variable de boucle et de la somme pour les trois premiers tours et pour le dernier tour.
\item Quelle valeur finale est affichée ?
\end{enumerate}
\end{exercice}

\begin{exercice}[Le vendeur de crédit téléphonique]
Un vendeur de crédit téléphonique de quartier part de 0 FCFA d'économies et met de côté 2 500 FCFA chaque jour. Il veut savoir au bout de combien de jours il disposera d'au moins 30 000 FCFA pour acheter une mallette de recharge.
\begin{enumerate}
\item Écris en pseudo-code un algorithme utilisant une boucle \texttt{TantQue} qui calcule et affiche le nombre de jours nécessaires.
\item Vérifie par le calcul le résultat affiché par ton algorithme.
\item Que se passerait-il si l'on oubliait l'instruction qui augmente le nombre de jours dans la boucle ?
\end{enumerate}
\end{exercice}

\begin{exercice}[Le plus grand de deux nombres]
\begin{enumerate}
\item Écris en pseudo-code un algorithme qui lit deux nombres réels $a$ et $b$ et affiche le plus grand des deux.
\item Représente cet algorithme par un algorigramme (ovales, rectangles, losange et flèches).
\item Que faudrait-il ajouter pour traiter correctement le cas où les deux nombres sont égaux ?
\end{enumerate}
\end{exercice}

\courssub{Je me prépare au Probatoire}

\begin{exercice}[Chasse aux erreurs]
Le directeur d'une école demande à un élève de 3\ieme{} un algorithme qui lit un nombre entier $n$ et affiche la somme $1 + 2 + \dots + n$. L'élève propose le pseudo-code et l'algorigramme (décrit) ci-dessous, qui contiennent \textbf{trois erreurs} :
\begin{lstlisting}
Algorithme Somme
Variable n, i, s : Entier
Debut
    Lire(n)
    s <- 0
    i <- 1
    TantQue i <= n Faire
        s <- s + i
    FinTantQue
    Si s > 100 Alors
        Ecrire("Grande somme")
    Ecrire(s)
Fin
\end{lstlisting}
L'algorigramme tracé par l'élève représente le test \texttt{i <= n} par un \textbf{rectangle} au lieu de la forme correcte.
\begin{enumerate}
\item Releve les trois erreurs commises (une boucle infinie, un \texttt{FinSi} manquant, un mauvais symbole dans l'algorigramme) en expliquant pour chacune pourquoi c'est une erreur.
\item Écris le pseudo-code corrigé.
\item Indique le symbole correct pour représenter le test dans l'algorigramme, puis trace l'algorigramme complet corrigé.
\end{enumerate}
\end{exercice}

\begin{exercice}[Caisse d'épargne du quartier]
Mama Rose dépose 5 000 FCFA chaque semaine dans une caisse d'épargne. Elle souhaite acheter une machine à coudre coûtant 85 000 FCFA.
\begin{enumerate}
\item Écris en pseudo-code un algorithme qui lit le montant de l'objectif (85 000 FCFA) et le dépôt hebdomadaire (5 000 FCFA), puis calcule et affiche le nombre de semaines nécessaires pour atteindre ou dépasser l'objectif, en utilisant une boucle \texttt{TantQue}.
\item Trace l'exécution de ton algorithme pour les trois premières semaines (montant épargné et nombre de semaines à chaque tour de boucle).
\item Représente ton algorithme par un algorigramme.
\item Mama Rose décide finalement de déposer 10 000 FCFA par semaine. Sans réécrire tout l'algorithme, indique quelle donnée change et donne le nouveau nombre de semaines affiché.
\end{enumerate}
\end{exercice}

\begin{exercice}[Type BEPC : bulletin de notes]
Au collège bilingue de Bafoussam, le professeur principal veut un algorithme qui lit les trois notes d'un élève (sur 20) en informatique, en mathématiques et en français, calcule la moyenne, puis affiche :
\begin{itemize}
\item la décision : \og Admis \fg{} si la moyenne est supérieure ou égale à 10, \og Refusé \fg{} sinon ;
\item la mention : \og Passable \fg{} si $10 \le m < 12$, \og Assez bien \fg{} si $12 \le m < 14$, \og Bien \fg{} si $14 \le m < 16$, \og Très bien \fg{} si $m \ge 16$ (la mention n'est affichée que si l'élève est admis).
\end{itemize}
\begin{enumerate}
\item Écris en pseudo-code l'algorithme complet (déclaration des variables, lecture des trois notes, calcul de la moyenne, structures alternatives pour la décision et la mention).
\item Trace l'exécution de ton algorithme pour un élève ayant obtenu 13, 15 et 11 : donne la moyenne, la décision et la mention affichées.
\item Représente cet algorithme par un algorigramme complet (début, lectures, traitement, tests, affichages, fin).
\item Le proviseur souhaite maintenant que le programme refuse toute note inférieure à 0 ou supérieure à 20 en affichant \og Note invalide \fg{}. Indique où placer ce contrôle dans ton algorithme et écris la structure alternative correspondante.
\end{enumerate}
\end{exercice}

\newpage

\coursec{Corrigés des exercices}

\coursec{Utilisation des structures de contrôle}

\courssub{Je vérifie mes connaissances}

\begin{corrige}[Exercice 1 — QCM : structures de contrôle]
\begin{enumerate}
\item \textbf{b.} Un algorithme est une suite finie et ordonnée d'instructions permettant de résoudre un problème. Ce n'est pas un langage de programmation (c'est une description abstraite) ni un composant matériel.
\item \textbf{a.} La structure alternative simple s'écrit : \texttt{Si \textless condition\textgreater Alors \textless instructions\textgreater Sinon \textless instructions\textgreater FinSi}. La proposition b correspond à une boucle \texttt{TantQue}, la c à des entrées/sorties.
\item \textbf{a.} La boucle \texttt{Pour} est utilisée lorsque le nombre de répétitions est connu à l'avance (borne de début, borne de fin, pas). La boucle \texttt{TantQue} convient quand la condition d'arrêt dépend d'un calcul imprévisible.
\item \textbf{b.} Dans un algorigramme, un test (décision) se représente par un \textbf{losange}. Le rectangle représente une action/traitement, l'ovale le début/fin.
\end{enumerate}
\end{corrige}

\begin{corrige}[Exercice 2 — Vrai ou faux ? Justifie]
\begin{enumerate}
\item \textbf{Vrai.} Dans une boucle \texttt{TantQue}, la condition est évaluée \emph{avant} chaque itération. Si elle est fausse dès la première évaluation, le corps de la boucle n'est jamais exécuté (zéro itération).
\item \textbf{Faux.} Il est \emph{indispensable} de modifier la variable intervenant dans la condition à l'intérieur de la boucle, sinon la condition reste vraie indéfiniment et la boucle ne s'arrête jamais (boucle infinie).
\item \textbf{Vrai.} Par convention, tout algorigramme commence par un ovale \og Début \fg{} et se termine par un ovale \og Fin \fg{}.
\item \textbf{Faux.} L'instruction \texttt{Sinon} est \emph{optionnelle}. On parle d'alternative simple (\texttt{Si ... Alors ... FinSi}) quand il n'y a pas de \texttt{Sinon}, et d'alternative complète quand il y en a un.
\end{enumerate}
\end{corrige}

\begin{corrige}[Exercice 3 — Définitions]
\begin{enumerate}
\item \textbf{Algorithme :} Description finie, non ambiguë et ordonnée d'une suite d'opérations élémentaires permettant de résoudre un problème ou d'obtenir un résultat à partir de données d'entrée.
\emph{Exemple :} La recette de cuisine pour préparer le \og ndolé \fg{} (entrées : ingrédients ; sorties : plat prêt ; étapes : laver, couper, cuire\dots).

\item \textbf{Structure alternative :} Structure de contrôle qui permet d'exécuter un bloc d'instructions \textbf{si une condition est vraie}, et éventuellement un autre bloc \textbf{si elle est fausse}. Elle réalise un choix.
\emph{Exemple :} \og Si la note $\ge$ 10 Alors \og Admis \fg{} Sinon \og Refusé \fg{} FinSi \fg{}.

\item \textbf{Structure itérative (boucle) :} Structure de contrôle qui permet de répéter un bloc d'instructions \textbf{tant qu'une condition est vraie} (\texttt{TantQue}) ou \textbf{un nombre déterminé de fois} (\texttt{Pour}). Elle évite la réécriture du code.
\emph{Exemple :} Afficher les nombres de 1 à 10 avec \texttt{Pour i de 1 à 10 Faire Ecrire(i) FinPour}.

\item \textbf{Algorigramme :} Représentation graphique normalisée d'un algorithme à l'aide de symboles géométriques (ovale, rectangle, losange, flèches) reliés entre eux, qui montre la logique et le flux de contrôle sans dépendre d'un langage de programmation.
\emph{Exemple :} Le diagramme en losange/rectangles/ovales de l'exercice 8 ci-dessous.
\end{enumerate}
\end{corrige}

\begin{corrige}[Exercice 4 — Trace d'exécution]
\begin{enumerate}
\item Table d'exécution pour $n = 4$ :
\begin{center}
\begin{tabularx}{\linewidth}{|l|X|X|}
\hline
\rowcolor{popblueL} \textbf{Étape} & \textbf{Valeur de i} & \textbf{Valeur de s} \\
\hline
Avant la boucle & -- & 0 \\
\hline
1\textsuperscript{er} tour & 1 & 1 \\
\hline
2\textsuperscript{e} tour & 2 & 3 \\
\hline
3\textsuperscript{e} tour & 3 & 6 \\
\hline
4\textsuperscript{e} tour & 4 & 10 \\
\hline
\end{tabularx}
\end{center}
\item L'algorithme affiche \textbf{10}.
En général, cet algorithme calcule la \textbf{somme des $n$ premiers entiers naturels non nuls} : $S = 1 + 2 + \dots + n = \frac{n(n+1)}{2}$.
Vérification : pour $n=4$, $\frac{4 \times 5}{2} = 10$.
\end{enumerate}
\end{corrige}

\courssub{Je m'entraîne}

\begin{corrige}[Exercice 5 — Mineur ou majeur]
\begin{enumerate}
\item Algorithme en pseudo-code :
\begin{lstlisting}[language=]
Algorithme MajeurMineur
Variable age : Entier
Debut
    Ecrire("Entrez l'âge de l'élève : ")
    Lire(age)
    Si age < 18 Alors
        Ecrire("Mineur")
    Sinon
        Ecrire("Majeur")
    FinSi
Fin
\end{lstlisting}
\item Traces d'exécution :
\begin{itemize}
\item \textbf{Cas 1 :} Âge = 15 ans.
\begin{itemize}
\item Lecture : \texttt{age} $\leftarrow$ 15.
\item Test : \texttt{15 < 18} est \textbf{VRAI}.
\item Exécution du \texttt{Alors} : Affichage \og Mineur \fg{}.
\end{itemize}
\item \textbf{Cas 2 :} Âge = 19 ans.
\begin{itemize}
\item Lecture : \texttt{age} $\leftarrow$ 19.
\item Test : \texttt{19 < 18} est \textbf{FAUX}.
\item Exécution du \texttt{Sinon} : Affichage \og Majeur \fg{}.
\end{itemize}
\end{itemize}
\end{enumerate}
\end{corrige}

\begin{corrige}[Exercice 6 — Somme des nombres pairs]
\begin{enumerate}
\item Algorithme (boucle \texttt{Pour} avec pas de 2) :
\begin{lstlisting}[language=]
Algorithme SommePairs
Variable i, s : Entier
Debut
    s <- 0
    Pour i de 2 a 50 pas 2 Faire
        s <- s + i
    FinPour
    Ecrire("La somme des nombres pairs de 2 a 50 est : ", s)
Fin
\end{lstlisting}
\emph{Variante possible :} \texttt{Pour i de 1 a 25 Faire s <- s + 2*i FinPour}.
\item Trace partielle (valeurs de \texttt{i} et \texttt{s}) :
\begin{center}
\begin{tabularx}{\linewidth}{|l|X|X|}
\hline
\rowcolor{popblueL} \textbf{Tour} & \textbf{i} & \textbf{s} \\
\hline
1\textsuperscript{er} & 2 & 2 \\
\hline
2\textsuperscript{e} & 4 & 6 \\
\hline
3\textsuperscript{e} & 6 & 12 \\
\hline
$\vdots$ & $\vdots$ & $\vdots$ \\
\hline
Dernier (25\textsuperscript{e}) & 50 & 650 \\
\hline
\end{tabularx}
\end{center}
\item Valeur finale affichée : \textbf{650}.
Vérification : Somme = $2+4+\dots+50 = 2(1+2+\dots+25) = 2 \times \frac{25 \times 26}{2} = 25 \times 26 = 650$.
\end{enumerate}
\end{corrige}

\begin{corrige}[Exercice 7 — Le vendeur de crédit téléphonique]
\begin{enumerate}
\item Algorithme avec \texttt{TantQue} :
\begin{lstlisting}[language=]
Algorithme JoursEpargne
Variable epargne, jours : Entier
Debut
    epargne <- 0
    jours <- 0
    TantQue epargne < 30000 Faire
        epargne <- epargne + 2500
        jours <- jours + 1
    FinTantQue
    Ecrire("Nombre de jours nécessaires : ", jours)
    Ecrire("Épargne finale : ", epargne, " FCFA")
Fin
\end{lstlisting}
\item Vérification par le calcul :
$30\,000 \div 2\,500 = 12$ jours exactement.
Après 12 jours : $12 \times 2\,500 = 30\,000$ FCFA. La condition \texttt{epargne < 30000} devient fausse, la boucle s'arrête. L'algorithme affiche 12. \checkmark
\item Si l'on oublie \texttt{jours <- jours + 1} dans la boucle :
La variable \texttt{jours} reste à 0. La condition \texttt{epargne < 30000} ne dépend pas de \texttt{jours}, donc l'épargne augmente normalement et la boucle s'arrête au bout de 12 tours. \textbf{Mais} l'affichage final donnera \og Nombre de jours nécessaires : 0 \fg{}, ce qui est \textbf{incorrect}. Si par contre on oublie \texttt{epargne <- epargne + 2500}, c'est une \textbf{boucle infinie} car la condition reste vraie éternellement.
\end{enumerate}
\end{corrige}

\begin{corrige}[Exercice 8 — Le plus grand de deux nombres]
\begin{enumerate}
\item Algorithme :
\begin{lstlisting}[language=]
Algorithme PlusGrand
Variable a, b : Reel
Debut
    Ecrire("Entrez a : "); Lire(a)
    Ecrire("Entrez b : "); Lire(b)
    Si a > b Alors
        Ecrire("Le plus grand est : ", a)
    Sinon
        Si a = b Alors
            Ecrire("Les deux nombres sont égaux : ", a)
        Sinon
            Ecrire("Le plus grand est : ", b)
        FinSi
    FinSi
Fin
\end{lstlisting}
\emph{Note :} La version simple (sans gestion de l'égalité explicite) affiche \texttt{b} si \texttt{a <= b}, ce qui est correct pour "le plus grand" mais ne signale pas l'égalité.
\item Algorigramme :
\begin{popfigure}
\begin{tikzpicture}[
    node distance=1.2cm and 1.5cm,
    startstop/.style={ellipse, draw=popblue, thick, fill=popblueL, text width=2.5cm, align=center, minimum height=1cm},
    process/.style={rectangle, draw=popdark, thick, fill=popcream, text width=3.5cm, align=center, minimum height=1cm, rounded corners=3pt},
    decision/.style={diamond, draw=poporange, thick, fill=poporangeL, text width=2.8cm, align=center, minimum height=1.5cm, aspect=2},
    arrow/.style={-Stealth, thick, popdark},
    label/.style={font=\small, text=poppurple, midway, fill=white, inner sep=1pt}
]
% Nodes
\node (start) [startstop] {Début};
\node (read) [process, below of=start] {Lire a, b};
\node (test1) [decision, below of=read] {a > b ?};
\node (printa) [process, left of=test1, xshift=-2.5cm] {Afficher a};
\node (test2) [decision, below of=test1, yshift=-1.5cm] {a = b ?};
\node (printeq) [process, left of=test2, xshift=-2.5cm] {Afficher "Égaux"};
\node (printb) [process, right of=test2, xshift=2.5cm] {Afficher b};
\node (end) [startstop, below of=test2, yshift=-1.5cm] {Fin};
% Arrows
\draw [arrow] (start) -- (read);
\draw [arrow] (read) -- (test1);
\draw [arrow] (test1) -- node[label, left] {Oui} (printa);
\draw [arrow] (test1) -- node[label, right] {Non} (test2);
\draw [arrow] (test2) -- node[label, left] {Oui} (printeq);
\draw [arrow] (test2) -- node[label, right] {Non} (printb);
\draw [arrow] (printa) |- (end);
\draw [arrow] (printeq) |- (end);
\draw [arrow] (printb) |- (end);
\end{tikzpicture}
\legende{Algorigramme : recherche du plus grand de deux nombres avec gestion de l'égalité}
\end{popfigure}
\item Pour traiter correctement le cas d'égalité, il faut ajouter un second test \texttt{Si a = b Alors ... Sinon ... FinSi} à l'intérieur de la branche \texttt{Sinon} du premier test (comme montré dans le pseudo-code et l'algorigramme ci-dessus). Sans cela, l'algorithme considère que \texttt{b} est le plus grand quand \texttt{a = b}, ce qui est inexact.
\end{enumerate}
\end{corrige}

\courssub{Je me prépare au Probatoire}

\begin{corrige}[Exercice 9 — Chasse aux erreurs]
\begin{enumerate}
\item \textbf{Les trois erreurs :}
\begin{enumerate}
\item \textbf{Boucle infinie (erreur logique) :} Dans la boucle \texttt{TantQue i <= n Faire ... FinTantQue}, la variable \texttt{i} n'est \textbf{jamais incrémentée}. Si \texttt{n $\ge$ 1}, la condition \texttt{i <= n} reste vraie pour toujours $\rightarrow$ la boucle ne s'arrête jamais.
\item \textbf{\texttt{FinSi} manquant (erreur de syntaxe/structure) :} L'instruction \texttt{Si s > 100 Alors Ecrire("Grande somme")} n'est pas fermée par un \texttt{FinSi}. En pseudo-code structuré, tout \texttt{Si} doit avoir son \texttt{FinSi} (même sans \texttt{Sinon}).
\item \textbf{Mauvais symbole dans l'algorigramme (erreur de représentation) :} Le test \texttt{i <= n} est représenté par un \textbf{rectangle}. Or, un test (décision) doit \textbf{obligatoirement} être représenté par un \textbf{losange}. Le rectangle est réservé aux traitements/affectations.
\end{enumerate}
\item \textbf{Pseudo-code corrigé :}
\begin{lstlisting}[language=]
Algorithme Somme
Variable n, i, s : Entier
Debut
    Lire(n)
    s <- 0
    i <- 1
    TantQue i <= n Faire
        s <- s + i
        i <- i + 1     // Correction 1 : incrémentation de i
    FinTantQue
    Si s > 100 Alors
        Ecrire("Grande somme")
    FinSi              // Correction 2 : FinSi ajouté
    Ecrire(s)
Fin
\end{lstlisting}
\item \textbf{Symbole correct :} Le \textbf{losange} (diamant).
\textbf{Algorigramme corrigé :}
\begin{popfigure}
\begin{tikzpicture}[
    node distance=1.1cm and 1.3cm,
    startstop/.style={ellipse, draw=popblue, thick, fill=popblueL, text width=2.2cm, align=center, minimum height=0.9cm},
    process/.style={rectangle, draw=popdark, thick, fill=popcream, text width=3cm, align=center, minimum height=0.8cm, rounded corners=3pt},
    decision/.style={diamond, draw=poporange, thick, fill=poporangeL, text width=2.5cm, align=center, minimum height=1.2cm, aspect=1.8},
    arrow/.style={-Stealth, thick, popdark},
    label/.style={font=\small, text=poppurple, midway, fill=white, inner sep=1pt}
]
\node (start) [startstop] {Début};
\node (read) [process, below of=start] {Lire n};
\node (init) [process, below of=read] {s $\leftarrow$ 0 ; i $\leftarrow$ 1};
\node (test) [decision, below of=init] {i $\le$ n ?};
\node (body) [process, below of=test, yshift=-0.5cm] {s $\leftarrow$ s + i ; i $\leftarrow$ i + 1};
\node (test2) [decision, below of=body, yshift=-0.5cm] {s > 100 ?};
\node (msg) [process, left of=test2, xshift=-2.5cm] {Afficher "Grande somme"};
\node (print) [process, below of=test2, yshift=-1cm] {Afficher s};
\node (end) [startstop, below of=print] {Fin};
\draw [arrow] (start) -- (read);
\draw [arrow] (read) -- (init);
\draw [arrow] (init) -- (test);
\draw [arrow] (test) -- node[label, left] {Oui} (body);
\draw [arrow] (body) -- (test);
\draw [arrow] (test) -- node[label, right] {Non} (test2);
\draw [arrow] (test2) -- node[label, left] {Oui} (msg);
\draw [arrow] (msg) |- (print);
\draw [arrow] (test2) -- node[label, right] {Non} (print);
\draw [arrow] (print) -- (end);
\end{tikzpicture}
\legende{Algorigramme corrigé : somme 1 à n avec test sur la somme}
\end{popfigure}
\end{enumerate}
\end{corrige}

\begin{corrige}[Exercice 10 — Caisse d'épargne du quartier]
\begin{enumerate}
\item Algorithme \texttt{TantQue} :
\begin{lstlisting}[language=]
Algorithme SemainesEpargne
Variable objectif, depot, epargne, semaines : Entier
Debut
    Ecrire("Objectif (FCFA) : "); Lire(objectif)
    Ecrire("Dépôt hebdomadaire (FCFA) : "); Lire(depot)
    epargne <- 0
    semaines <- 0
    TantQue epargne < objectif Faire
        epargne <- epargne + depot
        semaines <- semaines + 1
    FinTantQue
    Ecrire("Nombre de semaines nécessaires : ", semaines)
    Ecrire("Montant épargné : ", epargne, " FCFA")
Fin
\end{lstlisting}
\item Trace pour les 3 premières semaines (objectif=85000, depot=5000) :
\begin{center}
\begin{tabularx}{\linewidth}{|l|X|X|}
\hline
\rowcolor{popblueL} \textbf{Semaine (tour)} & \textbf{Épargne cumulée} & \textbf{Variable semaines} \\
\hline
1 & 5\,000 & 1 \\
\hline
2 & 10\,000 & 2 \\
\hline
3 & 15\,000 & 3 \\
\hline
\end{tabularx}
\end{center}
\item Algorigramme :
\begin{popfigure}
\begin{tikzpicture}[
    node distance=1.1cm and 1.3cm,
    startstop/.style={ellipse, draw=popblue, thick, fill=popblueL, text width=2.2cm, align=center, minimum height=0.9cm},
    process/.style={rectangle, draw=popdark, thick, fill=popcream, text width=3.5cm, align=center, minimum height=0.8cm, rounded corners=3pt},
    decision/.style={diamond, draw=poporange, thick, fill=poporangeL, text width=2.8cm, align=center, minimum height=1.2cm, aspect=2},
    arrow/.style={-Stealth, thick, popdark},
    label/.style={font=\small, text=poppurple, midway, fill=white, inner sep=1pt}
]
\node (start) [startstop] {Début};
\node (read) [process, below of=start] {Lire objectif, depot};
\node (init) [process, below of=read] {epargne $\leftarrow$ 0 ; semaines $\leftarrow$ 0};
\node (test) [decision, below of=init] {epargne < objectif ?};
\node (body) [process, below of=test, yshift=-0.5cm] {epargne $\leftarrow$ epargne + depot ; semaines $\leftarrow$ semaines + 1};
\node (print) [process, below of=test, yshift=-2.5cm] {Afficher semaines, epargne};
\node (end) [startstop, below of=print] {Fin};
\draw [arrow] (start) -- (read);
\draw [arrow] (read) -- (init);
\draw [arrow] (init) -- (test);
\draw [arrow] (test) -- node[label, left] {Oui} (body);
\draw [arrow] (body) -- (test);
\draw [arrow] (test) -- node[label, right] {Non} (print);
\draw [arrow] (print) -- (end);
\end{tikzpicture}
\legende{Algorigramme : calcul du nombre de semaines pour atteindre l'objectif}
\end{popfigure}
\item Changement : la valeur de la variable \texttt{depot} passe de 5\,000 à \textbf{10\,000} FCFA.
Nouveau calcul : $85\,000 \div 10\,000 = 8,5$. Il faut \textbf{9 semaines} (car au bout de 8 semaines : 80\,000 < 85\,000 ; au bout de 9 : 90\,000 $\ge$ 85\,000).
L'algorithme affichera \textbf{9}.
\end{enumerate}
\end{corrige}

\begin{corrige}[Exercice 11 — Type BEPC : bulletin de notes]
\begin{enumerate}
\item Algorithme complet :
\begin{lstlisting}[language=]
Algorithme Bulletin
Variable noteInfo, noteMaths, noteFrancais, moyenne : Reel
Variable decision, mention : Chaine
Debut
    // Lecture des notes
    Ecrire("Note Informatique (sur 20) : "); Lire(noteInfo)
    Ecrire("Note Mathématiques (sur 20) : "); Lire(noteMaths)
    Ecrire("Note Français (sur 20) : "); Lire(noteFrancais)
    
    // Calcul de la moyenne
    moyenne <- (noteInfo + noteMaths + noteFrancais) / 3
    
    // Décision Admis/Refusé
    Si moyenne >= 10 Alors
        decision <- "Admis"
        // Mention (seulement si admis)
        Si moyenne < 12 Alors
            mention <- "Passable"
        Sinon
            Si moyenne < 14 Alors
                mention <- "Assez bien"
            Sinon
                Si moyenne < 16 Alors
                    mention <- "Bien"
                Sinon
                    mention <- "Très bien"
                FinSi
            FinSi
        FinSi
        Ecrire("Décision : ", decision)
        Ecrire("Mention : ", mention)
    Sinon
        decision <- "Refusé"
        Ecrire("Décision : ", decision)
        // Pas de mention pour un refusé
    FinSi
Fin
\end{lstlisting}
\item Trace pour notes : 13, 15, 11.
\begin{itemize}
\item Somme = 13 + 15 + 11 = 39.
\item Moyenne = 39 / 3 = \textbf{13}.
\item Test \texttt{moyenne >= 10} : 13 $\ge$ 10 $\rightarrow$ \textbf{VRAI} $\rightarrow$ \texttt{decision} = "Admis".
\item Test mention :
\begin{itemize}
\item 13 < 12 ? FAUX.
\item 13 < 14 ? VRAI $\rightarrow$ \texttt{mention} = "Assez bien".
\end{itemize}
\item Affichage :
\og Décision : Admis \fg{}
\og Mention : Assez bien \fg{}
\end{itemize}
\item Algorigramme complet :
\begin{popfigure}
\begin{tikzpicture}[
    node distance=1.0cm and 1.2cm,
    startstop/.style={ellipse, draw=popblue, thick, fill=popblueL, text width=2cm, align=center, minimum height=0.8cm, font=\small},
    process/.style={rectangle, draw=popdark, thick, fill=popcream, text width=3.2cm, align=center, minimum height=0.7cm, rounded corners=3pt, font=\small},
    decision/.style={diamond, draw=poporange, thick, fill=poporangeL, text width=2.5cm, align=center, minimum height=1.1cm, aspect=2, font=\small},
    arrow/.style={-Stealth, thick, popdark},
    label/.style={font=\tiny, text=poppurple, midway, fill=white, inner sep=1pt}
]
% Nodes
\node (start) [startstop] {Début};
\node (read) [process, below of=start] {Lire noteInfo, noteMaths, noteFrancais};
\node (calc) [process, below of=read] {moyenne $\leftarrow$ (somme)/3};
\node (testAdmis) [decision, below of=calc] {moyenne $\ge$ 10 ?};
\node (refuse) [process, right of=testAdmis, xshift=2.8cm] {decision $\leftarrow$ "Refusé"\newline Afficher decision};
\node (admis) [process, left of=testAdmis, xshift=-2.8cm] {decision $\leftarrow$ "Admis"\newline Afficher decision};
\node (test12) [decision, below of=admis, yshift=-0.5cm] {moyenne < 12 ?};
\node (pass) [process, left of=test12, xshift=-2.2cm] {mention $\leftarrow$ "Passable"};
\node (test14) [decision, below of=test12, yshift=-0.5cm] {moyenne < 14 ?};
\node (ab) [process, left of=test14, xshift=-2.2cm] {mention $\leftarrow$ "Assez bien"};
\node (test16) [decision, below of=test14, yshift=-0.5cm] {moyenne < 16 ?};
\node (bien) [process, left of=test16, xshift=-2.2cm] {mention $\leftarrow$ "Bien"};
\node (tb) [process, right of=test16, xshift=2.2cm] {mention $\leftarrow$ "Très bien"};
\node (affMention) [process, below of=tb, yshift=-0.5cm] {Afficher mention};
\node (end) [startstop, below of=affMention] {Fin};
% Arrows - Main flow
\draw [arrow] (start) -- (read);
\draw [arrow] (read) -- (calc);
\draw [arrow] (calc) -- (testAdmis);
\draw [arrow] (testAdmis) -- node[label, above] {Non} (refuse);
\draw [arrow] (testAdmis) -- node[label, above] {Oui} (admis);
\draw [arrow] (refuse) |- (end);
\draw [arrow] (admis) -- (test12);
% Mention chain
\draw [arrow] (test12) -- node[label, above] {Oui} (pass);
\draw [arrow] (test12) -- node[label, right] {Non} (test14);
\draw [arrow] (test14) -- node[label, above] {Oui} (ab);
\draw [arrow] (test14) -- node[label, right] {Non} (test16);
\draw [arrow] (test16) -- node[label, above] {Oui} (bien);
\draw [arrow] (test16) -- node[label, above] {Non} (tb);
% Merge to affMention
\draw [arrow] (pass) |- (affMention);
\draw [arrow] (ab) |- (affMention);
\draw [arrow] (bien) |- (affMention);
\draw [arrow] (tb) -- (affMention);
\draw [arrow] (affMention) -- (end);
\end{tikzpicture}
\legende{Algorigramme complet : bulletin avec décision et mention imbriquée}
\end{popfigure}
\item Contrôle de validité des notes (0 à 20) :
Il faut placer ce contrôle \textbf{immédiatement après chaque lecture} (ou dans une boucle de lecture), \textbf{avant le calcul de la moyenne}. Si une note est invalide, on affiche le message et on arrête (ou on redemande).
Structure alternative à insérer après \texttt{Lire(noteInfo)} (et idem pour les deux autres) :
\begin{lstlisting}[language=]
    Si noteInfo < 0 Ou noteInfo > 20 Alors
        Ecrire("Note invalide")
        // Optionnel : arrêter le programme ou redemander la note
    FinSi
\end{lstlisting}
\emph{Point de vigilance :} L'opérateur \og Ou \fg{} (ou \texttt{OU} / \texttt{||}) est nécessaire car la note est invalide si \textbf{au moins une} des deux conditions est vraie (trop petite \textbf{OU} trop grande). L'utilisation de \og Et \fg{} serait une erreur logique (une note ne peut pas être \emph{à la fois} $<0$ et $>20$).
\end{enumerate}
\end{corrige}

\newpage

\coursec{Fiche bilan}

\begin{popfigure}
\begin{tikzpicture}[mindmap, grow cyclic, every node/.style={concept, text=white, font=\small\bfseries},
root concept/.append style={concept color=popdark, font=\large\bfseries},
level 1/.append style={level distance=3.6cm, sibling angle=60},
level 2/.append style={level distance=2.2cm, sibling angle=45, font=\scriptsize}]
\node[root concept]{Structures de contrôle}
  child[concept color=popblue]{ node {Algorithme et pseudo-code}
    child { node[concept color=popblueL, text=popdark] {Suite d'actions} }
    child { node[concept color=popblueL, text=popdark] {Variables} }
  }
  child[concept color=poporange]{ node {Structure alternative}
    child { node[concept color=poporangeL, text=popdark] {Si \dots Alors} }
    child { node[concept color=poporangeL, text=popdark] {Sinon} }
  }
  child[concept color=popgreen]{ node {Boucle Tant que}
    child { node[concept color=popgreenL, text=popdark] {Condition testée avant} }
    child { node[concept color=popgreenL, text=popdark] {0 ou $n$ passages} }
  }
  child[concept color=poppurple]{ node {Boucle Pour}
    child { node[concept color=poppurpleL, text=popdark] {Compteur} }
    child { node[concept color=poppurpleL, text=popdark] {Nombre de tours connu} }
  }
  child[concept color=poppink]{ node {Algorigramme}
    child { node[concept color=poppinkL, text=popdark] {Ovale : Début/Fin} }
    child { node[concept color=poppinkL, text=popdark] {Losange : test} }
  }
  child[concept color=popgold]{ node {Résolution de problèmes}
    child { node[concept color=popgoldL, text=popdark] {Analyser} }
    child { node[concept color=popgoldL, text=popdark] {Justifier} }
  };
\end{tikzpicture}
\legende{Carte mentale de la leçon : les structures de contrôle en algorithmique.}
\end{popfigure}

\begin{bilan}
\textbf{Définitions clés}
\begin{itemize}
  \item \cle{Algorithme} : suite finie et ordonnée d'actions (instructions) permettant de résoudre un problème.
  \item \cle{Pseudo-code} : langage simple, en français, utilisé pour écrire un algorithme (mots-clés : Algorithme, Début, Fin, Lire, Écrire).
  \item \cle{Variable} : zone mémoire nommée qui contient une valeur pouvant changer pendant l'exécution.
  \item \cle{Structure alternative} : permet d'exécuter des actions \emph{selon qu'une condition est vraie ou fausse}.
  \item \cle{Structure itérative} (boucle) : permet de répéter des actions plusieurs fois.
  \item \cle{Algorigramme} : représentation graphique d'un algorithme à l'aide de symboles normalisés reliés par des flèches.
\end{itemize}

\textbf{Syntaxes à connaître par c\oe ur}
\begin{center}
\begin{tabularx}{\linewidth}{|l|X|}
\hline
\rowcolor{popblueL} \textbf{Structure} & \textbf{Pseudo-code} \\
\hline
Alternative simple & Si \emph{condition} Alors \emph{actions} FinSi \\
\hline
Alternative complète & Si \emph{condition} Alors \emph{actions} Sinon \emph{autres actions} FinSi \\
\hline
Boucle Tant que & Tant que \emph{condition} Faire \emph{actions} FinTantQue \\
\hline
Boucle Pour & Pour $i$ de $1$ à $n$ Faire \emph{actions} FinPour \\
\hline
\end{tabularx}
\end{center}

\textbf{Symboles de l'algorigramme}
\begin{itemize}
  \item \textbf{Ovale} : Début ou Fin ; \quad \textbf{Rectangle} : traitement (calcul, affectation) ;
  \item \textbf{Parallélogramme} : entrée/sortie (Lire, Écrire) ; \quad \textbf{Losange} : test (condition) avec deux sorties Oui/Non ;
  \item \textbf{Flèches} : sens d'exécution.
\end{itemize}

\textbf{Méthodes express}
\begin{itemize}
  \item Choisir la boucle : nombre de répétitions \emph{connu à l'avance} $\rightarrow$ boucle \textbf{Pour} ; répétition \emph{selon une condition} $\rightarrow$ boucle \textbf{Tant que}.
  \item Rédiger un algorithme : (1) comprendre le problème, (2) identifier les données (Lire) et les résultats (Écrire), (3) écrire les traitements, (4) tester à la main avec des valeurs.
  \item Justifier : expliquer le choix de la structure utilisée et vérifier le résultat sur un exemple concret.
\end{itemize}

\begin{attention}[Pièges fréquents]
Dans un \textbf{Tant que}, une action à l'intérieur de la boucle doit faire évoluer la condition, sinon la boucle ne s'arrête jamais (\emph{boucle infinie}). Ne pas oublier le \textbf{Sinon} quand les deux cas doivent être traités, ni le \textbf{FinSi} / \textbf{FinPour} / \textbf{FinTantQue}.
\end{attention}
\end{bilan}

\begin{aretenir}[Checklist avant le BEPC]
\begin{itemize}
  \item $\square$ Je sais définir un algorithme et donner un exemple de la vie courante (recette, itinéraire).
  \item $\square$ Je sais écrire un algorithme en pseudo-code avec Algorithme, Début, Fin.
  \item $\square$ Je sais utiliser \texttt{Lire} et \texttt{Écrire} pour les entrées et sorties.
  \item $\square$ Je sais écrire une structure alternative Si \dots Alors \dots Sinon \dots FinSi.
  \item $\square$ Je sais choisir entre la boucle Pour et la boucle Tant que.
  \item $\square$ Je sais écrire une boucle Pour avec un compteur (de $1$ à $n$).
  \item $\square$ Je sais écrire une boucle Tant que et éviter la boucle infinie.
  \item $\square$ Je connais les symboles de l'algorigramme (ovale, rectangle, parallélogramme, losange, flèches).
  \item $\square$ Je sais dessiner l'algorigramme d'un algorithme simple.
  \item $\square$ Je sais tester mon algorithme à la main avec des valeurs concrètes.
  \item $\square$ Je sais rédiger et justifier ma démarche face à un problème concret.
\end{itemize}
\end{aretenir}

\findecours

\end{document}
